%----------------------------------------------------------------------------------------------------%
%----------------------------------------------------------------------------------------------------%
%%%%%%%%%%%                                              HEADER                                                %%%%%%%%%%%
%----------------------------------------------------------------------------------------------------%
%----------------------------------------------------------------------------------------------------%
%
% Universale header.tex – gilt für Arbeitsblätter, Leistungskontrollen, Handreichungen, Nachschlagewerke.
% Neue Pakete und TikZ-Bibliotheken werden HIER ergänzt (nicht im Dokument) und unten protokolliert.
% Quelle: https://github.com/mq90/latex-skills-schule  (CC BY 4.0, Martin Quandt)
%
% Änderungsprotokoll
% 2026-11-09  zusammengeführt aus allen Handreichungs-Headern: TikZ-Bibliotheken fit, shapes.multipart,
%             arrows.meta, shapes.callouts, circuits.logic.US, circuits.logic.IEC; Pakete array, booktabs,
%             longtable; xskak, chessboard, sudoku, tikzpeople, tipa (safe); MusiXTeX als auskommentierte
%             Option (definiert \cup als Achtelnote um – bricht $A \cup B$)
% 2026-11-11  Naturwissenschaften: circuitikz [european] (Schaltpläne), mhchem [version=4] (\ce),
%             chemfig (Strukturformeln), pgf-PeriodicTable (\pgfPT), genealogytree [all] (Stammbäume).
%             NICHT siunitx: kollidiert mit units (\unit[10]{cm}); siunitx nur in eigenständigen Dokumenten
% 2026-11-11  TikZ-Bibliothek decorations.pathmorphing (Federn, Wellen: decoration=coil/snake)
% 2026-11-23  twemojis: farbige Emojis mit pdflatex, z.B. \twemoji[height=1cm]{man standing} (Skill latex-sachskizzen)
%






%----------------------------------------------------------------------------------------------------%
%%%%%%%%%%%%%                            Gestaltung des Dokuments                                %%%%%%%%%%
%----------------------------------------------------------------------------------------------------%



%Definition der Dokumentenklasse
\documentclass[12pt]{article} %book, article, report, standalone



%Pakete für eine saubere Kodierung
\usepackage[utf8]{inputenc} %UTF-8 als Zeichenkodierung  ... ABER nicht immer notwendig, je nach Compiler
\usepackage[T1]{fontenc} %Verwendet westeuropäische Schriftkodierung
\usepackage[ngerman]{babel} %Ermöglicht deutsche Übersetzung und Silbentrennung
\usepackage{lmodern} %Verwendet moderne Schriftarten



%Pakete für die Gestaltung des Dokuments
\usepackage{geometry} %Ermöglicht die Ausrichtung des Blattes mit \geometry{...}
\geometry{a4paper, top=15mm, left=15mm, right=20mm, bottom=15mm, headsep=10mm, footskip=10mm} %gewählte Ausrichtung des Blattes
\pagestyle{empty} %Art der Seitennummerierung, z.B.: empty , plain
\usepackage{fancyhdr} %Ermöglicht die Gestaltung einer Kopfzeile für LK's oder KA's 
\usepackage{lscape} %Ermöglicht eventuelles Querformat des Blattes mit \begin{landscape} ... \end{landscape}
\usepackage{pdfpages} %Ermöglicht einbinden PDF-Dokumente bzw. einzelene Seiten mit \includepdf[pages={1,3,5}]{Dateiname} bzw. pages=-
\usepackage[german]{layout} %Ermöglicht es mit den Befehl \layout sich die Default-Layout beschreiben zu lassen (Bild vom Layout)
%----------------------------------------------------------------------------------------------------%








%----------------------------------------------------------------------------------------------------%
%%%%%%%%%%%%                                        Packages                                                  %%%%%%%%%%%
%----------------------------------------------------------------------------------------------------%



%Pakete für die Strukturierung des Dokuments
\usepackage[shortlabels]{enumitem} %Ermöglicht die Nummerierung selbst festzulegen mit \begin{enumerate}[Art der Nummerierung] ...   \end{enumerate}
\parindent=0px %Setzt die Absatzeinrückung auf 0 Pixel und ernfernt damit unnötige Leerzeichen vor dem ersten Wort eines neuen Abschnitts
\usepackage{setspace} %Ermöglicht verschiedene Funktion zur Abstandsregelung, z.B.: \hspace{2.5cm} ; \onehalfspacing
\onehalfspacing %Setzt den Zeilenabstand auf 1,5
\usepackage[style=german]{csquotes}%Ermöglicht den deutschen Zitierstil darstellen zu lassen mt \enquote{Zitat}
\usepackage[breaklinks=true,%Zeilenumbruch der Verlinkung ist bündig
			colorlinks=true,%Links können farbig sein
			linkcolor=blue,%Linkfarbe wird auf blau gesetzt
			urlcolor=darkgray,%Farb der URL wird auf dunkelgrau gesetzt
			plainpages=false,%
			hyperfootnotes=false,%
			pdfpagelabels=false%
		]{hyperref} %Ermöglicht die Verlinkung innerhalb des Dokuments 
\usepackage{multicol} %Ermöglicht gleichgroße Aufteilung des Dokuments in n Spalten mit \begin{multicols}{n} ... \end{multicols}
\usepackage{xcolor,colortbl} %Ermöglicht Zeilen, Spalten und Zellen einer Tabelle zu färben
\usepackage{caption} %Ermöglicht die Bezeichnung und Strukturierung von Abbildungen oder Quellcode
\usepackage{fancybox} % Ermöglicht die Umrandung einzelner Abschnitte
\usepackage[most]{tcolorbox} %Ermöglicht die Nutzung von Colorboxen für Merksätze, Lösungen, Definitionen, usw mit \newtcolorbox
\tcbuselibrary{breakable} %Aufteilen der Colorboxen, welche über mehrere Seiten gehen
\usepackage{ifthen} % Ermöglicht \ifthenelse Befehle
\allowdisplaybreaks %Erlaubt Seitenumbrüche innerhalb von align-Umgebung



%Nützliche Pakete für Symbole und mathematischen Formelsatz
\usepackage{amsmath, amssymb, mathrsfs, amsfonts, amsthm, dsfont} %sämtliche Pakete für mathematische Zeichen, Symbole und Formelsatz
\usepackage{units} %Ermöglicht die Angabe von Einheiten mit \unit[Wert]{Einheit}
\usepackage{gensymb} %Ermöglicht Nutzung von \degree für das macro des Grad-Symbols °
\usepackage[official]{eurosym} %Ermöglicht die Nutzung von \euro{}
\usepackage{pifont} %Ermöglicht die Nutzung verschiedener Fonts, z.B.: \ding{42} , \ding{52} , \ding{56}
\usepackage{tikzsymbols} %Ermöglicht die Nutzung von u.a. Smileys und anderen Symbolen, z.B.: \Winkey , \Smiley[1.5]



%Pakete zur Darstellung von Grafiken
\usepackage{tikz} %\begin{tikzpicture} .... \end{tikzpicture}
\usepackage{pgf, pgfplots} %Ermöglicht 2D und 3D Plof-Funktion
\pgfplotsset{compat=1.17}
\usetikzlibrary{arrows,%Pfeile
			calc,%Rechnen in Koordinaten, z.B. mit ($(1,2)+5*(0,1)$) auf Punkt (1,7) kommen
			angles, plotmarks, quotes, babel, intersections, shapes.geometric, decorations.pathreplacing, decorations.pathmorphing, automata, positioning, shadows, patterns, patterns.meta,%
			fit,%Rahmen um mehrere Knoten (UML, Automaten)
			shapes.multipart,%geteilte Knoten (UML-Klassen)
			arrows.meta,%moderne Pfeilspitzen, z.B. -{Triangle[open]}
			shapes.callouts,%Sprechblasen
			circuits.logic.US, circuits.logic.IEC%Logikgatter (digitale Schaltungen)
			} %nützliche Bibliotheken für TikZ
\usepackage{animate} %animierte Darstellungen
\usepackage{graphicx} %\includegraphics[width=..px,height=..px]{Dateiname}
\usepackage{multirow} %Zellen über mehrere Zeilen mit \multirow{n}{*}{Text}
\usepackage{array} %erweiterte Spaltentypen, z.B. >{\centering\arraybackslash}p{2cm}
\usepackage{booktabs} %schöne Tabellenlinien \toprule, \midrule, \bottomrule
\usepackage{longtable} %Tabellen über mehrere Seiten



%gelegentlich hilfreiche Zusatzpakete
\usepackage{hologo} %\hologo{LaTeX}
\usepackage{suetterl} %Sütterlin als Schriftart mit \mathfrak{Text}
\usepackage{qrcode} %Einbindung eines QR-Codes 
\usepackage{url} %Einbindung eines URL-Links
\usepackage{ccicons} %Creative Commons Symbole
\usepackage{listings} %\begin{lstlisting} ... \end{lstlisting}
\usepackage{struktex} %Struktogrammdarstellung
\usepackage{ulem} %Unterstreichung mit \uline{} ggf. mit \hfill im Inneren für Unterstrich bis zum Ende der Zeile
\usepackage{xskak} %Schach: Partien und Züge mit \newchessgame, \mainline{1. e4 e5} (lädt skak)
\usepackage{chessboard} %Schachbretter mit \chessboard[setfen=...]
\usepackage{sudoku} %Sudoku-Gitter mit \begin{sudoku} |5|3| | ... |. \end{sudoku}
\usepackage{tikzpeople} %Personen-Figuren als TikZ-Knoten, z.B. \node[alice, minimum size=2cm]{};
\usepackage{twemojis} %farbige Emojis (Twemoji, CC BY 4.0) als Vektorgrafik: \twemoji[height=1cm]{man standing}, im Text \texttwemoji{light bulb}
\usepackage[safe]{tipa} %Lautschrift mit \textipa{...}; [safe], weil tipa sonst \; \: \! \* \s umdefiniert (bricht z.B. \(\;\)) – Langformen siehe latex-lautschrift
\usepackage[european]{circuitikz} %Schaltpläne mit deutschen Symbolen: \begin{circuitikz} \draw (0,0) to[lamp] (2,0); \end{circuitikz}
\usepackage[version=4]{mhchem} %Reaktionsgleichungen mit \ce{2 H2 + O2 -> 2 H2O}
\usepackage{chemfig} %Strukturformeln mit \chemfig{H-C(-[2]H)(-[6]H)-H}, freie Elektronenpaare mit \charge
\usepackage{pgf-PeriodicTable} %Periodensystem mit \pgfPT[...], deutsche Namen mit \pgfPTsetLanguage{de}
\usepackage[all]{genealogytree} %Stammbäume mit \begin{genealogypicture} ... \end{genealogypicture}
%\usepackage{musixtex} %Notensatz – NUR für Musik-Dokumente einkommentieren: definiert u.a. \cup als Note um (bricht $A \cup B$)
%\usepackage[list=true, font=large, labelfont=bf, labelformat=brace, position=top]{subcaption}



%abgeänderte Befehle
\renewcommand{\arraystretch}{1.25} %
%\renewcommand{\familydefault}{\sfdefault} %serifenlose Schriftart
\renewcommand{\lstlistingname}{Quellcode} %Umbenennung
\renewcommand{\lstlistlistingname}{Quellcodeverzeichnis} %Umbenennung
%----------------------------------------------------------------------------------------------------%
%----------------------------------------------------------------------------------------------------%
%----------------------------------------------------------------------------------------------------%
%%%% %%%%%%                                     NEWCOMMANDS                                         %%%%%%%%%%%%
%----------------------------------------------------------------------------------------------------%
%----------------------------------------------------------------------------------------------------%
%
% Universale newcommands.tex – alle eigenen Befehle, Farben, Boxen, Kopfzeilen, TikZ- und Listing-Stile.
% Neue Befehle werden HIER ergänzt (nicht im Dokument) und unten protokolliert.
% Quelle: https://github.com/mq90/latex-skills-schule  (CC BY 4.0, Martin Quandt)
%
% Änderungsprotokoll
% 2026-11-09  zusammengeführt aus allen Handreichungen: \TikZ, Box »ausgabe«, Stile für
%             Programmablaufpläne (input, operation, condition, output) und UML (umlclass, inherit,
%             compose, aggregate, associate, depend); Listing-Stil »latex« wieder wie in den
%             Handreichungen ohne Zeilennummern, nummerierte Variante als »latexnummeriert«;
%             Umlaute und · € ° → ≤ ≥ ≠ – — „ “ … × ² ³ π in allen Listings erlaubt (\lstset literate); \lehrerblase, \schuelerblase; PAP-Stile mit align=center (Zeilenumbruch mit \\);
%             \ipa{...} für Lautschrift; \koordinatensystem[Schritt]{xmin}{xmax}{ymin}{ymax};
%             TikZ-Stile punkt, punktoffen, baumdiagramm, wkeit
% 2026-11-09  TikZ-Stil baumdiagrammvertikal (Baumdiagramm von oben nach unten)
% 2026-11-09  Formularfelder für ausfüllbare PDFs: \feld, \zeichen, \kreuz, \antwort, einheitliches
%             Aussehen über \DefaultOptionsofText/CheckBox/Radio/PopdownBox (nur wirksam in einer Form-Umgebung)
% 2026-11-11  Akzentfarben akzentblau, akzentrot, akzentgruen, akzentorange, akzentlila (Karten,
%             Naturwissenschaften, Schulorganisation); Box »hinweis« mit frei wählbarem Titel
% 2026-11-11  Umgebung »passend«: verkleinert den Inhalt auf \linewidth, aber nur wenn er breiter ist
% 2026-11-11  \oz{+II}{Cu}: Oxidationszahl über ein Elementsymbol (Chemie, auch innerhalb \ce{... $\oz{0}{Zn}$ ...})
% 2026-11-22  Farbige Schachbretter: \schachbrettfarben{braun|blau|gruen|grau|klassisch} (Standard: braun),
%             Farben brett...hell/brett...dunkel und brettmarke (braucht LSBC5 aus der header.tex)
%








%----------------------------------------------------------------------------------------------------%
%%%%%%%%%%%                                   Mathematische Kürzel                                       %%%%%%%%%%
%----------------------------------------------------------------------------------------------------%


%vereinfachte Schreibweise für den Displaystyle
\newcommand{\dis}[1]{\displaystyle{#1}}


%Zahlenbereiche
\newcommand{\N}{\mathbb{N}} %natürliche Zahlen
\newcommand{\Z}{\mathbb{Z}} %ganze Zahlen
\newcommand{\Q}{\mathbb{Q}} %rationale Zahlen
\newcommand{\R}{\mathbb{R}} %reelle Zahlen
\newcommand{\C}{\mathbb{C}} %komplexe Zahlen
\newcommand{\cP}{\mathcal{P}} %Potenzmenge


%Abkürzungen für Analysis
\newcommand{\eps}{\varepsilon} %Kürzel für Epsilon
\newcommand{\limn}{\underset{n\longrightarrow \infty}{\lim}} %Grenzwert für n gegen Unendlich
\newcommand{\limh}{\underset{h\longrightarrow 0}{\lim}} %Grenzwert für h gegen Null
\newcommand{\limxto}[1]{\underset{x\longrightarrow #1}{\lim}} %beliebiger funktionaler Grenzwert x gegen ....
\newcommand{\dx}{\text{\;d} x} %dx Operator für Integrale
\renewcommand{\d}[1]{\text{\;d}#1} %Operator für beliebige Integrationsvariable
\newcommand{\norm}[1]{\left\lVert#1\right\rVert} %\norm{...} für die Normdarstellung ||...||


%Abkürzungen für Analytische Geometrie
\newcommand{\vektor}[3]{\begin{tiny}\begin{pmatrix} \text{\begin{scriptsize}$#1$\end{scriptsize}}\\ \text{\begin{scriptsize}$#2$\end{scriptsize}}\\ \text{\begin{scriptsize}$#3$\end{scriptsize}}\end{pmatrix}\end{tiny}} %dreidimensionaler Vektor
\newcommand{\punktvektor}[2]{\begin{tiny}\begin{pmatrix} \text{\begin{scriptsize}$#1$\end{scriptsize}}\\ \text{\begin{scriptsize}$#2$\end{scriptsize}}\end{pmatrix}\end{tiny}} %zweidimensionaler Vektor


%Abkürzungen für Stochastik
\newcommand{\E}{\mathbb{E}} %Erwartungswert
\newcommand{\V}{\mathbb{V}} %Varianz


%Umgebung für gegeben und gesucht
\newcommand{\gegges}[2]{%
	\begin{multicols}{2}
		\begin{enumerate}
			\item[geg:] #1
			\item[ges:] #2
		\end{enumerate}
	\end{multicols}
}
%----------------------------------------------------------------------------------------------------%









%----------------------------------------------------------------------------------------------------%
%%%%%%%%%%%%                             Verschiedene  hilfreiche Kürzel                                %%%%%%%%%
%----------------------------------------------------------------------------------------------------%


%römische Darstellung einer Zahl \uproman{2021} ergibt MMXXI
\newcommand{\uproman}[1]{\uppercase\expandafter{\romannumeral#1}}


%Automatisch erstellte Lücke wie die Länge des Textes mit \luecke{<Text>}
\newcommand{\luecke}[1]{\hspace{0,1cm}\underline{\hspace{0,25cm}\hphantom{#1}\hspace{0,25cm}}\hspace{0,1cm}}


%fett-kursive Schrift für neu eingeführte Begriffe
\newcommand{\fk}[1]{\textbf{\textit{#1}}}


%Schwierigkeitsstufen für Aufgaben mit \leicht , \mittel , \schwer
\newcommand{\leicht}{\smash{\raisebox{-1pt}{{\color{black}\rule{4pt}{8pt}}\,{\color{lightgray}\rule{4pt}{8pt}}\,{\color{lightgray}\rule{4pt}{8pt}}}}\,}
\newcommand{\mittel}{\smash{\raisebox{-1pt}{{\color{black}\rule{4pt}{8pt}}\,{\color{black}\rule{4pt}{8pt}}\,{\color{lightgray}\rule{4pt}{8pt}}}}\,}
\newcommand{\schwer}{\smash{\raisebox{-1pt}{{\color{black}\rule{4pt}{8pt}}\,{\color{black}\rule{4pt}{8pt}}\,{\color{black}\rule{4pt}{8pt}}}}\,}
%----------------------------------------------------------------------------------------------------%










%----------------------------------------------------------------------------------------------------%
%%%%%%%                                            Verschiedene Farben                                                      %%%%%%%
%----------------------------------------------------------------------------------------------------%


%Formeln in herkömmlichen Farben schreiben
\newcommand{\frot}[1]{\text{\textcolor{red}{$#1$}}} % rote Formel für Lösungen
\newcommand{\fgruen}[1]{\text{\textcolor{green}{$#1$}}} % grüne Formel für Lösungen
\newcommand{\fblau}[1]{\text{\textcolor{blue}{$#1$}}} % blaue Formel für Lösungen
\newcommand{\fmagenta}[1]{\text{\textcolor{magenta}{$#1$}}} % magenta-farbige Formel für Lösungen
\newcommand{\forange}[1]{\text{\textcolor{orange}{$#1$}}} % orangene Formel für Lösungen
\newcommand{\fcyan}[1]{\text{\textcolor{cyan}{$#1$}}} % cyan-farbige Formel für Lösungen
\newcommand{\folive}[1]{\text{\textcolor{olive}{$#1$}}} % olive-farbige Formel für Lösungen


%Custom made colors
\definecolor{plum}{rgb}{0.36078, 0.20784, 0.4}
\definecolor{chameleon}{rgb}{0.30588, 0.60392, 0.023529}
\definecolor{cornflower}{rgb}{0.12549, 0.29020, 0.52941}
\definecolor{scarlet}{rgb}{0.8, 0, 0}
\definecolor{brick}{rgb}{0.64314, 0, 0}
\definecolor{sunrise}{rgb}{0.80784, 0.36078, 0}
\definecolor{lightblue}{rgb}{0.15,0.35,0.75}
\definecolor{carolina}{rgb}{153, 186, 221}

%Akzentfarben für Grafiken, Karten und Diagramme (kräftig, gut unterscheidbar; aufgehellt z.B. akzentblau!18)
\definecolor{akzentblau}{HTML}{2563EB}
\definecolor{akzentrot}{HTML}{DC2626}
\definecolor{akzentgruen}{HTML}{16A34A}
\definecolor{akzentorange}{HTML}{EA580C}
\definecolor{akzentlila}{HTML}{7C3AED}
%----------------------------------------------------------------------------------------------------%









%----------------------------------------------------------------------------------------------------%
%%%%%%%%                            Kopfzeilen für Leistungsüberprüfungen                              %%%%%%%%%
%----------------------------------------------------------------------------------------------------%


%Definition Kopfzeile für Leistungskontrolle mit \lk{Klasse}{Thema}
\newcommand{\lk}[2]{%
	\pagestyle{fancy}
	\fancyhead[L]{%
	\begin{tabular}[b]{l}
		\\Name:\\Punkte:\\Note:
	\end{tabular} 
	}
	\fancyhead[R]{%
	\begin{tabular}[b]{r}
		\\Leistungskontrolle\\Klasse #1\\ \today
	\end{tabular} 
	}
	\fancyhead[C]{%
	\textbf{\large{#2}}\\
	}
	\fancyfoot[C]{}
	\renewcommand{\headrulewidth}{1pt}
	\fancyheadoffset{1cm}
}


%Definition Kopfzeile für Klassenarbeit mit \ka{Klasse}{Thema}
\newcommand{\ka}[2]{%
	\pagestyle{fancy}
	\fancyhead[L]{%
	\begin{tabular}[b]{l}
		\\Name:\\Punkte:\\Note:
	\end{tabular} 
	}
	\fancyhead[R]{%
	\begin{tabular}[b]{r}
		\\Klassenarbeit\\Klasse #1\\ \today
	\end{tabular} 
	}
	\fancyhead[C]{%
	\textbf{\large{#2}}\\
	}
	\fancyfoot[C]{}
	\renewcommand{\headrulewidth}{1pt}
	\fancyheadoffset{1cm}
}


%Definition Kopfzeile für Klausur mit \kl{Klasse}{Thema}
\newcommand{\kl}[2]{%
	\pagestyle{fancy}
	\fancyhead[L]{%
	\begin{tabular}[b]{l}
		\\Name:\\Punkte:\\Note:
	\end{tabular} 
	}
	\fancyhead[R]{%
	\begin{tabular}[b]{r}
		\\Klausur\\Klasse #1\\ \today
	\end{tabular} 
	}
	\fancyhead[C]{%
	\textbf{\large{#2}}\\
	}
	\fancyfoot[C]{}
	\renewcommand{\headrulewidth}{1pt}
	\fancyheadoffset{1cm}
}


%Entfernt den Header für Leistungskontrollen, Klassenarbeiten oder Klausuren
\newcommand{\clearfancyheader}{%
	\fancyhf{}
	\renewcommand{\headrulewidth}{0pt}
	\renewcommand{\headheight}{0pt}
}


%Definition Kopfzeile für Tägliche Übungen mit \tu{Klasse}{Thema}
\newcommand{\tu}[2]{%
	\pagestyle{fancy}
	\fancyhead[L]{%
	\begin{tabular}[b]{l}
		\\Name:\\Punkte:\\Note:
	\end{tabular} 
	}
	\fancyhead[R]{%
	\begin{tabular}[b]{r}
		\\Tägliche Übung\\Klasse #1\\ \today
	\end{tabular} 
	}
	\fancyhead[C]{%
	\textbf{\large{#2}}\\
	}
	\fancyfoot[C]{}
	\renewcommand{\headrulewidth}{1pt}
	\fancyheadoffset{1cm}
}
%----------------------------------------------------------------------------------------------------%








%----------------------------------------------------------------------------------------------------%
%%%%%%%%%                                     Colorboxes für Hervorhebungen                              %%%%%%%%%
%----------------------------------------------------------------------------------------------------%


 %Colorboxes für Lösungen mit \begin{solution}<Text>\end{solution}
\newtcolorbox{solution}{%
        	lower separated=false,%
	colback=gray!5!white,%
	colframe=gray!75!black,%
        	fonttitle=\bfseries,%
	title=Lösung:,%
      	enhanced,%
	breakable,%
        	attach boxed title to top left={xshift=0.5cm,yshift=-2mm}%
}


%Colorboxes für das chronologische Vorgehen 
\newtcolorbox{todo}{%
        	lower separated=false,%
	colback=yellow!5!white,%
	colframe=orange!75!black,%
        	enhanced,%
	breakable%
}


%Colorboxes für Merksätze
\newtcolorbox{merke}{%
        	lower separated=false,%
	colback=red!5!white,%
	colframe=red!75!black,%
        	fonttitle=\bfseries,%
	title=Merke:,%
      	enhanced,%
	breakable,%
        	attach boxed title to top left={xshift=0.5cm,yshift=-2mm}%
}


%Hinweis-Box mit frei wählbarem Titel: \begin{hinweis}[Lizenz] ... \end{hinweis} (Standardtitel: Hinweis)
\newtcolorbox{hinweis}[1][Hinweis]{%
	enhanced,%
	breakable,%
	sharp corners,%
	colback=akzentblau!5,%
	colframe=akzentblau,%
	boxrule=0pt,%
	leftrule=3pt,%
	fonttitle=\bfseries,%
	coltitle=akzentblau!85!black,%
	attach title to upper={\ },%
	title={#1:}%
}


%Colorbox2 für Merksätze 
\newtcolorbox{merke2}{%
        	lower separated=false,%
        	colback=red!40,%
       	title=Merke,%
       	colframe=black,%
        	fonttitle=\bfseries,%
        	colbacktitle=red,%
        	coltitle=black,%
        	enhanced,%
        	breakable,%
        	attach boxed title to top left={xshift=0.5cm,yshift=-2mm},%
        	title=Merke:%
}
%----------------------------------------------------------------------------------------------------%










%----------------------------------------------------------------------------------------------------%
%%%%%%%%%%%%%                                         MC-Fragen                                    %%%%%%%%%%%%%
%----------------------------------------------------------------------------------------------------%

%Definition des Kästchens zum ankreuzen
\newcommand{\casket}{\LARGE{$\square$}}


%Umgebung für horizontale MC-Fragen
\newenvironment{mcquestionh}[1]{%
	#1%
	\begin{center}\begin{tabular}{>{\centering\arraybackslash}p{2.5cm}  >{\centering\arraybackslash}p{2.5cm} >{\centering\arraybackslash}p{2.5cm} >{\centering\arraybackslash}p{2.5cm} >{\centering\arraybackslash}p{2.5cm}}
	\casket & \casket & \casket & \casket & \casket\medskip\\
	}
	%hier wird die fehlende Tabellenzeile angegeben. Aufruf erfolgt durch:
	%\begin{mcquestionh}{<Frage>}
	%   	\answerh{<Antwort>} & \answerh{<Antwort>} & \answerh{<Antwort>} & \answerh{<Antwort>} & \answerh{<Antwort>}
	%\end{mcquestionh}
	{
	\\
	\end{tabular}\end{center}
}


%Antwort für horizontale MC-Fragen mit \answerh{<Antwort>}
\newcommand{\answerh}[1]{%
	\parbox{2.5cm}{\centering{#1}}
}


%Umgebung für vertikale MC-Fragen
\newenvironment{mcquestionv}[1]{%
    	#1%
    	\begin{enumerate}[\raisebox{-3pt}{\casket}]
    	}
    	%hier wird das fehlende \item angegeben werden. Aufruf erfolgt durch:
    	%\begin{mcquestionv}{<Frage>}
    	%   	\answerv{<Antwort>}
    	%  	\answerv{<Antwort>}
    	%   	\answerv{<Antwort>}
    	%   	\answerv{<Antwort>}
    	%   	\answerv{<Antwort>}
    	%\end{mcquestionv}
    	{
    	\end{enumerate}
}


%Antwort für vertikale MC-Frage mit \answerv{<Antwort>}
\newcommand{\answerv}[1]{%
	\item \hspace{0.1cm} #1
}
%----------------------------------------------------------------------------------------------------%









%----------------------------------------------------------------------------------------------------%
%%%%%%%%%%%                                          Tikz Settings                                          %%%%%%%%%%%%
%----------------------------------------------------------------------------------------------------%


%Koordinatensystem im Standardstil der Mathe-Arbeitsblätter (Gitter gray!30, Achsen -latex, Beschriftung \tiny)
%\koordinatensystem[Schritt]{xmin}{xmax}{ymin}{ymax};
%             TikZ-Stile punkt, punktoffen, baumdiagramm, wkeit
% 2026-11-09  TikZ-Stil baumdiagrammvertikal (Baumdiagramm von oben nach unten) – ganze Zahlen, Bereich muss 0 enthalten,
%Striche und Zahlen bei allen Vielfachen von Schritt (Standard 2). Innerhalb von tikzpicture[scale=0.42] verwenden.
\newcommand{\koordinatensystem}[5][2]{%
	\foreach \kosyx in {#2,...,#3} \draw[gray!30, thin] (\kosyx,#4) -- (\kosyx,#5);
	\foreach \kosyy in {#4,...,#5} \draw[gray!30, thin] (#2,\kosyy) -- (#3,\kosyy);
	\draw[-latex, thick] ({#2-0.3},0) -- ({#3+0.5},0) node[right] {$x$};
	\draw[-latex, thick] (0,{#4-0.3}) -- (0,{#5+0.5}) node[above] {$y$};
	\foreach \kosyx in {#2,...,#3} {%
		\pgfmathtruncatemacro{\kosyrest}{mod(abs(\kosyx),#1)}%
		\ifnum\kosyrest=0 \ifnum\kosyx=0 \else
			\draw (\kosyx,0.1) -- (\kosyx,-0.1) node[below, font=\tiny] {$\kosyx$};
		\fi\fi}
	\foreach \kosyy in {#4,...,#5} {%
		\pgfmathtruncatemacro{\kosyrest}{mod(abs(\kosyy),#1)}%
		\ifnum\kosyrest=0 \ifnum\kosyy=0 \else
			\draw (0.1,\kosyy) -- (-0.1,\kosyy) node[left, font=\tiny] {$\kosyy$};
		\fi\fi}
}


%optionale Eigenschaft für ein  Tikzpicture, damit dieses bündig am oberen Rand einer enumerate-Auflistung steht (benötigt tikzlibrary calc!!!)
\tikzset{
  topline/.style={baseline={([yshift=#1]current bounding box.north)}},
  topline/.default=-\ht\strutbox
}


%Punkte in Graphen und Figuren: \node[punkt, label=above:{$P$}] at (2,3) {}; (offen: punktoffen)
%Als Knoten, damit die Größe nicht mit scale=0.42 mitschrumpft (circle (3pt) würde mitskaliert)
\tikzset{
  punkt/.style={circle, fill, inner sep=0pt, minimum size=4pt},
  punktoffen/.style={circle, draw, fill=white, inner sep=0pt, minimum size=4pt}
}


%Baumdiagramme (Stochastik) von links nach rechts: \begin{tikzpicture}[baumdiagramm] \node[punkt]{} child{...}
%Wahrscheinlichkeiten an den Ästen: edge from parent node[wkeit, above] {$0{,}6$} (oberer Ast above, unterer below)
\tikzset{
  baumdiagramm/.style={grow'=right, sloped, level distance=2.6cm,
    level 1/.style={sibling distance=3cm}, level 2/.style={sibling distance=1.5cm},
    level 3/.style={sibling distance=0.75cm}, edge from parent/.style={draw, thin}},
  wkeit/.style={font=\scriptsize, inner sep=2pt}
}


%Baumdiagramme von oben nach unten: \begin{tikzpicture}[baumdiagrammvertikal] \node[punkt]{} child{...}
%erstes child = linker Ast; Wahrscheinlichkeit links mit node[wkeit, left], rechts mit node[wkeit, right]
\tikzset{
  baumdiagrammvertikal/.style={grow=down, level distance=1.6cm,
    level 1/.style={sibling distance=4.8cm}, level 2/.style={sibling distance=2.4cm},
    level 3/.style={sibling distance=1.2cm}, edge from parent/.style={draw, thin}}
}
%----------------------------------------------------------------------------------------------------%









%----------------------------------------------------------------------------------------------------%
%%%%%%%%%%%                                     Handreichungen                                        %%%%%%%%%%%
%----------------------------------------------------------------------------------------------------%


%Schriftzug TikZ im Fließtext (Leerzeichen danach ist eingebaut)
\newcommand{\TikZ}{Ti\textit{k}Z }


%Oxidationszahl über ein Elementsymbol: $\oz{+II}{Cu}$, in \ce{...} zwischen Dollarzeichen: \ce{$\oz{0}{Zn}$ + ...}
\newcommand{\oz}[2]{\overset{\scriptscriptstyle\mathrm{#1}}{\mathrm{\ce{#2}}}}


%Umgebung passend: Inhalt (z.B. eine große tikzpicture) höchstens so breit wie die Zeile setzen.
%Ist er schmaler, bleibt er unverändert. \begin{passend} ... \end{passend} (Inhalt nicht als Argument, # erlaubt)
\newsavebox{\passendbox}
\newenvironment{passend}{\begin{lrbox}{\passendbox}}{\end{lrbox}%
	\ifdim\wd\passendbox>\linewidth\resizebox{\linewidth}{!}{\usebox{\passendbox}}\else\usebox{\passendbox}\fi}


%Box für die Ausgabe unter einem Code-Beispiel: erst lstlisting[style=latex], darunter \begin{ausgabe} ... \end{ausgabe}
\newtcolorbox{ausgabe}{%
	colback=blue!5,%
	colframe=blue!50!black,%
	title=Ausgabe,%
	fonttitle=\bfseries%
}
%----------------------------------------------------------------------------------------------------%









%----------------------------------------------------------------------------------------------------%
%%%%%%%%%%%                                        Lautschrift                                         %%%%%%%%%%%
%----------------------------------------------------------------------------------------------------%


%Lautschrift in eckigen Klammern mit \ipa{"a:b@nt} – schaltet das babel-Kürzel " vorher aus (sonst wird "a zu ä)
\makeatletter
\newcommand{\ipa}{\begingroup\shorthandoff{"}\ipa@aux}
\newcommand{\ipa@aux}[1]{[\textipa{#1}]\endgroup}
\makeatother
%----------------------------------------------------------------------------------------------------%









%----------------------------------------------------------------------------------------------------%
%%%%%%%%%%%                                       Sprechblasen                                         %%%%%%%%%%%
%----------------------------------------------------------------------------------------------------%


%Sprechblase mit Zipfel links unten (Lehrkraft) mit \lehrerblase[TikZ-Optionen]{Text} (benötigt shapes.callouts)
\newcommand{\lehrerblase}[2][]{%
  \tikz[baseline=(B.base)]{%
    \node[rectangle callout, rounded corners=6pt, callout relative pointer={(-13mm,-7mm)},
      draw, thick, fill=blue!6, align=left, inner sep=3.5mm, #1] (B) {#2};}%
}


%Sprechblase mit Zipfel rechts unten (Schülerin/Schüler) mit \schuelerblase[TikZ-Optionen]{Text}
\newcommand{\schuelerblase}[2][]{%
  \tikz[baseline=(B.base)]{%
    \node[rectangle callout, rounded corners=6pt, callout relative pointer={(13mm,-7mm)},
      draw, thick, fill=green!6, align=left, inner sep=3.5mm, #1] (B) {#2};}%
}
%----------------------------------------------------------------------------------------------------%









%----------------------------------------------------------------------------------------------------%
%%%%%%%%%%%                            TikZ-Stile für Programmablaufpläne                            %%%%%%%%%%%
%----------------------------------------------------------------------------------------------------%


%Ein-/Ausgabe als Parallelogramm, Operation als Rechteck, Verzweigung als Raute (benötigt shapes.geometric)
%align=center, damit \\ im Knotentext eine neue Zeile ergibt (ohne align wird \\ stillschweigend verschluckt)
\tikzstyle{input}=[trapezium, trapezium left angle=60, trapezium right angle=120, fill=green!20, draw=green, align=center]
\tikzstyle{operation}=[rectangle, fill=orange!20, draw=orange, align=center]
\tikzstyle{condition}=[diamond, fill=magenta!20, draw=magenta, align=center]
\tikzstyle{output}=[trapezium, trapezium left angle=60, trapezium right angle=120, fill=red!20, draw=red, align=center]
%----------------------------------------------------------------------------------------------------%









%----------------------------------------------------------------------------------------------------%
%%%%%%%%%%%                            TikZ-Stile für UML-Klassendiagramme                           %%%%%%%%%%%
%----------------------------------------------------------------------------------------------------%


%Klasse mit drei Teilen: Name, Attribute, Methoden (benötigt shapes.multipart)
\tikzstyle{umlclass} = [rectangle split, rectangle split parts=3, rectangle split part align={center,left,left}, draw, thick, text width=4cm, minimum height=2cm]


%Beziehungen (benötigt arrows.meta)
\tikzstyle{inherit} = [-{Triangle[open]}, thick]	%Vererbung
\tikzstyle{compose} = [-{Diamond}, thick]	%Komposition
\tikzstyle{aggregate} = [-{Diamond[open]}, thick]	%Aggregation
\tikzstyle{associate} = [-, thick]	%Assoziation
\tikzstyle{depend} = [-, thick, dashed]	%Abhängigkeit
%----------------------------------------------------------------------------------------------------%








%----------------------------------------------------------------------------------------------------%
%%%%%%%%%%%                                         Listing Settings                                       %%%%%%%%%%%%
%----------------------------------------------------------------------------------------------------%


%Umlaute und häufige Sonderzeichen in ALLEN Listings (Python, Java, ...), sonst Fehler »Invalid UTF-8 byte sequence«
\lstset{literate=%
	{ä}{{\"a}}1 {ö}{{\"o}}1 {ü}{{\"u}}1 {Ä}{{\"A}}1 {Ö}{{\"O}}1 {Ü}{{\"U}}1 {ß}{{\ss}}1%
	{·}{{$\cdot$}}1 {€}{{\euro}}1 {°}{{\degree}}1 {→}{{$\rightarrow$}}2 {≤}{{$\leq$}}1 {≥}{{$\geq$}}1 {≠}{{$\neq$}}1%
	{–}{{--}}1 {—}{{---}}1 {„}{{\glqq}}1 {“}{{\grqq}}1 {…}{{\dots}}3 {×}{{$\times$}}1 {²}{{$^2$}}1 {³}{{$^3$}}1 {π}{{$\pi$}}1%
}


%Einstellung der Codeumgebung für LaTeX-Code-Beispiele (Stil der Handreichungen: ohne Zeilennummern)
%Verwendung: \begin{lstlisting}[style=latex] ... \end{lstlisting}, darunter \begin{ausgabe} ... \end{ausgabe}
\lstdefinestyle{latex}{%
  	language=TeX,%
  	frame=single,%
  	backgroundcolor=\color{yellow!20},%
  	rulecolor=\color{lightgray!40},%
  	breaklines=true,%
  	basicstyle=\scriptsize\ttfamily,%
  	keywordstyle=\bfseries\color{purple!40!black},%
  	commentstyle=\itshape\color{green!40!black},%
  	identifierstyle=\color{blue},%
	stringstyle=\color{black},%
	captionpos=b,%
	tabsize=2,%
	showspaces=false,%
	showstringspaces=false,%
  	showtabs=false,%
	deletekeywords={document},%
	morekeywords={begin}%
}


%Variante mit Zeilennummern: \begin{lstlisting}[style=latexnummeriert]
\lstdefinestyle{latexnummeriert}{%
	style=latex,%
	numbers=left,%
	numberstyle=\tiny\color{gray!90!black},%
  	numbersep=-5pt%
}
%----------------------------------------------------------------------------------------------------%






%----------------------------------------------------------------------------------------------------%
%%%%%%%%%%%                          Formularfelder (ausfüllbare PDFs)                          %%%%%%%%%%%
%----------------------------------------------------------------------------------------------------%
% Nur innerhalb von \begin{Form} ... \end{Form} (hyperref) verwendbar. Für Optionsfelder (radio)
% im Dokument zusätzlich \hypersetup{unicode=false}, sonst bekommen sie kaputte Namen.
% Jedes Feld steckt in einer eigenen Box (\raisebox/\hbox): Ohne sie übernimmt pdfTeX die Höhe
% der umgebenden Zeile, in Tabellen werden die Felder dann so hoch wie die ganze Zeile.

%Aussehen aller Felder: blauer Rand, hellblauer Hintergrund, Schrift 11 pt, wird mitgedruckt
\renewcommand{\DefaultOptionsofText}{print, bordercolor={0 0.353 0.612}, backgroundcolor={0.94 0.96 1}, charsize=11pt}
\renewcommand{\DefaultOptionsofCheckBox}{print, bordercolor={0 0.353 0.612}, backgroundcolor={0.94 0.96 1}, charsize=11pt}
\renewcommand{\DefaultOptionsofRadio}{print, notoggletooff, bordercolor={0 0.353 0.612}, backgroundcolor={0.94 0.96 1}, charsize=11pt}
\renewcommand{\DefaultOptionsofPopdownBox}{print, bordercolor={0 0.353 0.612}, backgroundcolor={0.94 0.96 1}, charsize=11pt}
\renewcommand{\DefaultHeightofText}{0.6cm}

%\feld[Breite]{Name} – einzeiliges Textfeld, sitzt auf der Grundlinie (Standardbreite 2,4 cm)
\newcommand{\feld}[2][2.4cm]{\raisebox{-0.17cm}{\TextField[name=#2, width=#1]{}}}

%\zeichen{Name} – Feld für genau ein Zeichen, zentriert (z. B. E/V/A, Buchstabe einer Reihenfolge)
\newcommand{\zeichen}[1]{\raisebox{-0.12cm}{\TextField[name=#1, width=0.8cm, height=0.5cm, maxlen=1, align=1]{}}}

%\kreuz{Name} – Ankreuzkästchen
\newcommand{\kreuz}[1]{\raisebox{-0.08cm}{\CheckBox[name=#1, width=0.42cm, height=0.42cm]{}}}

%\antwort[Höhe]{Name} – mehrzeiliges Antwortfeld über die ganze Zeilenbreite (Standardhöhe 3,2 cm ≈ 4 Zeilen)
\newcommand{\antwort}[2][3.2cm]{\par\vspace{0.15cm}\noindent\hbox{\TextField[name=#2, multiline=true, width=\linewidth, height=#1]{}}}
%----------------------------------------------------------------------------------------------------%




%----------------------------------------------------------------------------------------------------%
%%%%%%%%%%%                                   Schach (farbige Bretter)                                   %%%%%%%%%%
%----------------------------------------------------------------------------------------------------%
% Die Standardschrift von skak schraffiert die dunklen Felder; Figur und Schraffur sind dort ein
% Zeichen und lassen sich nicht getrennt einfärben. Deshalb: Schrift skaknew in der Kodierung LSBC5
% (zeichnet nur die Figuren, weiße Figuren weiß gefüllt), die Felder malt chessboard als Flächen
% darunter. Eigene Markierungen mit colorbackfields liegen über den Feldern und unter den Figuren.

%Feldfarben der Brettschemata (hell/dunkel) und Markierungsfarbe für den letzten Zug
\definecolor{brettblauhell}{HTML}{EEF3FD} %Schema blau, helle Felder
\definecolor{brettblaudunkel}{HTML}{A9B4CC} %Schema blau, dunkle Felder
\definecolor{brettbraunhell}{HTML}{F0D9B5} %Schema braun (klassisches Holzbrett), helle Felder
\definecolor{brettbraundunkel}{HTML}{B58863} %Schema braun, dunkle Felder
\definecolor{brettgruenhell}{HTML}{EEEED2} %Schema gruen (Turnierbrett), helle Felder
\definecolor{brettgruendunkel}{HTML}{769656} %Schema gruen, dunkle Felder
\definecolor{brettgrauhell}{HTML}{F2F2F2} %Schema grau (für Schwarz-Weiß-Druck), helle Felder
\definecolor{brettgraudunkel}{HTML}{B4B4B4} %Schema grau, dunkle Felder
\definecolor{brettmarke}{HTML}{F6D55C} %Start- und Zielfeld eines Zuges hervorheben: pgfstyle=color, color=brettmarke, colorbackfields={e2,e4}

% Setzt die Voreinstellungen von \setchessboard zurück – eigene \setchessboard{...} deshalb DANACH schreiben.
%\schachbrettfarben{Schema} – Farben aller folgenden Schachbretter: braun (Standard), blau, gruen, grau oder klassisch (schraffiert wie in skak)
\makeatletter
\let\brett@grundeinstellung\board@val@savedkeylist
\newcommand{\schachbrettfarben}[1]{%
	\let\board@val@savedkeylist\brett@grundeinstellung
	\ifthenelse{\equal{#1}{klassisch}}{}%
	{\setchessboard{boardfontfamily=skaknew, boardfontencoding=LSBC5,
		piececolor=black, piecemaskcolor=white, setfontcolors,
		pgfstyle=color, opacity=1, color=brett#1hell, colorbackboard, color=brett#1dunkel, colorblackbackfields,
		color=black, labelfont=\sffamily\color{black!65}}}}
\makeatother
\schachbrettfarben{braun}
%----------------------------------------------------------------------------------------------------%


% Mathematisches Wörterbuch Deutsch – Englisch mit Lautschrift
% Skills: latex-grundstil, latex-lautschrift (tipa, \ipa), latex-grundlagen (Formelsatz), latex-geometrie,
%         latex-funktionsgraphen, latex-venn-diagramme
% Kompilieren mit pdflatex (zweimal, wegen Inhaltsverzeichnis und Seitenzahlen im Register).

%----------------------------------------------------------------------------------------------------%
% Dokumentspezifische Befehle (nur für dieses Wörterbuch, deshalb lokal und nicht in newcommands.tex)
%----------------------------------------------------------------------------------------------------%

% \eintrag{Marke}{Artikel}{Begriff}{english term}{Lautschrift EN}{Zeichen/Skizze}
% Schaltet das babel-Kürzel " aus, BEVOR die Argumente gelesen werden – dadurch darf die Lautschrift
% wie bei \ipa{...} mit " (Hauptbetonung) und "" (Nebenbetonung) geschrieben werden.
\makeatletter
\newcommand\eintrag{\begingroup\shorthandoff{"}\eintrag@aux}
\newcommand\eintrag@aux[6]{\endgroup
	#2 \textbf{#3}\phantomsection\label{e:#1}
	& \textbf{#4}\newline{\small\color{black!70}\ipa{#5}}
	& #6 \\ \midrule}
\makeatother

% Tabellenkopf für jedes Kapitel (wird auf jeder Folgeseite wiederholt)
\newenvironment{woerter}{%
	\arrayrulecolor{black!20}%
	\begin{longtable}{>{\raggedright\arraybackslash}m{5.3cm} >{\raggedright\arraybackslash}m{6.1cm} >{\centering\arraybackslash}m{4.6cm}}
		\rowcolor{cornflower!15}\textbf{Deutsch} & \textbf{Englisch} & \textbf{Zeichen \,/\, Skizze} \\ \midrule
		\endhead}{%
	\end{longtable}\arrayrulecolor{black}}

% Kleine Skizzen (Gestaltung wie im Skill latex-geometrie: Figur blau very thick, Maße/Hervorhebung rot)
\tikzset{
	fig/.style={blue, very thick},
	mass/.style={red, thick},
	hilf/.style={dashed},
	achse/.style={-{Stealth[length=4pt]}, thin},
}
\newcommand{\skizze}[1]{\tikz[baseline=(current bounding box.center), scale=0.45, line cap=round,
	line join=round, every node/.style={font=\scriptsize}, outer sep=0pt]{#1}}

% Grundbegriffe
\newcommand{\skPunkt}{\skizze{\fill[blue] (0,0) circle (2.5pt); \node[above right] at (0,0) {$P$};}}
\newcommand{\skGerade}{\skizze{\draw[fig] (0,0) -- (3.6,0.8); \node[below right] at (3.2,0.7) {$g$};}}
\newcommand{\skStrecke}{\skizze{\draw[fig] (0,0) -- (3.2,0.7); \fill (0,0) circle (2pt) (3.2,0.7) circle (2pt);
	\node[below] at (0,0) {$A$}; \node[below] at (3.2,0.7) {$B$};}}
\newcommand{\skStrahl}{\skizze{\draw[fig, -{Stealth}] (0,0) -- (3.6,0.8); \fill (0,0) circle (2pt);
	\node[below] at (0,0) {$S$};}}
\newcommand{\skParallel}{\skizze{\draw[fig] (0,0) -- (3.5,0.6) node[right] {$g$};
	\draw[fig] (0,-0.9) -- (3.5,-0.3) node[right] {$h$};}}
\newcommand{\skSenkrecht}{\skizze{\draw[fig] (-1.6,0) -- (1.6,0) node[right] {$g$};
	\draw[fig] (0,-0.5) -- (0,1.8) node[right] {$h$};
	\draw[mass] (0.45,0) arc[start angle=0, end angle=90, radius=0.45]; \fill[red] (45:0.24) circle (1pt);}}
\newcommand{\skAbstand}{\skizze{\draw[fig] (-0.4,0) -- (3.2,0) node[right] {$g$}; \fill (1.4,1.6) circle (2pt) node[above] {$P$};
	\draw[mass, hilf] (1.4,1.6) -- node[right] {$d$} (1.4,0);
	\draw[mass] (1.75,0) arc[start angle=0, end angle=90, radius=0.35]; \fill[red] (1.4,0) ++(45:0.19) circle (0.8pt);}}

% Winkel
\newcommand{\skWinkel}[1]{\skizze{\draw[mass] (0.7,0) arc[start angle=0, end angle=#1, radius=0.7];
	\draw[fig] (2.3,0) -- (0,0) -- (#1:2.3); \node at ({#1/2}:1.15) {$\alpha$};}}
\newcommand{\skRechterWinkel}{\skizze{\draw[mass] (0.6,0) arc[start angle=0, end angle=90, radius=0.6];
	\fill[red] (45:0.32) circle (1pt); \draw[fig] (2.3,0) -- (0,0) -- (0,2.3);}}
\newcommand{\skVollwinkel}{\skizze{\draw[mass] (0.6,0) arc[start angle=0, end angle=355, radius=0.6];
	\draw[fig] (0,0) -- (2.3,0); \fill[blue] (0,0) circle (1.5pt); \node at (135:1.0) {$\alpha$};}}
\newcommand{\skNebenwinkel}{\skizze{\draw[mass] (0.7,0) arc[start angle=0, end angle=60, radius=0.7];
	\draw[mass] (60:0.55) arc[start angle=60, end angle=180, radius=0.55];
	\draw[fig] (-2.2,0) -- (2.2,0); \draw[fig] (0,0) -- (60:2.2);
	\node at (28:1.15) {$\alpha$}; \node at (125:1.0) {$\beta$};}}
\newcommand{\skScheitelwinkel}{\skizze{\draw[mass] (143.1:0.6) arc[start angle=143.1, end angle=201.8, radius=0.6];
	\draw[mass] (323.1:0.6) arc[start angle=323.1, end angle=381.8, radius=0.6];
	\draw[fig] (-2,-0.8) -- (2,0.8); \draw[fig] (-1.6,1.2) -- (1.6,-1.2);
	\node at (172:1.05) {$\alpha$}; \node at (352:1.1) {$\alpha'$};}}
\newcommand{\skStufenwinkel}{\skizze{\draw[mass] (0.5,0) arc[start angle=0, end angle=53.13, radius=0.5];
	\draw[mass] (1.7,1.6) arc[start angle=0, end angle=53.13, radius=0.5];
	\draw[fig] (-1.2,0) -- (2.6,0); \draw[fig] (-1.2,1.6) -- (2.6,1.6); \draw[fig] (-0.6,-0.8) -- (1.8,2.4);
	\node at (26:0.85) {$\alpha$}; \node at ($(1.2,1.6)+(26:0.85)$) {$\beta$};}}
\newcommand{\skWechselwinkel}{\skizze{\draw[mass] (0.5,0) arc[start angle=0, end angle=53.13, radius=0.5];
	\draw[mass] (0.7,1.6) arc[start angle=180, end angle=233.13, radius=0.5];
	\draw[fig] (-1.2,0) -- (2.6,0); \draw[fig] (-1.2,1.6) -- (2.6,1.6); \draw[fig] (-0.6,-0.8) -- (1.8,2.4);
	\node at (26:0.85) {$\alpha$}; \node at ($(1.2,1.6)+(206:0.85)$) {$\beta$};}}

% Abbildungen
\newcommand{\skSymmetrieachse}{\skizze{\draw[fig] (0,0) -- (2.6,0) -- (1.3,2.1) -- cycle;
	\draw[mass, hilf] (1.3,-0.5) -- (1.3,2.5);}}
\newcommand{\skSpiegelung}{\skizze{\draw[mass, hilf] (0,-0.4) -- (0,1.8);
	\draw[fig] (-2,0) -- (-0.5,0.3) -- (-1.5,1.4) -- cycle;
	\draw[fig, blue!45] (2,0) -- (0.5,0.3) -- (1.5,1.4) -- cycle;}}
\newcommand{\skVerschiebung}{\skizze{\draw[fig] (0,0) -- (1.3,0) -- (0.3,1.1) -- cycle;
	\draw[fig, blue!45] (2,0.6) -- (3.3,0.6) -- (2.3,1.7) -- cycle;
	\draw[mass, -{Stealth}] (0.3,1.1) -- (2.3,1.7);}}
\newcommand{\skDrehung}{\skizze{\draw[fig] (0.6,0) -- (1.9,0) -- (0.8,0.9) -- cycle;
	\draw[fig, blue!45] (0,0.6) -- (0,1.9) -- (-0.9,0.8) -- cycle;
	\fill (0,0) circle (2pt) node[below left] {$Z$};
	\draw[mass, -{Stealth}] (10:2.2) arc[start angle=10, end angle=80, radius=2.2];}}

% Dreiecke
\newcommand{\skDreieck}{\skizze{\draw[fig] (0,0) -- (3,0) -- (1,2) -- cycle;
	\node[below left] at (0,0) {$A$}; \node[below right] at (3,0) {$B$}; \node[above] at (1,2) {$C$};}}
\newcommand{\skRechtwinklig}{\skizze{\draw[mass] (0.5,0) arc[start angle=0, end angle=90, radius=0.5];
	\fill[red] (45:0.27) circle (1pt); \draw[fig] (0,0) -- (3,0) -- (0,2) -- cycle;}}
\newcommand{\skGleichschenklig}{\skizze{\draw[fig] (0,0) -- (2.6,0) -- (1.3,2.2) -- cycle;
	\draw[mass] (0.52,1.176) -- (0.78,1.024); \draw[mass] (1.82,1.024) -- (2.08,1.176);}}
\newcommand{\skGleichseitig}{\skizze{\draw[fig] (0,0) -- (2.6,0) -- (1.3,2.252) -- cycle;
	\draw[mass] (1.3,-0.15) -- (1.3,0.15); \draw[mass] (0.52,1.201) -- (0.78,1.051);
	\draw[mass] (1.82,1.051) -- (2.08,1.201);}}
\newcommand{\skHypotenuse}{\skizze{\draw[fig] (0,2) -- (0,0) -- (3,0); \draw[red, very thick] (3,0) -- (0,2);
	\draw[mass] (0.5,0) arc[start angle=0, end angle=90, radius=0.5]; \fill[red] (45:0.27) circle (1pt);
	\node[above right] at (1.5,1) {$c$};}}
\newcommand{\skKathete}{\skizze{\draw[fig] (3,0) -- (0,2); \draw[red, very thick] (0,2) -- (0,0) -- (3,0);
	\draw[mass] (0.5,0) arc[start angle=0, end angle=90, radius=0.5]; \fill[red] (45:0.27) circle (1pt);
	\node[left] at (0,1) {$b$}; \node[below] at (1.5,0) {$a$};}}
\newcommand{\skPythagoras}{\skizze{\draw[fig] (0,0) -- (2,0) -- (0,1.5) -- cycle;
	\draw[fig, fill=blue!10] (0,0) rectangle (2,-2); \draw[fig, fill=blue!10] (0,0) rectangle (-1.5,1.5);
	\draw[fig, fill=red!10] (2,0) -- (3.5,2) -- (1.5,3.5) -- (0,1.5) -- cycle;
	\node at (1,-1) {$a^2$}; \node at (-0.75,0.75) {$b^2$}; \node at (1.75,1.75) {$c^2$};}}
\newcommand{\skHoehe}{\skizze{\draw[fig] (0,0) -- (3.2,0) -- (1,2) -- cycle;
	\draw[mass, hilf] (1,2) -- node[right] {$h$} (1,0);
	\draw[mass] (1.35,0) arc[start angle=0, end angle=90, radius=0.35]; \fill[red] (1,0) ++(45:0.19) circle (0.8pt);}}

% Vierecke, Vielecke
\newcommand{\skViereck}{\skizze{\draw[fig] (0,0) -- (2.8,0.3) -- (2.3,1.9) -- (0.4,1.5) -- cycle;}}
\newcommand{\skQuadrat}{\skizze{\draw[fig] (0,0) rectangle (1.9,1.9);}}
\newcommand{\skRechteck}{\skizze{\draw[fig] (0,0) rectangle (3.2,1.6);}}
\newcommand{\skParallelogramm}{\skizze{\draw[fig] (0,0) -- (2.6,0) -- (3.4,1.5) -- (0.8,1.5) -- cycle;}}
\newcommand{\skRaute}{\skizze{\draw[fig] (0,0) -- (1.4,-1) -- (2.8,0) -- (1.4,1) -- cycle;}}
\newcommand{\skTrapez}{\skizze{\draw[fig] (0,0) -- (3.2,0) -- (2.4,1.5) -- (0.6,1.5) -- cycle;}}
\newcommand{\skDrachen}{\skizze{\draw[fig] (1,0) -- (2,1.6) -- (1,2.3) -- (0,1.6) -- cycle;
	\draw[mass, hilf] (1,-0.3) -- (1,2.6);}}
\newcommand{\skFuenfeck}{\skizze{\draw[fig] (90:1.2) -- (162:1.2) -- (234:1.2) -- (306:1.2) -- (18:1.2) -- cycle;}}
\newcommand{\skSechseck}{\skizze{\draw[fig] (0:1.2) -- (60:1.2) -- (120:1.2) -- (180:1.2) -- (240:1.2) -- (300:1.2) -- cycle;}}
\newcommand{\skFlaeche}{\skizze{\filldraw[gray, opacity=0.5] (0,0) rectangle (3,1.6);
	\draw[fig] (0,0) rectangle (3,1.6); \node at (1.5,0.8) {$A$};}}
\newcommand{\skUmfangViereck}{\skizze{\draw[red, very thick] (0,0) rectangle (3,1.6); \node at (1.5,0.8) {$u$};}}

% Kreis
\newcommand{\skKreis}{\skizze{\draw[fig] (0,0) circle (1.2); \fill (0,0) circle (1.5pt);}}
\newcommand{\skMittelpunkt}{\skizze{\draw[fig] (0,0) circle (1.2); \fill[red] (0,0) circle (2.5pt) node[below] {$M$};}}
\newcommand{\skRadius}{\skizze{\draw[fig] (0,0) circle (1.2); \fill (0,0) circle (1.5pt) node[below] {$M$};
	\draw[mass] (0,0) -- node[above] {$r$} (30:1.2);}}
\newcommand{\skDurchmesser}{\skizze{\draw[fig] (0,0) circle (1.2); \fill (0,0) circle (1.5pt);
	\draw[mass] (200:1.2) -- (20:1.2); \node[above] at (100:0.25) {$d$};}}
\newcommand{\skKreisumfang}{\skizze{\draw[red, very thick] (0,0) circle (1.2); \fill (0,0) circle (1.5pt);}}
\newcommand{\skSektor}{\skizze{\filldraw[gray, opacity=0.5] (0,0) -- (1.2,0) arc[start angle=0, end angle=70, radius=1.2] -- cycle;
	\draw[fig] (0,0) circle (1.2); \draw[mass] (1.2,0) -- (0,0) -- (70:1.2);}}
\newcommand{\skSehne}{\skizze{\draw[fig] (0,0) circle (1.2); \draw[mass] (160:1.2) -- (40:1.2);}}
\newcommand{\skTangente}{\skizze{\draw[fig] (0,0) circle (1.1); \draw[mass] ($(50:1.1)+(140:1.5)$) -- ($(50:1.1)+(-40:1.5)$);
	\fill[red] (50:1.1) circle (2.5pt);}}

% Körper (Schrägbild: Tiefe unter 45° und halb so lang, verdeckte Kanten gestrichelt)
\newcommand{\skWuerfel}{\skizze{\def\d{0.57}
	\draw[fig, hilf, thick] (0,0) -- (\d,\d) -- (1.6+\d,\d) (\d,\d) -- (\d,1.6+\d);
	\draw[fig] (0,0) rectangle (1.6,1.6); \draw[fig] (1.6,0) -- (1.6+\d,\d) -- (1.6+\d,1.6+\d) -- (\d,1.6+\d) -- (0,1.6)
		(1.6,1.6) -- (1.6+\d,1.6+\d);}}
\newcommand{\skQuader}{\skizze{\def\d{0.6}
	\draw[fig, hilf, thick] (0,0) -- (\d,\d) -- (2.6+\d,\d) (\d,\d) -- (\d,1.3+\d);
	\draw[fig] (0,0) rectangle (2.6,1.3); \draw[fig] (2.6,0) -- (2.6+\d,\d) -- (2.6+\d,1.3+\d) -- (\d,1.3+\d) -- (0,1.3)
		(2.6,1.3) -- (2.6+\d,1.3+\d);}}
\newcommand{\skPrisma}{\skizze{
	\draw[fig, hilf, thick] (0,0) -- (1.4,0.7) -- (3,0.7) (1.4,0.7) -- (2.2,2);
	\draw[fig] (0,0) -- (1.6,0) -- (0.8,1.3) -- cycle;
	\draw[fig] (1.6,0) -- (3,0.7) -- (2.2,2) -- (0.8,1.3);}}
\newcommand{\skZylinder}{\skizze{
	\draw[fig, hilf, thick] (1,0) arc[start angle=0, end angle=180, x radius=1, y radius=0.3];
	\draw[fig] (-1,0) arc[start angle=180, end angle=360, x radius=1, y radius=0.3];
	\draw[fig] (0,1.9) ellipse[x radius=1, y radius=0.3];
	\draw[fig] (-1,0) -- (-1,1.9) (1,0) -- (1,1.9);}}
\newcommand{\skKegel}{\skizze{
	\draw[fig, hilf, thick] (1,0) arc[start angle=0, end angle=180, x radius=1, y radius=0.3];
	\draw[fig] (-1,0) arc[start angle=180, end angle=360, x radius=1, y radius=0.3];
	\draw[fig] (-1,0) -- (0,2.2) -- (1,0);}}
\newcommand{\skPyramide}{\skizze{
	\draw[fig, hilf, thick] (0,0) -- (0.8,0.6) -- (2.8,0.6) (0.8,0.6) -- (1.4,2.3);
	\draw[fig] (0,0) -- (2,0) -- (2.8,0.6) -- (1.4,2.3) -- cycle (2,0) -- (1.4,2.3);}}
\newcommand{\skKugel}{\skizze{\draw[fig] (0,0) circle (1.1);
	\draw[fig, hilf, thick] (1.1,0) arc[start angle=0, end angle=180, x radius=1.1, y radius=0.3];
	\draw[fig] (-1.1,0) arc[start angle=180, end angle=360, x radius=1.1, y radius=0.3];}}
\newcommand{\skKante}{\skizze{\def\d{0.57}
	\draw[fig, hilf, thick] (0,0) -- (\d,\d) -- (1.6+\d,\d) (\d,\d) -- (\d,1.6+\d);
	\draw[fig] (0,0) rectangle (1.6,1.6); \draw[fig] (1.6,0) -- (1.6+\d,\d) -- (1.6+\d,1.6+\d) -- (\d,1.6+\d) -- (0,1.6);
	\draw[red, very thick] (1.6,1.6) -- (1.6+\d,1.6+\d);}}
\newcommand{\skEcke}{\skizze{\def\d{0.57}
	\draw[fig, hilf, thick] (0,0) -- (\d,\d) -- (1.6+\d,\d) (\d,\d) -- (\d,1.6+\d);
	\draw[fig] (0,0) rectangle (1.6,1.6); \draw[fig] (1.6,0) -- (1.6+\d,\d) -- (1.6+\d,1.6+\d) -- (\d,1.6+\d) -- (0,1.6)
		(1.6,1.6) -- (1.6+\d,1.6+\d); \fill[red] (1.6+\d,1.6+\d) circle (2.5pt);}}
\newcommand{\skSeitenflaeche}{\skizze{\def\d{0.57}
	\draw[fig, hilf, thick] (0,0) -- (\d,\d) -- (1.6+\d,\d) (\d,\d) -- (\d,1.6+\d);
	\fill[gray, opacity=0.5] (0,0) rectangle (1.6,1.6);
	\draw[fig] (0,0) rectangle (1.6,1.6); \draw[fig] (1.6,0) -- (1.6+\d,\d) -- (1.6+\d,1.6+\d) -- (\d,1.6+\d) -- (0,1.6)
		(1.6,1.6) -- (1.6+\d,1.6+\d);}}
\newcommand{\skNetz}{\skizze{\begin{scope}[scale=0.62]
	\foreach \x/\y in {1/0, 0/1, 1/1, 2/1, 3/1, 1/2} {\draw[fig, fill=blue!8] (\x,\y) rectangle ++(1,1);}
	\end{scope}}}

% Funktionen (Achsen dünn mit Pfeil, Graph blau, Hervorhebungen rot)
\newcommand{\achsen}[4]{\draw[achse] (#1,0) -- (#2,0) node[below left] {$x$}; \draw[achse] (0,#3) -- (0,#4) node[below left] {$y$};}
\newcommand{\skKoordinaten}{\skizze{\draw[black!25, very thin] (-1.5,-1.5) grid (2.5,2); \achsen{-1.6}{2.7}{-1.6}{2.2}
	\fill[red] (1,1.5) circle (2pt) node[right] {$P$};}}
\newcommand{\skGraph}{\skizze{\achsen{-1.8}{2.4}{-0.6}{2.4}
	\begin{scope}\clip (-1.8,-0.6) rectangle (2.2,2.2); \draw[fig, domain=-1.8:2.2, samples=40] plot (\x, {0.3*\x*\x*\x - 0.6*\x + 0.9});\end{scope}}}
\newcommand{\skXAchse}{\skizze{\draw[achse, red, thick] (-1.8,0) -- (2.4,0) node[below left] {$x$}; \draw[achse] (0,-1.2) -- (0,1.6) node[below left] {$y$};}}
\newcommand{\skParabel}{\skizze{\achsen{-1.9}{1.9}{-0.6}{2.4}
	\begin{scope}\clip (-1.9,-0.6) rectangle (1.9,2.2); \draw[fig, domain=-1.8:1.8, samples=40] plot (\x, {0.7*\x*\x});\end{scope}}}
\newcommand{\skScheitelpunkt}{\skizze{\achsen{-0.6}{3}{-0.6}{2.4}
	\begin{scope}\clip (-0.6,-0.6) rectangle (3,2.2); \draw[fig, domain=-0.3:3, samples=40] plot (\x, {0.7*(\x-1.4)*(\x-1.4)+0.4});\end{scope}
	\fill[red] (1.4,0.4) circle (2.5pt) node[below] {$S$};}}
\newcommand{\skNullstelle}{\skizze{\achsen{-1.9}{1.9}{-1.2}{1.8}
	\begin{scope}\clip (-1.9,-1.2) rectangle (1.9,1.6); \draw[fig, domain=-1.8:1.8, samples=40] plot (\x, {0.6*\x*\x-0.9});\end{scope}
	\fill[red] (-1.2247,0) circle (2.5pt) (1.2247,0) circle (2.5pt);}}
\newcommand{\skLinear}{\skizze{\achsen{-0.8}{3}{-0.6}{2.4}
	\begin{scope}\clip (-0.8,-0.6) rectangle (3,2.2); \draw[fig, domain=-0.8:3] plot (\x, {0.5*\x+0.5});\end{scope}}}
\newcommand{\skSteigung}{\skizze{\achsen{-0.8}{3}{-0.6}{2.4}
	\begin{scope}\clip (-0.8,-0.6) rectangle (3,2.2); \draw[fig, domain=-0.8:3] plot (\x, {0.5*\x+0.5});\end{scope}
	\draw[mass] (0.6,0.8) -- node[below] {$\Delta x$} (2.2,0.8) -- node[right] {$\Delta y$} (2.2,1.6);}}
\newcommand{\skAchsenabschnitt}{\skizze{\achsen{-0.8}{3}{-0.6}{2.4}
	\begin{scope}\clip (-0.8,-0.6) rectangle (3,2.2); \draw[fig, domain=-0.8:3] plot (\x, {0.5*\x+0.5});\end{scope}
	\fill[red] (0,0.5) circle (2.5pt) node[left] {$n$};}}
\newcommand{\skSchnittpunkt}{\skizze{\draw[fig] (-0.2,0) -- (3,1.6); \draw[fig] (-0.2,1.8) -- (3,0.2);
	\fill[red] (1.6,0.9) circle (2.5pt) node[above=2pt] {$S$};}}
\newcommand{\skExponential}{\skizze{\achsen{-1.9}{1.9}{-0.6}{2.4}
	\begin{scope}\clip (-1.9,-0.6) rectangle (1.9,2.2); \draw[fig, domain=-1.9:1.9, samples=40] plot (\x, {0.5*2^\x});\end{scope}}}
\newcommand{\skLogarithmus}{\skizze{\achsen{-0.6}{3.2}{-1.4}{1.6}
	\begin{scope}\clip (-0.6,-1.4) rectangle (3.2,1.4); \draw[fig, domain=0.2:3.2, samples=40] plot (\x, {ln(\x)});\end{scope}}}
\newcommand{\sinusgraph}{\achsen{-0.3}{3.5}{-1.2}{1.4}\draw[fig, domain=0:3.3, samples=50] plot (\x, {0.9*sin(2*\x r)});}
\newcommand{\skSinus}{\skizze{\sinusgraph}}
\newcommand{\skKosinus}{\skizze{\achsen{-0.3}{3.5}{-1.2}{1.4}\draw[fig, domain=0:3.3, samples=50] plot (\x, {0.9*cos(2*\x r)});}}
\newcommand{\skHochpunkt}{\skizze{\sinusgraph \fill[red] (0.785,0.9) circle (2.5pt) node[above] {$H$};}}
\newcommand{\skTiefpunkt}{\skizze{\sinusgraph \fill[red] (2.356,-0.9) circle (2.5pt) node[below] {$T$};}}
\newcommand{\skWendepunkt}{\skizze{\sinusgraph \fill[red] (1.571,0) circle (2.5pt) node[above right] {$W$};}}
\newcommand{\skMonotonSteigend}{\skizze{\achsen{-0.4}{3}{-0.4}{2.2}\draw[fig, domain=0:2.7, samples=40] plot (\x, {0.25*\x*\x+0.2});}}
\newcommand{\skMonotonFallend}{\skizze{\achsen{-0.4}{3}{-0.4}{2.2}\draw[fig, domain=0.25:2.8, samples=40] plot (\x, {0.5/\x+0.1});}}
\newcommand{\skAchsensymmetrisch}{\skizze{\draw[mass, hilf] (0,-0.6) -- (0,2.2); \achsen{-1.9}{1.9}{-0.6}{2.4}
	\begin{scope}\clip (-1.9,-0.6) rectangle (1.9,2.2); \draw[fig, domain=-1.8:1.8, samples=40] plot (\x, {0.3*\x*\x*\x*\x - 0.9*\x*\x + 1});\end{scope}}}
\newcommand{\skPunktsymmetrisch}{\skizze{\achsen{-1.9}{1.9}{-1.4}{1.6}
	\begin{scope}\clip (-1.9,-1.4) rectangle (1.9,1.4); \draw[fig, domain=-1.8:1.8, samples=40] plot (\x, {0.4*\x*\x*\x - 0.9*\x});\end{scope}
	\fill[red] (0,0) circle (2.5pt);}}
\newcommand{\skAsymptote}{\skizze{\achsen{-0.4}{3.2}{-0.4}{2.4}\draw[mass, hilf] (-0.3,0.4) -- (3.1,0.4);
	\draw[fig, domain=0.3:3.1, samples=40] plot (\x, {0.5/\x+0.4});}}
\newcommand{\skIntegral}{\skizze{\fill[gray, opacity=0.5] (0.5,0) -- plot[domain=0.5:2.5, samples=30] (\x, {0.3*(\x-1.5)*(\x-1.5)+0.6}) -- (2.5,0) -- cycle;
	\achsen{-0.4}{3.2}{-0.4}{2.2}\draw[fig, domain=0.1:2.9, samples=40] plot (\x, {0.3*(\x-1.5)*(\x-1.5)+0.6});
	\node[below] at (0.5,0) {$a$}; \node[below] at (2.5,0) {$b$};}}
\newcommand{\skTangenteGraph}{\skizze{\achsen{-0.4}{3.2}{-0.4}{2.4}
	\begin{scope}\clip (-0.4,-0.4) rectangle (3.2,2.2); \draw[fig, domain=0:3, samples=40] plot (\x, {0.25*\x*\x});
	\draw[mass] (0.4,-0.4) -- (3.2,2.4);\end{scope} \fill[red] (2,1) circle (2pt);}}
\newcommand{\skFolge}{\skizze{\achsen{-0.3}{3.4}{-0.3}{2.2}
	\foreach \n in {1,...,10} {\fill[blue] ({0.3*\n},{1.6-1.2/\n}) circle (1.6pt);}
	\draw[mass, hilf] (0,1.6) -- (3.2,1.6);}}

% Vektoren
\newcommand{\skVektor}{\skizze{\draw[mass, very thick, -{Stealth}] (0,0) -- node[above left] {$\vec{a}$} (2.6,1.3);}}
\newcommand{\skOrtsvektor}{\skizze{\achsen{-0.4}{3}{-0.4}{2.2}\draw[mass, very thick, -{Stealth}] (0,0) -- (2.2,1.5);
	\fill (2.2,1.5) circle (2pt) node[above] {$P$}; \node[below left] at (0,0) {$O$};}}
\newcommand{\skKreuzprodukt}{\skizze{\draw[fig, -{Stealth}] (0,0) -- (2.2,0) node[below] {$\vec{a}$};
	\draw[fig, -{Stealth}] (0,0) -- (1,-0.8) node[below] {$\vec{b}$};
	\draw[mass, very thick, -{Stealth}] (0,0) -- (0,2) node[right] {$\vec{a}\times\vec{b}$};}}
\newcommand{\skEbene}{\skizze{\filldraw[blue!10, draw=blue, very thick] (0,0) -- (2.8,0) -- (3.6,1.4) -- (0.8,1.4) -- cycle;
	\node at (1.8,0.7) {$E$};}}

% Mengen (wie im Skill latex-venn-diagramme: Schnitt über \clip)
\newcommand{\venn}{\draw[fig] (0,0) circle (1) (1.2,0) circle (1); \node at (-0.5,0) {$A$}; \node at (1.7,0) {$B$};}
\newcommand{\skSchnittmenge}{\skizze{\begin{scope}\clip (0,0) circle (1); \fill[gray, opacity=0.5] (1.2,0) circle (1);\end{scope}\venn}}
\newcommand{\skVereinigung}{\skizze{\fill[gray, opacity=0.5] (0,0) circle (1) (1.2,0) circle (1);\venn}}
\newcommand{\skTeilmenge}{\skizze{\draw[fig] (0,0) ellipse[x radius=1.6, y radius=1.1]; \draw[fig, fill=gray!40] (-0.5,-0.1) circle (0.55);
	\node at (-0.5,-0.1) {$A$}; \node at (0.8,0.3) {$B$};}}

% Stochastik
\newcommand{\skBaum}{\skizze{\draw[fig, thick] (0,0) -- (1.4,0.9) (0,0) -- (1.4,-0.9) (1.4,0.9) -- (2.8,1.35) (1.4,0.9) -- (2.8,0.45)
	(1.4,-0.9) -- (2.8,-0.45) (1.4,-0.9) -- (2.8,-1.35);
	\foreach \p in {(0,0),(1.4,0.9),(1.4,-0.9),(2.8,1.35),(2.8,0.45),(2.8,-0.45),(2.8,-1.35)} {\fill[blue] \p circle (1.6pt);}}}
\newcommand{\skVierfeldertafel}{{\scriptsize\begin{tabular}{c|cc|c}
	& $B$ & $\overline{B}$ & \\ \hline $A$ & & & \\ $\overline{A}$ & & & \\ \hline & & & $1$\end{tabular}}}
\newcommand{\skWertetabelle}{{\scriptsize\begin{tabular}{c|c|c|c|c} $x$ & $0$ & $1$ & $2$ & $3$ \\ \hline $y$ & $1$ & $3$ & $5$ & $7$\end{tabular}}}
\newcommand{\skSaeulen}{\skizze{\draw[achse] (-0.2,0) -- (3.4,0); \draw[achse] (0,-0.2) -- (0,2.2);
	\foreach \k/\h in {1/0.4, 2/1.2, 3/1.8, 4/1.2, 5/0.4} {\filldraw[fill=blue!25, draw=blue, thick] ({0.55*\k-0.2},0) rectangle ++(0.4,\h);}}}

% Ober- und Untersumme (Ergänzung aus Teil II)
\newcommand{\rechtecksumme}[2]{\skizze{\foreach \i in {0,...,3} {
		\pgfmathsetmacro{\xl}{0.4+0.6*\i} \pgfmathsetmacro{\xr}{1.0+0.6*\i}
		\pgfmathsetmacro{\h}{0.18*(#1)*(#1)+0.3}
		\filldraw[fill=#2!30, draw=#2] (\xl,0) rectangle (\xr,\h);}
	\achsen{-0.3}{3.2}{-0.3}{2.4}\draw[fig, domain=0:2.9, samples=30] plot (\x, {0.18*\x*\x+0.3});}}
\newcommand{\skUntersumme}{\rechtecksumme{\xl}{gray}}
\newcommand{\skObersumme}{\rechtecksumme{\xr}{red}}

%----------------------------------------------------------------------------------------------------%
% Ergänzungen: Geometrie, Körper, analytische Geometrie
%----------------------------------------------------------------------------------------------------%
\newcommand{\skScheitel}{\skizze{\draw[fig] (2.3,0) -- (0,0) -- (50:2.3); \fill[red] (0,0) circle (3pt) node[below] {$S$};}}
\newcommand{\skSchenkel}{\skizze{\draw[fig] (0,0) -- (50:2.3); \draw[red, very thick] (0,0) -- (2.3,0); \fill (0,0) circle (1.5pt);}}
\newcommand{\skKomplement}{\skizze{\draw[mass] (0.8,0) arc[start angle=0, end angle=35, radius=0.8];
	\draw[mass] (35:1.1) arc[start angle=35, end angle=90, radius=1.1];
	\draw[fig] (2.3,0) -- (0,0) -- (0,2.3); \draw[fig] (0,0) -- (35:2.3);
	\node at (17:1.25) {$\alpha$}; \node at (62:1.5) {$\beta$};}}
\newcommand{\skEntgegen}{\skizze{\draw[mass] (0.5,0) arc[start angle=0, end angle=53.13, radius=0.5];
	\draw[mass] ($(1.2,1.6)+(233.13:0.45)$) arc[start angle=233.13, end angle=360, radius=0.45];
	\draw[fig] (-1.2,0) -- (2.6,0); \draw[fig] (-1.2,1.6) -- (2.6,1.6); \draw[fig] (-0.6,-0.8) -- (1.8,2.4);
	\node at (26:0.85) {$\alpha$}; \node at ($(1.2,1.6)+(300:0.8)$) {$\beta$};}}
\newcommand{\skMittelsenkrechte}{\skizze{\draw[fig] (0,0) -- (3,0); \fill (0,0) circle (2pt) (3,0) circle (2pt);
	\draw[mass] (1.5,-0.8) -- (1.5,1.4); \draw[mass] (1.85,0) arc[start angle=0, end angle=90, radius=0.35];
	\fill[red] (1.5,0) ++(45:0.19) circle (0.8pt); \draw[mass] (0.75,-0.13) -- (0.75,0.13) (2.25,-0.13) -- (2.25,0.13);}}
\newcommand{\skWinkelhalbierende}{\skizze{\draw[mass] (0.8,0) arc[start angle=0, end angle=30, radius=0.8];
	\draw[mass] (30:0.95) arc[start angle=30, end angle=60, radius=0.95];
	\draw[fig] (2.3,0) -- (0,0) -- (60:2.3); \draw[mass, hilf] (0,0) -- (30:2.6);}}
\newcommand{\skInnenwinkel}{\skizze{\draw[mass] (0.55,0) arc[start angle=0, end angle=63.43, radius=0.55];
	\draw[mass] ($(3,0)+(135:0.55)$) arc[start angle=135, end angle=180, radius=0.55];
	\draw[mass] ($(1,2)+(243.43:0.5)$) arc[start angle=243.43, end angle=315, radius=0.5];
	\draw[fig] (0,0) -- (3,0) -- (1,2) -- cycle;}}
\newcommand{\skAussenwinkel}{\skizze{\draw[mass] (3.6,0) arc[start angle=0, end angle=135, radius=0.6];
	\draw[fig] (0,0) -- (3,0) -- (1,2) -- cycle; \draw[fig, hilf, thick] (3,0) -- (4.3,0);
	\node at ($(3,0)+(60:0.98)$) {$\delta$};}}
\newcommand{\skBasis}{\skizze{\draw[fig] (0,0) -- (1.3,2.2) -- (2.6,0); \draw[red, very thick] (0,0) -- (2.6,0);}}
\newcommand{\skBasiswinkel}{\skizze{\draw[mass] (0.6,0) arc[start angle=0, end angle=59.4, radius=0.6];
	\draw[mass] ($(2.6,0)+(120.6:0.6)$) arc[start angle=120.6, end angle=180, radius=0.6];
	\draw[fig] (0,0) -- (2.6,0) -- (1.3,2.2) -- cycle;}}
\newcommand{\rwdreieck}{\draw[mass] (0.85,0) arc[start angle=0, end angle=30.96, radius=0.85]; \node at (14:1.2) {$\alpha$};
	\draw[black!60, thick] (3,0.4) arc[start angle=90, end angle=180, radius=0.4]; \fill[black!60] ($(3,0)+(135:0.2)$) circle (0.8pt);}
\newcommand{\skAnkathete}{\skizze{\rwdreieck \draw[fig] (0,0) -- (3,1.8) -- (3,0); \draw[red, very thick] (0,0) -- (3,0);}}
\newcommand{\skGegenkathete}{\skizze{\rwdreieck \draw[fig] (3,0) -- (0,0) -- (3,1.8); \draw[red, very thick] (3,0) -- (3,1.8);}}
\newcommand{\skSeitenhalbierende}{\skizze{\draw[fig] (0,0) -- (3,0) -- (1,2) -- cycle; \draw[mass] (1,2) -- (1.5,0);
	\draw[mass] (0.75,-0.13) -- (0.75,0.13) (2.25,-0.13) -- (2.25,0.13);}}
\newcommand{\skInkreis}{\skizze{\draw[fig] (0,0) -- (3,0) -- (1,2) -- cycle; \draw[mass] (1.204,0.744) circle (0.744);}}
\newcommand{\skUmkreis}{\skizze{\draw[mass] (1.5,0.5) circle (1.581); \draw[fig] (0,0) -- (3,0) -- (1,2) -- cycle;}}
\newcommand{\skDiagonale}{\skizze{\draw[fig] (0,0) rectangle (3,1.6); \draw[mass] (0,0) -- (3,1.6);}}
\newcommand{\skNeck}[1]{\skizze{\draw[fig] (90:1.2) \foreach \k in {1,...,#1} { -- ({90+\k*360/#1}:1.2) } -- cycle;}}
\newcommand{\skSekante}{\skizze{\draw[fig] (0,0) circle (1.2); \draw[mass] ($(150:1.2)!-0.35!(30:1.2)$) -- ($(150:1.2)!1.35!(30:1.2)$);
	\fill[red] (150:1.2) circle (1.6pt) (30:1.2) circle (1.6pt);}}
\newcommand{\skBeruehrpunkt}{\skizze{\draw[fig] (0,0) circle (1.1); \draw[fig, thick] ($(50:1.1)+(140:1.5)$) -- ($(50:1.1)+(-40:1.5)$);
	\fill[red] (50:1.1) circle (3pt) node[above right] {$B$};}}
\newcommand{\skKreisbogen}{\skizze{\draw[fig] (0,0) circle (1.2); \draw[red, line width=1.8pt] (20:1.2) arc[start angle=20, end angle=110, radius=1.2];}}
\newcommand{\skKreisabschnitt}{\skizze{\fill[gray, opacity=0.5] (150:1.2) arc[start angle=150, end angle=30, radius=1.2] -- cycle;
	\draw[fig] (0,0) circle (1.2); \draw[mass] (150:1.2) -- (30:1.2);}}
\newcommand{\skKreisflaeche}{\skizze{\filldraw[fill=gray!50, draw=blue, very thick] (0,0) circle (1.2);}}

% Körper
\newcommand{\wuerfelbasis}{\def\d{0.57}\draw[fig, hilf, thick] (0,0) -- (\d,\d) -- (1.6+\d,\d) (\d,\d) -- (\d,1.6+\d);
	\draw[fig] (0,0) rectangle (1.6,1.6); \draw[fig] (1.6,0) -- (1.6+\d,\d) -- (1.6+\d,1.6+\d) -- (\d,1.6+\d) -- (0,1.6) (1.6,1.6) -- (1.6+\d,1.6+\d);}
\newcommand{\quaderbasis}{\def\d{0.6}\draw[fig, hilf, thick] (0,0) -- (\d,\d) -- (2.6+\d,\d) (\d,\d) -- (\d,1.3+\d);
	\draw[fig] (0,0) rectangle (2.6,1.3); \draw[fig] (2.6,0) -- (2.6+\d,\d) -- (2.6+\d,1.3+\d) -- (\d,1.3+\d) -- (0,1.3) (2.6,1.3) -- (2.6+\d,1.3+\d);}
\newcommand{\zylinderbasis}{\draw[fig, hilf, thick] (1,0) arc[start angle=0, end angle=180, x radius=1, y radius=0.3];
	\draw[fig] (-1,0) arc[start angle=180, end angle=360, x radius=1, y radius=0.3];
	\draw[fig] (0,1.9) ellipse[x radius=1, y radius=0.3]; \draw[fig] (-1,0) -- (-1,1.9) (1,0) -- (1,1.9);}
\newcommand{\skFlaechendiagonale}{\skizze{\wuerfelbasis \draw[red, very thick] (0,0) -- (1.6,1.6);}}
\newcommand{\skRaumdiagonale}{\skizze{\wuerfelbasis \draw[red, very thick] (0,0) -- (1.6+\d,1.6+\d);}}
\newcommand{\skGrundflaeche}{\skizze{\quaderbasis \fill[red, opacity=0.4] (0,0) -- (2.6,0) -- (2.6+\d,\d) -- (\d,\d) -- cycle;}}
\newcommand{\skDeckflaeche}{\skizze{\quaderbasis \fill[red, opacity=0.4] (0,1.3) -- (2.6,1.3) -- (2.6+\d,1.3+\d) -- (\d,1.3+\d) -- cycle;}}
\newcommand{\skMantel}{\skizze{\fill[red, opacity=0.35] (-1,0) arc[start angle=180, end angle=360, x radius=1, y radius=0.3] -- (1,1.9)
	arc[start angle=0, end angle=-180, x radius=1, y radius=0.3] -- cycle; \zylinderbasis}}
\newcommand{\skPyramidenstumpf}{\skizze{\draw[fig, hilf, thick] (0,0) -- (0.8,0.6) -- (2.8,0.6) (0.8,0.6) -- (1.1,1.45);
	\draw[fig] (0,0) -- (2,0) -- (2.8,0.6) -- (2.1,1.45) -- (1.1,1.45) -- (0.7,1.15) -- cycle (0.7,1.15) -- (1.7,1.15) -- (2.1,1.45) (1.7,1.15) -- (2,0);}}
\newcommand{\skKegelstumpf}{\skizze{\draw[fig, hilf, thick] (1.1,0) arc[start angle=0, end angle=180, x radius=1.1, y radius=0.3];
	\draw[fig] (-1.1,0) arc[start angle=180, end angle=360, x radius=1.1, y radius=0.3];
	\draw[fig] (0,1.5) ellipse[x radius=0.5, y radius=0.15]; \draw[fig] (-1.1,0) -- (-0.5,1.5) (1.1,0) -- (0.5,1.5);}}
\newcommand{\skTetraeder}{\skizze{\draw[fig, hilf, thick] (0,0) -- (2.6,0.3);
	\draw[fig] (0,0) -- (1.6,-0.4) -- (2.6,0.3) -- (1.3,1.9) -- cycle (1.6,-0.4) -- (1.3,1.9);}}
\newcommand{\skHohlzylinder}{\skizze{\zylinderbasis \draw[fig] (0,1.9) ellipse[x radius=0.5, y radius=0.15];
	\draw[fig, hilf, thin] (-0.5,0) -- (-0.5,1.9) (0.5,0) -- (0.5,1.9);}}
\newcommand{\skSchieferZylinder}{\skizze{\draw[fig, hilf, thick] (1,0) arc[start angle=0, end angle=180, x radius=1, y radius=0.3];
	\draw[fig] (-1,0) arc[start angle=180, end angle=360, x radius=1, y radius=0.3];
	\draw[fig] (0.8,1.8) ellipse[x radius=1, y radius=0.3]; \draw[fig] (-1,0) -- (-0.2,1.8) (1,0) -- (1.8,1.8);}}
\newcommand{\drehachse}[2]{\draw[dash dot] (0,#1) -- (0,#2); \draw[mass, -{Stealth[length=3pt]}] (-0.3,#2-0.15) arc[start angle=180, end angle=500, x radius=0.3, y radius=0.1];}
\newcommand{\skRotation}{\skizze{\filldraw[fill=gray!50, draw=blue, very thick] (0,-1.1) arc[start angle=-90, end angle=90, radius=1.1] -- cycle;
	\draw[fig, hilf, thick] (0,1.1) arc[start angle=90, end angle=270, radius=1.1]; \drehachse{-1.4}{1.6}}}
\newcommand{\skKugelsegment}{\skizze{\fill[gray, opacity=0.5] (27.04:1.1) arc[start angle=27.04, end angle=152.96, radius=1.1] -- cycle;
	\draw[fig] (0,0) circle (1.1); \draw[fig] (-0.98,0.5) arc[start angle=180, end angle=360, x radius=0.98, y radius=0.2];
	\draw[fig, hilf, thick] (0.98,0.5) arc[start angle=0, end angle=180, x radius=0.98, y radius=0.2];}}
\newcommand{\skKugelsektor}{\skizze{\fill[gray, opacity=0.5] (0,0) -- (55:1.1) arc[start angle=55, end angle=125, radius=1.1] -- cycle;
	\draw[fig] (0,0) circle (1.1); \draw[fig, thick] (55:1.1) -- (0,0) -- (125:1.1);
	\draw[fig, thick] (-0.63,0.9) arc[start angle=180, end angle=360, x radius=0.63, y radius=0.13];}}
\newcommand{\skEllipsoid}{\skizze{\draw[fig] (0,0) ellipse[x radius=1.5, y radius=0.9];
	\draw[fig, hilf, thick] (1.5,0) arc[start angle=0, end angle=180, x radius=1.5, y radius=0.3];
	\draw[fig] (-1.5,0) arc[start angle=180, end angle=360, x radius=1.5, y radius=0.3];}}
\newcommand{\skParaboloid}{\skizze{\draw[fig, domain=-1.3:1.3, samples=30] plot (\x, {0.9*\x*\x});
	\draw[fig] (0,1.521) ellipse[x radius=1.3, y radius=0.22];}}
\newcommand{\skHyperboloid}{\skizze{\draw[fig, domain=-1.2:1.2, samples=30, variable=\t] plot ({0.55*cosh(\t)}, {0.9*sinh(\t)});
	\draw[fig, domain=-1.2:1.2, samples=30, variable=\t] plot ({-0.55*cosh(\t)}, {0.9*sinh(\t)});
	\draw[fig] (0,1.358) ellipse[x radius=0.996, y radius=0.2];
	\draw[fig] (-0.996,-1.358) arc[start angle=180, end angle=360, x radius=0.996, y radius=0.2];
	\draw[fig, hilf, thick] (0.996,-1.358) arc[start angle=0, end angle=180, x radius=0.996, y radius=0.2];}}

% Trigonometrie und analytische Geometrie
\newcommand{\skEinheitskreis}{\skizze{\achsen{-1.6}{1.9}{-1.5}{1.8}\draw[fig] (0,0) circle (1.1); \node[below right] at (1.05,0) {$1$};}}
\newcommand{\skBogenmass}{\skizze{\draw[fig] (1.6,0) -- (0,0) -- (60:1.6); \draw[red, very thick] (1.6,0) arc[start angle=0, end angle=60, radius=1.6];
	\node at (30:2.0) {$b$}; \node[below] at (0.8,0) {$r$}; \node at (30:0.6) {$\alpha$};}}
\newcommand{\skPolar}{\skizze{\achsen{-0.4}{3}{-0.4}{2.4}\draw[mass] (0.9,0) arc[start angle=0, end angle=36, radius=0.9];
	\draw[mass, very thick] (0,0) -- node[above left, inner sep=1pt] {$r$} (2.2,1.6); \fill (2.2,1.6) circle (2pt) node[above] {$P$};
	\node at (17:1.35) {$\varphi$};}}
\newcommand{\skKreisgleichung}{\skizze{\achsen{-1.7}{1.9}{-1.6}{1.8}\draw[fig] (0,0) circle (1.2); \draw[mass] (0,0) -- node[above left, inner sep=1pt] {$r$} (35:1.2);}}
\newcommand{\ellipsebasis}{\achsen{-2}{2.2}{-1.4}{1.6}\draw[fig] (0,0) ellipse[x radius=1.6, y radius=1];}
\newcommand{\skEllipse}{\skizze{\ellipsebasis}}
\newcommand{\skBrennpunkt}{\skizze{\ellipsebasis \fill[red] (-1.249,0) circle (2.5pt) node[below] {$F_1$} (1.249,0) circle (2.5pt) node[below] {$F_2$};}}
\newcommand{\skHalbachse}{\skizze{\ellipsebasis \draw[red, very thick] (0,0) -- node[below] {$a$} (1.6,0) (0,0) -- node[left] {$b$} (0,1);}}
\newcommand{\skHyperbel}{\skizze{\achsen{-2.2}{2.4}{-1.6}{1.8}
	\draw[fig, domain=-1.6:1.6, samples=30, variable=\t] plot ({0.7*cosh(\t)}, {0.6*sinh(\t)});
	\draw[fig, domain=-1.6:1.6, samples=30, variable=\t] plot ({-0.7*cosh(\t)}, {0.6*sinh(\t)});}}
\newcommand{\skLeitlinie}{\skizze{\achsen{-1.2}{2.4}{-1.7}{1.9}
	\draw[fig, domain=-1.5:1.5, samples=30, variable=\t] plot ({\t*\t/1.6}, {\t});
	\draw[mass, hilf] (-0.4,-1.6) -- (-0.4,1.7) node[left] {$l$}; \fill (0.4,0) circle (2pt) node[below right] {$F$};}}

%----------------------------------------------------------------------------------------------------%
% Teil B: Lesen und Sprechen (dokumentspezifisch)
%----------------------------------------------------------------------------------------------------%

% \lies{Ausdruck}{english reading}{deutsche Lesart}
% Wie \eintrag: " ist vor dem Lesen der Argumente ausgeschaltet, damit \ipa{...} mit " funktioniert.
\makeatletter
\newcommand\lies{\begingroup\shorthandoff{"}\lies@aux}
\newcommand\lies@aux[3]{\endgroup #1 & #2\newline{\color{black!60}#3} \\[4.5pt]}
\makeatother

% Liste ohne Linien und Rahmen: links der Ausdruck, rechts Englisch (schwarz) und Deutsch (grau)
\newenvironment{lesen}{%
	\begin{longtable}{@{}>{\raggedright\arraybackslash}p{4.6cm} >{\raggedright\arraybackslash}p{12.1cm}@{}}}{%
	\end{longtable}\vspace{-1.2em}}

% Kurzer Hinweis im Fließtext
\newcommand{\anmerkung}[1]{\par\smallskip\noindent{\small\sloppy\textbf{Hinweis:} #1\par}\smallskip}

% Zahlenstrahl für Intervalle: \intervall{g|o}{g|o} (g = Randpunkt gehört dazu, o = gehört nicht dazu)
\newcommand{\randpunkt}[2]{\if#1g\fill[blue] (#2,0) circle (3.2pt);\else\filldraw[fill=white, draw=blue, very thick] (#2,0) circle (3.2pt);\fi}
\newcommand{\intervall}[2]{\tikz[baseline=-0.6ex, scale=0.6]{\draw[achse] (-0.3,0) -- (5.6,0);
	\draw[blue, line width=2pt] (1,0) -- (4,0); \randpunkt{#1}{1} \randpunkt{#2}{4};
	\node[below=2pt, font=\scriptsize] at (1,0) {$a$}; \node[below=2pt, font=\scriptsize] at (4,0) {$b$};}}
\newcommand{\intervallR}{\tikz[baseline=-0.6ex, scale=0.6]{\draw[achse] (-0.3,0) -- (5.6,0);
	\draw[blue, line width=2pt, -{Stealth}] (1,0) -- (5.3,0); \randpunkt{g}{1};
	\node[below=2pt, font=\scriptsize] at (1,0) {$a$};}}

% Rechengröße mit Namen (englisch/deutsch) unter einer Klammer
\newcommand{\ub}[3]{\underbrace{\text{\Large$#1$}}_{\text{\footnotesize\shortstack{#2\\\color{black!60}#3}}}}

% Verhindert, dass eine Zwischenüberschrift allein am Seitenende steht
\newcommand{\platz}{\par\ifdim\dimexpr\pagegoal-\pagetotal\relax<5.5cm \newpage\fi}


\title{Mathematisches Wörterbuch\\[0.3em] \large Deutsch -- Englisch mit englischer Lautschrift, Zeichen und Skizzen}
\author{Martin Quandt}
\date{\today}
\pagestyle{plain}
\setlength{\tabcolsep}{4pt}

\begin{document}

\maketitle

Dieses Wörterbuch enthält 353 mathematische Fachbegriffe aus der Sekundarstufe I und II,
geordnet nach Sachgebieten. Zu jedem Begriff stehen der englische Fachausdruck mit seiner Aussprache und
-- wo es hilft -- das Zeichen, eine Formel oder eine kleine Skizze. Teil~II zeigt, wie man Zeichen,
Formeln und Rechenwege auf Englisch vorliest, mit vielen Beispielen. Am Ende finden Sie je ein
alphabetisches Register Deutsch -- Englisch und Englisch -- Deutsch.

\subsection*{Hinweise zur Lautschrift}

Die Lautschrift folgt dem Internationalen Phonetischen Alphabet (IPA) und gibt die britische
Aussprache wieder, wie sie in Schulwörterbüchern steht (dort wird \textipa{r} statt
\textipa{\textturnr} geschrieben). Der Strich \textipa{\textprimstress} steht \textit{vor} der
betonten Silbe, \textipa{\textsecstress} vor einer Silbe mit Nebenbetonung, \textipa{:} bedeutet lang.

\medskip
{\small
\begin{tabular}{>{\centering\arraybackslash}p{0.9cm} l @{\hspace{1.2cm}} >{\centering\arraybackslash}p{0.9cm} l}
	\toprule
	\textbf{Laut} & \textbf{Beispiel} & \textbf{Laut} & \textbf{Beispiel} \\ \midrule
	\textipa{@} & \textbf{a}rea \ipa{"e@ri@} & \textipa{T} & \textbf{th}eorem \ipa{"TI@r@m} \\
	\textipa{\ae} & \textbf{a}ngle \ipa{"\ae Ng\textsyllabic{l}} & \textipa{D} & less \textbf{th}an \ipa{"les D@n} \\
	\textipa{2} & n\textbf{u}mber \ipa{"n2mb@} & \textipa{S} & frac\textbf{t}ion \ipa{"fr\ae kS\textsyllabic{n}} \\
	\textipa{6} & pr\textbf{o}duct \ipa{"pr6d2kt} & \textipa{dZ} & e\textbf{dge} \ipa{edZ} \\
	\textipa{3:} & v\textbf{er}tex \ipa{"v3:teks} & \textipa{N} & a\textbf{n}gle \ipa{"\ae Ng\textsyllabic{l}} \\
	\textipa{@U} & c\textbf{o}ne \ipa{k@Un} & \textipa{\textsyllabic{l}} & circ\textbf{le} \ipa{"s3:k\textsyllabic{l}} \\
	\bottomrule
\end{tabular}}

\medskip
\textbf{Farben in den Skizzen:} Figuren sind blau, hervorgehobene Teile (der Begriff selbst, Winkel,
Maße) rot, verdeckte Kanten und Hilfslinien gestrichelt.

\newpage
{\small\setlength{\parskip}{0pt}\tableofcontents}
\newpage

\part{Wörterbuch}

%%%%  ZAHLEN UND RECHENARTEN  %%%%
\section{Zahlen und Rechenarten}

\begin{woerter}
\eintrag{1}{die}{Zahl}{number}{"n2mb@}{$1;\ 2;\ 3;\ \ldots$}
\eintrag{2}{die}{Ziffer}{digit}{"dIdZIt}{$0;\ 1;\ \ldots;\ 9$}
\eintrag{3}{die}{Ordnungszahl}{ordinal number}{"O:dIn\textsyllabic{l} "n2mb@}{1.\ 2.\ 3. \newline \footnotesize engl. 1st, 2nd, 3rd}
\eintrag{4}{der}{Index}{subscript, index}{"s2bskrIpt, "Indeks}{$x_1;\ x_2;\ \ldots;\ x_n$}
\eintrag{5}{die}{natürliche Zahl}{natural number}{"n\ae tSr@l "n2mb@}{$\N$}
\eintrag{6}{die}{ganze Zahl}{integer}{"IntIdZ@}{$\Z$}
\eintrag{7}{die}{rationale Zahl}{rational number}{"r\ae Sn@l "n2mb@}{$\Q$}
\eintrag{8}{die}{reelle Zahl}{real number}{"ri:@l "n2mb@}{$\R$}
\eintrag{9}{die}{irrationale Zahl}{irrational number}{I"r\ae S@n\textsyllabic{l} "n2mb@}{$\sqrt{2};\ \pi$}
\eintrag{10}{die}{komplexe Zahl}{complex number}{"k6mpleks "n2mb@}{$\C$ \newline $z = a + b\,\mathrm{i}$}
\eintrag{11}{die}{imaginäre Einheit}{imaginary unit}{I"m\ae dZIn@ri "ju:nIt}{$\mathrm{i}^2 = -1$}
\eintrag{12}{die}{Primzahl}{prime number}{"praIm "n2mb@}{$2;\ 3;\ 5;\ 7;\ 11;\ \ldots$}
\eintrag{13}{die}{Dezimalzahl}{decimal number}{"desIm\textsyllabic{l} "n2mb@}{$3{,}14$}
\eintrag{14}{}{addieren}{add}{\ae d}{$a + b$}
\eintrag{15}{die}{Summe}{sum}{s2m}{$\underbrace{a + b}_{\text{Summe}}$}
\eintrag{16}{der}{Summand}{summand, addend}{"s2m\ae nd, "\ae dend}{$\underbrace{a}_{\text{Summand}} + \underbrace{b}_{\text{Summand}}$}
\eintrag{17}{}{subtrahieren}{subtract}{s@b"tr\ae kt}{$a - b$}
\eintrag{18}{die}{Differenz}{difference}{"dIfr@ns}{$\underbrace{a - b}_{\text{Differenz}}$}
\eintrag{19}{der}{Minuend}{minuend}{"mInjuend}{$\underbrace{a}_{\text{Minuend}} - \underbrace{b}_{\text{Subtrahend}}$}
\eintrag{20}{der}{Subtrahend}{subtrahend}{"s2btr@hend}{$a - \frot{b}$}
\eintrag{21}{}{multiplizieren}{multiply}{"m2ltIplaI}{$a \cdot b$}
\eintrag{22}{das}{Produkt}{product}{"pr6d2kt}{$\underbrace{a \cdot b}_{\text{Produkt}}$}
\eintrag{23}{der}{Faktor}{factor}{"f\ae kt@}{$\underbrace{a}_{\text{Faktor}} \cdot \underbrace{b}_{\text{Faktor}}$}
\eintrag{24}{}{dividieren}{divide}{dI"vaId}{$a : b$}
\eintrag{25}{der}{Quotient}{quotient}{"kw@US\textsyllabic{n}t}{$\underbrace{a : b}_{\text{Quotient}}$}
\eintrag{26}{der}{Dividend}{dividend}{"dIvIdend}{$\underbrace{a}_{\text{Dividend}} : \underbrace{b}_{\text{Divisor}}$}
\eintrag{27}{der}{Divisor}{divisor}{dI"vaIz@}{$a : \frot{b}$}
\eintrag{28}{der}{Bruch}{fraction}{"fr\ae kS\textsyllabic{n}}{$\dfrac{a}{b}$}
\eintrag{29}{der}{Zähler}{numerator}{"nju:m@reIt@}{$\dfrac{\frot{a}}{b}$}
\eintrag{30}{der}{Nenner}{denominator}{dI"n6mIneIt@}{$\dfrac{a}{\frot{b}}$}
\eintrag{31}{der}{echte Bruch}{proper fraction}{"pr6p@ "fr\ae kS\textsyllabic{n}}{$\dfrac{3}{4}$}
\eintrag{32}{der}{unechte Bruch}{improper fraction}{Im"pr6p@ "fr\ae kS\textsyllabic{n}}{$\dfrac{7}{4}$}
\eintrag{33}{die}{gemischte Zahl}{mixed number}{"mIkst "n2mb@}{$1\tfrac{3}{4}$}
\eintrag{34}{}{gleichnamige Brüche}{like fractions}{"laIk "fr\ae kS\textsyllabic{n}z}{$\dfrac{1}{5};\ \dfrac{3}{5}$}
\eintrag{35}{}{ungleichnamige Brüche}{unlike fractions}{""2n"laIk "fr\ae kS\textsyllabic{n}z}{$\dfrac{1}{2};\ \dfrac{1}{3}$}
\eintrag{36}{der}{Hauptnenner}{lowest common denominator}{"l@UIst "k6m@n dI"n6mIneIt@}{$\dfrac{1}{2} + \dfrac{1}{3} = \dfrac{3}{\frot{6}} + \dfrac{2}{\frot{6}}$}
\eintrag{37}{}{kürzen}{cancel, reduce}{"k\ae ns\textsyllabic{l}, rI"dju:s}{$\dfrac{6}{8} = \dfrac{3}{4}$}
\eintrag{38}{}{erweitern}{expand (a fraction)}{Ik"sp\ae nd}{$\dfrac{3}{4} = \dfrac{15}{20}$}
\eintrag{39}{der}{Kehrwert}{reciprocal}{rI"sIpr@k\textsyllabic{l}}{$\dfrac{a}{b} \;\rightarrow\; \dfrac{b}{a}$}
\eintrag{40}{das}{Verhältnis}{ratio}{"reISi@U}{$a : b$}
\eintrag{41}{die}{Verhältnisgleichung}{proportion}{pr@"pO:S\textsyllabic{n}}{$a : b = c : d$}
\eintrag{42}{die}{Potenz}{power}{"paU@}{$a^n$}
\eintrag{43}{die}{Basis}{base}{beIs}{$\frot{a}^{n}$}
\eintrag{44}{der}{Exponent}{exponent}{Ik"sp@Un@nt}{$a^{\frot{n}}$}
\eintrag{45}{die}{Quadratwurzel}{square root}{""skwe@ "ru:t}{$\sqrt{a}$}
\eintrag{46}{die}{Kubikwurzel}{cube root}{"kju:b "ru:t}{$\sqrt[3]{a}$}
\eintrag{47}{die}{$n$-te Wurzel}{$n$th root}{"enT "ru:t}{$\sqrt[n]{a}$}
\eintrag{48}{der}{Radikand}{radicand}{"r\ae dIk\ae nd}{$\sqrt[n]{\frot{a}}$}
\eintrag{49}{der}{Wurzelexponent}{index (of a root)}{"Indeks}{$\sqrt[\frot{n}]{a}$}
\eintrag{50}{}{potenzieren}{raise to a power}{"reIz tu: @ "paU@}{$2^5 = 32$}
\eintrag{51}{}{quadrieren}{square}{skwe@}{$7^2 = 49$}
\eintrag{52}{}{radizieren, Wurzel ziehen}{extract a root}{Ik"str\ae kt @ "ru:t}{$\sqrt{49} = 7$}
\eintrag{53}{der}{Betrag}{absolute value}{"\ae bs@lu:t "v\ae lju:}{$|a|$}
\eintrag{54}{der}{Teiler}{divisor}{dI"vaIz@}{$3 \mid 12$}
\eintrag{55}{das}{Vielfache}{multiple}{"m2ltIp\textsyllabic{l}}{$3;\ 6;\ 9;\ 12;\ \ldots$}
\eintrag{56}{der}{größte gemeinsame Teiler}{greatest common divisor}{"greItIst "k6m@n dI"vaIz@}{$\operatorname{ggT}(12;18) = 6$ \newline \footnotesize engl. $\operatorname{gcd}$}
\eintrag{57}{das}{kleinste gemeinsame Vielfache}{least common multiple}{"li:st "k6m@n "m2ltIp\textsyllabic{l}}{$\operatorname{kgV}(4;6) = 12$ \newline \footnotesize engl. $\operatorname{lcm}$}
\eintrag{58}{der}{Prozentsatz}{percentage}{p@"sentIdZ}{$p\,\%$}
\eintrag{59}{}{runden}{round}{raUnd}{$3{,}14 \approx 3{,}1$}
\eintrag{60}{}{gleich}{equal}{"i:kw@l}{$=$}
\eintrag{61}{}{ungleich}{not equal}{n6t "i:kw@l}{$\neq$}
\eintrag{62}{}{identisch gleich}{identically equal}{aI"dentIkli "i:kw@l}{$\equiv$}
\eintrag{63}{}{kleiner als}{less than}{"les D@n}{$<$}
\eintrag{64}{}{größer als}{greater than}{"greIt@ D@n}{$>$}
\eintrag{65}{}{kleiner oder gleich}{less than or equal to}{"les D@n O:r "i:kw@l tu:}{$\le$}
\eintrag{66}{}{größer oder gleich}{greater than or equal to}{"greIt@ D@n O:r "i:kw@l tu:}{$\ge$}
\eintrag{67}{}{viel kleiner als}{much less than}{"m2tS "les D@n}{$\ll$}
\eintrag{68}{}{proportional zu}{proportional to}{pr@"pO:S@n\textsyllabic{l} tu:}{$y \sim x$ \newline \footnotesize engl. $y \propto x$}
\eintrag{69}{}{ungefähr}{approximately}{@"pr6ksIm@tli}{$\approx$}
\eintrag{70}{}{unendlich}{infinite}{"InfIn@t}{$\infty$}
\end{woerter}

\platz\subsection*{Beispiele}

\begin{lesen}
\lies{$12 + 30 = 42$}{The sum of twelve and thirty is forty-two.}{Die Summe von zwölf und dreißig ist zweiundvierzig.}
\lies{$3 \mid 12$}{Three is a divisor (a factor) of twelve; twelve is a multiple of three; twelve is divisible by three.}{Drei ist ein Teiler von zwölf; zwölf ist ein Vielfaches von drei; zwölf ist durch drei teilbar.}
\lies{$\operatorname{ggT}(12;\,18) = 6$}{The greatest common divisor of twelve and eighteen is six.}{Der größte gemeinsame Teiler von zwölf und achtzehn ist sechs.}
\lies{$17$}{Seventeen is a prime number: it has exactly two divisors, one and itself.}{Siebzehn ist eine Primzahl: Sie hat genau zwei Teiler, die Eins und sich selbst.}
\lies{$3{,}14159 \approx 3{,}14$}{Round the number to two decimal places.}{Runde die Zahl auf zwei Nachkommastellen.}
\lies{$|-7| = 7$}{The absolute value of minus seven is seven.}{Der Betrag von minus sieben ist sieben.}
\lies{$25\,\%$ von $80$ sind $20$}{Twenty-five per cent of eighty is twenty.}{Fünfundzwanzig Prozent von achtzig sind zwanzig.}
\end{lesen}


%%%%  ALGEBRA, GLEICHUNGEN UND MENGEN  %%%%
\newpage
\section{Algebra, Gleichungen und Mengen}

\begin{woerter}
\eintrag{71}{die}{Variable}{variable}{"ve@ri@b\textsyllabic{l}}{$x;\ y;\ a$}
\eintrag{72}{der}{Term}{expression}{Ik"spreS\textsyllabic{n}}{$3x + 2$}
\eintrag{73}{die}{Unbekannte}{unknown}{""2n"n@Un}{$x$}
\eintrag{74}{der}{Koeffizient}{coefficient}{""k@UI"fIS\textsyllabic{n}t}{$\frot{3}x^2$}
\eintrag{75}{der}{Grad}{degree}{dI"gri:}{$x^3 + x$: 3.\,Grad}
\eintrag{76}{die}{Gleichung}{equation}{I"kweIZ\textsyllabic{n}}{$2x + 3 = 7$}
\eintrag{77}{die}{Ungleichung}{inequality}{""InI"kw6l@ti}{$2x + 3 < 7$}
\eintrag{78}{die}{Lösung}{solution}{s@"lu:S\textsyllabic{n}}{$x = 2$}
\eintrag{79}{die}{Lösungsmenge}{solution set}{s@"lu:S\textsyllabic{n} set}{$L = \{2\}$}
\eintrag{80}{das}{Gleichungssystem}{system of equations}{"sIst@m @v I"kweIZ\textsyllabic{n}z}{$\begin{aligned} x + y &= 5\\ x - y &= 1\end{aligned}$}
\eintrag{81}{die}{lineare Gleichung}{linear equation}{"lIni@ I"kweIZ\textsyllabic{n}}{$3x - 7 = 11$}
\eintrag{82}{die}{quadratische Gleichung}{quadratic equation}{kw6"dr\ae tIk I"kweIZ\textsyllabic{n}}{$x^2 - 5x + 6 = 0$}
\eintrag{83}{die}{Wurzelgleichung}{radical equation}{"r\ae dIk\textsyllabic{l} I"kweIZ\textsyllabic{n}}{$\sqrt{x + 3} = 5$}
\eintrag{84}{die}{Exponentialgleichung}{exponential equation}{""eksp@"nenS\textsyllabic{l} I"kweIZ\textsyllabic{n}}{$2^x = 64$}
\eintrag{85}{die}{logarithmische Gleichung}{logarithmic equation}{""l6g@"rIDmIk I"kweIZ\textsyllabic{n}}{$\log_3 x = 4$}
\eintrag{86}{die}{trigonometrische Gleichung}{trigonometric equation}{""trIg@n@"metrIk I"kweIZ\textsyllabic{n}}{$2\sin x = 1$}
\eintrag{87}{die}{Klammer}{bracket}{"br\ae kIt}{$(\ \ )\quad [\ \ ]\quad \{\ \ \}$}
\eintrag{88}{die}{runden Klammern}{round brackets, parentheses}{"raUnd "br\ae kIts, p@"renT@si:z}{$(\ \ )$}
\eintrag{89}{die}{eckigen Klammern}{square brackets}{"skwe@ "br\ae kIts}{$[\ \ ]$}
\eintrag{90}{die}{geschweiften Klammern}{curly brackets, braces}{"k3:li "br\ae kIts, "breIsIz}{$\{\ \ \}$}
\eintrag{91}{}{ausmultiplizieren}{expand}{Ik"sp\ae nd}{$a(b + c) = ab + ac$}
\eintrag{92}{}{ausklammern}{factorise}{"f\ae kt@raIz}{$ab + ac = a(b + c)$}
\eintrag{93}{die}{binomische Formel}{binomial formula}{baI"n@Umi@l "fO:mj@l@}{$(a + b)^2 = a^2 + 2ab + b^2$}
\eintrag{94}{die}{Proportionalität}{proportionality}{pr@""pO:S@"n\ae l@ti}{$y = k \cdot x$}
\eintrag{95}{der}{Logarithmus}{logarithm}{"l6g@rID@m}{$\log_b a$}
\eintrag{96}{die}{Menge}{set}{set}{$M = \{1;\ 2;\ 3\}$}
\eintrag{97}{das}{Element}{element}{"elIm@nt}{$2 \in M$}
\eintrag{98}{die}{Teilmenge}{subset}{"s2bset}{$A \subseteq B$ \newline \skTeilmenge}
\eintrag{99}{die}{echte Teilmenge}{proper subset}{"pr6p@ "s2bset}{$A \subset B$}
\eintrag{100}{die}{Schnittmenge}{intersection}{""Int@"sekS\textsyllabic{n}}{$A \cap B$ \newline \skSchnittmenge}
\eintrag{101}{die}{Vereinigungsmenge}{union}{"ju:ni@n}{$A \cup B$ \newline \skVereinigung}
\eintrag{102}{die}{Differenzmenge}{set difference}{"set "dIfr@ns}{$A \setminus B$}
\eintrag{103}{das}{kartesische Produkt}{Cartesian product}{kA:"ti:zi@n "pr6d2kt}{$A \times B$}
\eintrag{104}{die}{leere Menge}{empty set}{"empti set}{$\emptyset = \{\,\}$}
\end{woerter}

\platz\subsection*{Beispiele}

\begin{lesen}
\lies{$3(x + 2) = 3x + 6$}{Expand the brackets.}{Multipliziere die Klammer aus.}
\lies{$6x + 9 = 3(2x + 3)$}{Factorise the expression: take out the common factor three.}{Klammere aus: Der gemeinsame Faktor ist drei.}
\lies{$(x + 3)^2 = x^2 + 6x + 9$}{$x$ plus three, all squared, equals $x$ squared plus six $x$ plus nine.}{$x$ plus drei in Klammern zum Quadrat ist gleich $x$ Quadrat plus sechs $x$ plus neun.}
\lies{$2x + 3 = 11$ \newline $x = 4$}{Solve the equation for $x$. The solution is $x$ equals four.}{Löse die Gleichung nach $x$ auf. Die Lösung ist $x$ gleich vier.}
\lies{$x + y = 5$ \newline $x - y = 1$}{Solve the system of equations: $x$ equals three and $y$ equals two.}{Löse das Gleichungssystem: $x$ gleich drei und $y$ gleich zwei.}
\lies{$A = \{1;\,2;\,3\},\ B = \{3;\,4\}$ \newline $A \cap B = \{3\}$}{The intersection of $A$ and $B$ contains only the element three.}{Die Schnittmenge von $A$ und $B$ enthält nur das Element drei.}
\end{lesen}


%%%%  FUNKTIONEN  %%%%
\newpage
\section{Funktionen}

\begin{woerter}
\eintrag{105}{die}{Funktion}{function}{"f2NkS\textsyllabic{n}}{$f(x)$}
\eintrag{106}{die}{unabhängige Variable}{independent variable}{""IndI"pend@nt "ve@ri@b\textsyllabic{l}}{$x$}
\eintrag{107}{die}{abhängige Variable}{dependent variable}{dI"pend@nt "ve@ri@b\textsyllabic{l}}{$y = f(x)$}
\eintrag{108}{der}{Funktionswert}{value of a function}{"v\ae lju: @v @ "f2NkS\textsyllabic{n}}{$f(3) = 15$}
\eintrag{109}{der}{Definitionsbereich}{domain}{d@"meIn}{$D_f$}
\eintrag{110}{der}{Wertebereich}{range}{reIndZ}{$W_f$}
\eintrag{111}{der}{Funktionsgraph}{graph}{grA:f}{\skGraph}
\eintrag{112}{die}{Wertetabelle}{table of values}{"teIb\textsyllabic{l} @v "v\ae lju:z}{\skWertetabelle}
\eintrag{113}{das}{Koordinatensystem}{coordinate system}{k@U"O:dIn@t "sIst@m}{\skKoordinaten}
\eintrag{114}{die}{$x$-Achse}{$x$-axis}{"eks ""\ae ksIs}{\skXAchse}
\eintrag{115}{der}{Ursprung}{origin}{"6rIdZIn}{$O(0 \mid 0)$}
\eintrag{116}{die}{Koordinaten}{coordinates}{k@U"O:dIn@ts}{$P(x_P \mid y_P)$}
\eintrag{117}{die}{Abszisse}{abscissa}{\ae b"sIs@}{$P(\frot{x_P} \mid y_P)$}
\eintrag{118}{die}{Ordinate}{ordinate}{"O:dIn@t}{$P(x_P \mid \frot{y_P})$}
\eintrag{119}{die}{lineare Funktion}{linear function}{"lIni@ "f2NkS\textsyllabic{n}}{$f(x) = mx + n$ \newline \skLinear}
\eintrag{120}{die}{Steigung}{gradient}{"greIdi@nt}{$m = \dfrac{\Delta y}{\Delta x}$ \newline \skSteigung}
\eintrag{121}{der}{$y$-Achsenabschnitt}{$y$-intercept}{"waI "Int@sept}{\skAchsenabschnitt}
\eintrag{122}{die}{Nullstelle}{zero, root}{"zI@r@U, ru:t}{$f(x_0) = 0$ \newline \skNullstelle}
\eintrag{123}{der}{Schnittpunkt}{point of intersection}{"pOInt @v ""Int@"sekS\textsyllabic{n}}{\skSchnittpunkt}
\eintrag{124}{die}{quadratische Funktion}{quadratic function}{kw6"dr\ae tIk "f2NkS\textsyllabic{n}}{$f(x) = ax^2 + bx + c$}
\eintrag{125}{die}{Parabel}{parabola}{p@"r\ae b@l@}{\skParabel}
\eintrag{126}{die}{ganzrationale Funktion}{polynomial function}{""p6li"n@Umi@l "f2NkS\textsyllabic{n}}{$f(x) = x^3 - 2x$}
\eintrag{127}{die}{gebrochenrationale Funktion}{rational function}{"r\ae Sn@l "f2NkS\textsyllabic{n}}{$f(x) = \dfrac{x + 1}{x - 2}$}
\eintrag{128}{die}{Wurzelfunktion}{square root function}{""skwe@ "ru:t "f2NkS\textsyllabic{n}}{$f(x) = \sqrt{x}$}
\eintrag{129}{der}{Scheitelpunkt}{vertex}{"v3:teks}{\skScheitelpunkt}
\eintrag{130}{die}{Exponentialfunktion}{exponential function}{""eksp@"nenS\textsyllabic{l} "f2NkS\textsyllabic{n}}{$f(x) = a \cdot b^x$ \newline \skExponential}
\eintrag{131}{die}{Logarithmusfunktion}{logarithmic function}{""l6g@"rIDmIk "f2NkS\textsyllabic{n}}{$f(x) = \log_b x$ \newline \skLogarithmus}
\eintrag{132}{der}{Sinus}{sine}{saIn}{$\sin x$ \newline \skSinus}
\eintrag{133}{der}{Kosinus}{cosine}{"k@UsaIn}{$\cos x$ \newline \skKosinus}
\eintrag{134}{der}{Tangens}{tangent}{"t\ae ndZ@nt}{$\tan x = \dfrac{\sin x}{\cos x}$}
\eintrag{135}{der}{Kotangens}{cotangent}{""k@U"t\ae ndZ@nt}{$\cot x = \dfrac{\cos x}{\sin x}$}
\eintrag{136}{die}{Arkusfunktion}{inverse trigonometric function}{""Inv3:s ""trIg@n@"metrIk "f2NkS\textsyllabic{n}}{$\arcsin x$ \newline \footnotesize engl. auch $\sin^{-1} x$}
\eintrag{137}{der}{Einheitskreis}{unit circle}{"ju:nIt ""s3:k\textsyllabic{l}}{\skEinheitskreis}
\eintrag{138}{das}{Gradmaß}{degree measure}{dI"gri: ""meZ@}{$360\degree$}
\eintrag{139}{das}{Bogenmaß}{radian measure}{"reIdi@n ""meZ@}{$180\degree \mathrel{\widehat{=}} \pi$ \newline \skBogenmass}
\eintrag{140}{der}{Sinussatz}{sine rule, law of sines}{"saIn ru:l}{$\frac{a}{\sin\alpha} = \frac{b}{\sin\beta} = \frac{c}{\sin\gamma}$}
\eintrag{141}{der}{Kosinussatz}{cosine rule, law of cosines}{"k@UsaIn ru:l}{\small$c^2 = a^2 + b^2 - 2ab\cos\gamma$}
\eintrag{142}{}{monoton steigend}{increasing}{In"kri:sIN}{$x_1 < x_2 \Rightarrow f(x_1) \le f(x_2)$ \newline \skMonotonSteigend}
\eintrag{143}{}{monoton fallend}{decreasing}{dI"kri:sIN}{$x_1 < x_2 \Rightarrow f(x_1) \ge f(x_2)$ \newline \skMonotonFallend}
\eintrag{144}{}{achsensymmetrisch zur $y$-Achse}{symmetric about the $y$-axis}{sI"metrIk @"baUt D@ "waI ""\ae ksIs}{$f(-x) = f(x)$ \newline \skAchsensymmetrisch}
\eintrag{145}{}{punktsymmetrisch zum Ursprung}{symmetric about the origin}{sI"metrIk @"baUt Di "6rIdZIn}{$f(-x) = -f(x)$ \newline \skPunktsymmetrisch}
\end{woerter}

\platz\subsection*{Beispiele}

\begin{lesen}
\lies{$f(x) = 2x - 4$}{$f$ is a linear function with the gradient two and the $y$-intercept minus four. Its zero is $x$ equals two.}{$f$ ist eine lineare Funktion mit der Steigung zwei und dem $y$-Achsenabschnitt minus vier. Ihre Nullstelle ist $x$ gleich zwei.}
\lies{$f(x) = (x - 1)^2 - 4$}{The graph is a parabola with the vertex $S(1 \mid -4)$. The zeros are minus one and three.}{Der Graph ist eine Parabel mit dem Scheitelpunkt $S(1 \mid -4)$. Die Nullstellen sind minus eins und drei.}
\lies{$f(x) = x^2$}{The graph is symmetric about the $y$-axis. The function is decreasing for $-\infty < x \le 0$ and increasing for $0 \le x < \infty$.}{Der Graph ist achsensymmetrisch zur $y$-Achse. Die Funktion ist fallend für $-\infty < x \le 0$ und steigend für $0 \le x < \infty$.}
\lies{$f(x) = 2^x$}{The graph of this exponential function passes through the point $(0 \mid 1)$; the $x$-axis is an asymptote.}{Der Graph dieser Exponentialfunktion verläuft durch den Punkt $(0 \mid 1)$; die $x$-Achse ist Asymptote.}
\lies{$D_f = \R$ \newline $W_f = \{y \in \R \mid y \ge -4\}$}{The domain is the set of all real numbers; the range is the set of all real numbers greater than or equal to minus four.}{Der Definitionsbereich ist die Menge aller reellen Zahlen; der Wertebereich ist die Menge aller reellen Zahlen, die größer oder gleich minus vier sind.}
\end{lesen}


%%%%  GEOMETRIE: GRUNDBEGRIFFE, WINKEL, ABBILDUNGEN  %%%%
\newpage
\section{Geometrie: Grundbegriffe, Winkel, Abbildungen}

\begin{woerter}
\eintrag{146}{der}{Punkt}{point}{pOInt}{\skPunkt}
\eintrag{147}{die}{Gerade}{straight line}{""streIt "laIn}{\skGerade}
\eintrag{148}{die}{Strecke}{line segment}{"laIn ""segm@nt}{$\overline{AB}$ \newline \skStrecke}
\eintrag{149}{der}{Strahl}{ray}{reI}{\skStrahl}
\eintrag{150}{}{parallel}{parallel}{"p\ae r@lel}{$g \parallel h$ \newline \skParallel}
\eintrag{151}{}{senkrecht}{perpendicular}{""p3:p@n"dIkj@l@}{$g \perp h$ \newline \skSenkrecht}
\eintrag{152}{der}{Abstand}{distance}{"dIst@ns}{\skAbstand}
\eintrag{153}{der}{Winkel}{angle}{"\ae Ng\textsyllabic{l}}{$\alpha;\ \beta;\ \gamma$ \newline \skWinkel{50}}
\eintrag{154}{der}{Scheitel}{vertex}{"v3:teks}{\skScheitel}
\eintrag{155}{der}{Schenkel}{arm, leg}{A:m, leg}{\skSchenkel}
\eintrag{156}{die}{Winkelhalbierende}{angle bisector}{"\ae Ng\textsyllabic{l} baI"sekt@}{\skWinkelhalbierende}
\eintrag{157}{der}{spitze Winkel}{acute angle}{@"kju:t "\ae Ng\textsyllabic{l}}{$0\degree < \alpha < 90\degree$ \newline \skWinkel{50}}
\eintrag{158}{der}{rechte Winkel}{right angle}{"raIt "\ae Ng\textsyllabic{l}}{$\alpha = 90\degree$ \newline \skRechterWinkel}
\eintrag{159}{der}{stumpfe Winkel}{obtuse angle}{@b"tju:s "\ae Ng\textsyllabic{l}}{$90\degree < \alpha < 180\degree$ \newline \skWinkel{130}}
\eintrag{160}{der}{gestreckte Winkel}{straight angle}{"streIt "\ae Ng\textsyllabic{l}}{$\alpha = 180\degree$ \newline \skWinkel{180}}
\eintrag{161}{der}{überstumpfe Winkel}{reflex angle}{"ri:fleks "\ae Ng\textsyllabic{l}}{$180\degree < \alpha < 360\degree$ \newline \skWinkel{250}}
\eintrag{162}{der}{Vollwinkel}{full angle}{"fUl "\ae Ng\textsyllabic{l}}{$\alpha = 360\degree$ \newline \skVollwinkel}
\eintrag{163}{der}{Nebenwinkel}{supplementary angle}{""s2plI"mentri "\ae Ng\textsyllabic{l}}{$\alpha + \beta = 180\degree$ \newline \skNebenwinkel}
\eintrag{164}{der}{Komplementwinkel}{complementary angle}{""k6mplI"mentri "\ae Ng\textsyllabic{l}}{$\alpha + \beta = 90\degree$ \newline \skKomplement}
\eintrag{165}{der}{Supplementwinkel}{supplementary angle}{""s2plI"mentri "\ae Ng\textsyllabic{l}}{$\alpha + \beta = 180\degree$}
\eintrag{166}{der}{Scheitelwinkel}{vertically opposite angle}{"v3:tIkli "6p@zIt "\ae Ng\textsyllabic{l}}{$\alpha = \alpha'$ \newline \skScheitelwinkel}
\eintrag{167}{der}{Stufenwinkel}{corresponding angle}{""k6r@"sp6ndIN "\ae Ng\textsyllabic{l}}{$\alpha = \beta$ \newline \skStufenwinkel}
\eintrag{168}{der}{Wechselwinkel}{alternate angle}{O:l"t3:n@t "\ae Ng\textsyllabic{l}}{$\alpha = \beta$ \newline \skWechselwinkel}
\eintrag{169}{}{entgegengesetzt liegende Winkel}{co-interior angles, allied angles}{""k@UIn"tI@ri@ "\ae Ng\textsyllabic{l}z, "\ae laId "\ae Ng\textsyllabic{l}z}{$\alpha + \beta = 180\degree$ \newline \skEntgegen}
\eintrag{170}{die}{Mittelsenkrechte}{perpendicular bisector}{""p3:p@n"dIkj@l@ baI"sekt@}{\skMittelsenkrechte}
\eintrag{171}{die}{Symmetrieachse}{axis of symmetry}{"\ae ksIs @v "sIm@tri}{\skSymmetrieachse}
\eintrag{172}{die}{Spiegelung}{reflection}{rI"flekS\textsyllabic{n}}{\skSpiegelung}
\eintrag{173}{die}{Verschiebung}{translation}{tr\ae ns"leIS\textsyllabic{n}}{\skVerschiebung}
\eintrag{174}{die}{Drehung}{rotation}{r@U"teIS\textsyllabic{n}}{\skDrehung}
\eintrag{175}{}{kongruent}{congruent}{"k6Ngru@nt}{$\triangle ABC \cong \triangle A'B'C'$}
\eintrag{176}{}{ähnlich}{similar}{"sIm@l@}{$\triangle ABC \sim \triangle A'B'C'$}
\end{woerter}

\platz\subsection*{Beispiele}

\begin{lesen}
\lies{$\overline{AB} = 5\,\mathrm{cm}$}{The line segment $A$ $B$ is five centimetres long.}{Die Strecke $A$ $B$ ist fünf Zentimeter lang.}
\lies{$g \parallel h$}{Draw a line $h$ through the point $P$ parallel to $g$.}{Zeichne durch den Punkt $P$ die Parallele $h$ zu $g$.}
\lies{$\alpha = 35\degree$}{Alpha is an acute angle of thirty-five degrees.}{Alpha ist ein spitzer Winkel von fünfunddreißig Grad.}
\lies{$\beta = 180\degree - 70\degree = 110\degree$}{Adjacent angles on a straight line add up to one hundred and eighty degrees.}{Nebenwinkel ergänzen sich zu hundertachtzig Grad.}
\lies{\skSpiegelung}{Reflect the triangle in the line $g$.}{Spiegele das Dreieck an der Geraden $g$.}
\lies{\skDrehung}{Rotate the triangle through ninety degrees about the point $Z$.}{Drehe das Dreieck um neunzig Grad um den Punkt $Z$.}
\lies{\skMittelsenkrechte}{Construct the perpendicular bisector of the line segment $A$ $B$.}{Konstruiere die Mittelsenkrechte der Strecke $A$ $B$.}
\end{lesen}


%%%%  EBENE FIGUREN  %%%%
\newpage
\section{Ebene Figuren}

\begin{woerter}
\eintrag{177}{das}{Dreieck}{triangle}{"traI\ae Ng\textsyllabic{l}}{$\triangle ABC$ \newline \skDreieck}
\eintrag{178}{der}{Innenwinkel}{interior angle}{In"tI@ri@ "\ae Ng\textsyllabic{l}}{$\alpha + \beta + \gamma = 180\degree$ \newline \skInnenwinkel}
\eintrag{179}{der}{Außenwinkel}{exterior angle}{Ik"stI@ri@ "\ae Ng\textsyllabic{l}}{\skAussenwinkel}
\eintrag{180}{das}{unregelmäßige Dreieck}{scalene triangle}{"skeIli:n "traI\ae Ng\textsyllabic{l}}{\skDreieck}
\eintrag{181}{das}{rechtwinklige Dreieck}{right-angled triangle}{"raIt ""\ae Ng\textsyllabic{l}d "traI\ae Ng\textsyllabic{l}}{\skRechtwinklig}
\eintrag{182}{das}{gleichschenklige Dreieck}{isosceles triangle}{aI"s6s@li:z "traI\ae Ng\textsyllabic{l}}{\skGleichschenklig}
\eintrag{183}{die}{Basis (im Dreieck)}{base}{beIs}{\skBasis}
\eintrag{184}{der}{Basiswinkel}{base angle}{"beIs ""\ae Ng\textsyllabic{l}}{\skBasiswinkel}
\eintrag{185}{das}{gleichseitige Dreieck}{equilateral triangle}{""i:kwI"l\ae t@r@l "traI\ae Ng\textsyllabic{l}}{\skGleichseitig}
\eintrag{186}{die}{Hypotenuse}{hypotenuse}{haI"p6t@nju:z}{\skHypotenuse}
\eintrag{187}{die}{Kathete}{leg}{leg}{\skKathete}
\eintrag{188}{die}{Ankathete}{adjacent side}{@"dZeIs\textsyllabic{n}t saId}{\skAnkathete}
\eintrag{189}{die}{Gegenkathete}{opposite side}{"6p@zIt saId}{\skGegenkathete}
\eintrag{190}{der}{Satz des Pythagoras}{Pythagorean theorem}{paI""T\ae g@"ri:@n "TI@r@m}{$a^2 + b^2 = c^2$ \newline \skPythagoras}
\eintrag{191}{die}{Höhe}{height, altitude}{haIt, "\ae ltItju:d}{\skHoehe}
\eintrag{192}{die}{Seitenhalbierende}{median}{"mi:di@n}{\skSeitenhalbierende}
\eintrag{193}{der}{Inkreis}{incircle, inscribed circle}{"Ins3:k\textsyllabic{l}, In"skraIbd "s3:k\textsyllabic{l}}{\skInkreis}
\eintrag{194}{der}{Umkreis}{circumcircle}{"s3:k@m""s3:k\textsyllabic{l}}{\skUmkreis}
\eintrag{195}{}{flächengleich}{equal in area}{"i:kw@l In "e@ri@}{$A_1 = A_2$}
\eintrag{196}{das}{Viereck}{quadrilateral}{""kw6drI"l\ae t@r@l}{\skViereck}
\eintrag{197}{die}{Diagonale}{diagonal}{daI"\ae g@n\textsyllabic{l}}{\skDiagonale}
\eintrag{198}{das}{Quadrat}{square}{skwe@}{\skQuadrat}
\eintrag{199}{das}{Rechteck}{rectangle}{"rekt\ae Ng\textsyllabic{l}}{\skRechteck}
\eintrag{200}{das}{Parallelogramm}{parallelogram}{""p\ae r@"lel@gr\ae m}{\skParallelogramm}
\eintrag{201}{die}{Raute}{rhombus}{"r6mb@s}{\skRaute}
\eintrag{202}{das}{Trapez}{trapezium}{tr@"pi:zi@m}{\skTrapez}
\eintrag{203}{das}{Drachenviereck}{kite}{kaIt}{\skDrachen}
\eintrag{204}{das}{Vieleck}{polygon}{"p6lIg@n}{\skFuenfeck}
\eintrag{205}{das}{regelmäßige Sechseck}{regular hexagon}{"regj@l@ "heks@g@n}{\skSechseck}
\eintrag{206}{das}{Fünfeck}{pentagon}{"pent@g@n}{\skNeck{5}}
\eintrag{207}{das}{Achteck}{octagon}{"6kt@g@n}{\skNeck{8}}
\eintrag{208}{das}{Zwölfeck}{dodecagon}{d@U"dek@g@n}{\skNeck{12}}
\eintrag{209}{der}{Umfang}{perimeter}{p@"rImIt@}{$u = 2a + 2b$ \newline \skUmfangViereck}
\eintrag{210}{der}{Flächeninhalt}{area}{"e@ri@}{$A = a \cdot b$ \newline \skFlaeche}
\eintrag{211}{der}{Kreis}{circle}{"s3:k\textsyllabic{l}}{\skKreis}
\eintrag{212}{der}{Mittelpunkt}{centre}{"sent@}{\skMittelpunkt}
\eintrag{213}{der}{Radius}{radius}{"reIdi@s}{\skRadius}
\eintrag{214}{der}{Durchmesser}{diameter}{daI"\ae mIt@}{$d = 2r$ \newline \skDurchmesser}
\eintrag{215}{der}{Kreisumfang}{circumference}{s@"k2mf@r@ns}{$u = 2\pi r$ \newline \skKreisumfang}
\eintrag{216}{die}{Kreiszahl}{pi}{paI}{$\pi \approx 3{,}14159$}
\eintrag{217}{der}{Kreissektor}{sector}{"sekt@}{\skSektor}
\eintrag{218}{die}{Sehne}{chord}{kO:d}{\skSehne}
\eintrag{219}{die}{Sekante}{secant}{"si:k@nt}{\skSekante}
\eintrag{220}{der}{Berührungspunkt}{point of tangency}{"pOInt @v "t\ae ndZ@nsi}{\skBeruehrpunkt}
\eintrag{221}{der}{Kreisbogen}{arc}{A:k}{\skKreisbogen}
\eintrag{222}{der}{Kreisabschnitt}{segment (of a circle)}{"segm@nt}{\skKreisabschnitt}
\eintrag{223}{die}{Kreisfläche}{area of a circle}{"e@ri@ @v @ "s3:k\textsyllabic{l}}{$A = \pi r^2$ \newline \skKreisflaeche}
\eintrag{224}{die}{Tangente}{tangent}{"t\ae ndZ@nt}{\skTangente}
\end{woerter}

\platz\subsection*{Beispiele}

\begin{lesen}
\lies{$a = 6\,\mathrm{cm},\ b = 4\,\mathrm{cm}$ \newline $A = 24\,\mathrm{cm}^2$}{The area of the rectangle is six centimetres times four centimetres, that is twenty-four square centimetres.}{Der Flächeninhalt des Rechtecks ist sechs Zentimeter mal vier Zentimeter, also vierundzwanzig Quadratzentimeter.}
\lies{$3^2 + 4^2 = 5^2$}{A triangle with the sides three, four and five is right-angled.}{Ein Dreieck mit den Seitenlängen drei, vier und fünf ist rechtwinklig.}
\lies{\skGleichschenklig}{An isosceles triangle has two equal sides and two equal base angles.}{Ein gleichschenkliges Dreieck hat zwei gleich lange Seiten und zwei gleich große Basiswinkel.}
\lies{\skDiagonale}{The diagonals of a rectangle are equal in length and bisect each other.}{Die Diagonalen eines Rechtecks sind gleich lang und halbieren einander.}
\lies{$d = 10\,\mathrm{cm}$ \newline $u = \pi \cdot d \approx 31{,}4\,\mathrm{cm}$}{The circumference of a circle with a diameter of ten centimetres is approximately thirty-one point four centimetres.}{Der Umfang eines Kreises mit zehn Zentimetern Durchmesser beträgt ungefähr einunddreißig Komma vier Zentimeter.}
\end{lesen}


%%%%  KÖRPER  %%%%
\newpage
\section{Körper}

\begin{woerter}
\eintrag{225}{der}{Körper}{solid}{"s6lId}{}
\eintrag{226}{der}{Würfel}{cube}{kju:b}{\skWuerfel}
\eintrag{227}{der}{Quader}{cuboid}{"kju:bOId}{\skQuader}
\eintrag{228}{das}{Prisma}{prism}{"prIz@m}{\skPrisma}
\eintrag{229}{der}{Zylinder}{cylinder}{"sIlInd@}{\skZylinder}
\eintrag{230}{die}{Pyramide}{pyramid}{"pIr@mId}{\skPyramide}
\eintrag{231}{der}{Kegel}{cone}{k@Un}{\skKegel}
\eintrag{232}{die}{Kugel}{sphere}{sfI@}{\skKugel}
\eintrag{233}{das}{Polyeder}{polyhedron}{""p6li"hi:dr@n}{}
\eintrag{234}{das}{Tetraeder}{tetrahedron}{""tetr@"hi:dr@n}{\skTetraeder}
\eintrag{235}{der}{Pyramidenstumpf}{frustum of a pyramid}{"fr2st@m @v @ "pIr@mId}{\skPyramidenstumpf}
\eintrag{236}{der}{Kegelstumpf}{frustum of a cone}{"fr2st@m @v @ "k@Un}{\skKegelstumpf}
\eintrag{237}{der}{Hohlzylinder}{hollow cylinder}{"h6l@U "sIlInd@}{\skHohlzylinder}
\eintrag{238}{der}{schiefe Zylinder}{oblique cylinder}{@"bli:k "sIlInd@}{\skSchieferZylinder}
\eintrag{239}{der}{Rotationskörper}{solid of revolution}{"s6lId @v ""rev@"lu:S\textsyllabic{n}}{\skRotation}
\eintrag{240}{das}{Kugelsegment}{spherical cap, spherical segment}{"sferIk\textsyllabic{l} k\ae p}{\skKugelsegment}
\eintrag{241}{der}{Kugelsektor}{spherical sector}{"sferIk\textsyllabic{l} "sekt@}{\skKugelsektor}
\eintrag{242}{das}{Rotationsellipsoid}{spheroid}{"sfI@rOId}{\skEllipsoid}
\eintrag{243}{das}{Rotationsparaboloid}{paraboloid of revolution}{p@"r\ae b@lOId}{\skParaboloid}
\eintrag{244}{das}{Rotationshyperboloid}{hyperboloid of revolution}{haI"p3:b@lOId}{\skHyperboloid}
\eintrag{245}{die}{Ecke}{vertex}{"v3:teks}{\skEcke}
\eintrag{246}{die}{Kante}{edge}{edZ}{\skKante}
\eintrag{247}{die}{Seitenfläche}{face}{feIs}{\skSeitenflaeche}
\eintrag{248}{die}{Grundfläche}{base}{beIs}{\skGrundflaeche}
\eintrag{249}{die}{Deckfläche}{top face}{"t6p feIs}{\skDeckflaeche}
\eintrag{250}{die}{Mantelfläche}{lateral surface, curved surface}{"l\ae t@r@l "s3:fIs, "k3:vd "s3:fIs}{$M = 2\pi r h$ \newline \skMantel}
\eintrag{251}{die}{Flächendiagonale}{face diagonal}{"feIs daI"\ae g@n\textsyllabic{l}}{\skFlaechendiagonale}
\eintrag{252}{die}{Raumdiagonale}{space diagonal}{"speIs daI"\ae g@n\textsyllabic{l}}{$d = \sqrt{a^2 + b^2 + c^2}$ \newline \skRaumdiagonale}
\eintrag{253}{das}{Netz}{net}{net}{\skNetz}
\eintrag{254}{das}{Schrägbild}{oblique drawing}{@"bli:k "drO:IN}{\skQuader}
\eintrag{255}{das}{Volumen}{volume}{"v6lju:m}{$V = a \cdot b \cdot c$}
\eintrag{256}{der}{Oberflächeninhalt}{surface area}{"s3:fIs ""e@ri@}{$O = 2(ab + ac + bc)$}
\end{woerter}

\platz\subsection*{Beispiele}

\begin{lesen}
\lies{$a = 3\,\mathrm{cm}$ \newline $V = a^3 = 27\,\mathrm{cm}^3$}{The volume of a cube with an edge length of three centimetres is twenty-seven cubic centimetres.}{Das Volumen eines Würfels mit drei Zentimetern Kantenlänge beträgt siebenundzwanzig Kubikzentimeter.}
\lies{$O = 6a^2 = 54\,\mathrm{cm}^2$}{Its surface area is fifty-four square centimetres.}{Sein Oberflächeninhalt beträgt vierundfünfzig Quadratzentimeter.}
\lies{\skNetz}{The net of a cube consists of six squares.}{Das Netz eines Würfels besteht aus sechs Quadraten.}
\lies{\skPyramide}{A square-based pyramid has five vertices, eight edges and five faces.}{Eine quadratische Pyramide hat fünf Ecken, acht Kanten und fünf Flächen.}
\lies{$r = 2\,\mathrm{cm},\ h = 5\,\mathrm{cm}$ \newline $V = \pi r^2 h \approx 62{,}8\,\mathrm{cm}^3$}{The volume of the cylinder is twenty pi, approximately sixty-two point eight cubic centimetres.}{Das Volumen des Zylinders ist zwanzig pi, ungefähr zweiundsechzig Komma acht Kubikzentimeter.}
\end{lesen}


%%%%  ANALYSIS  %%%%
\newpage
\section{Analysis}

\begin{woerter}
\eintrag{257}{die}{Folge}{sequence}{"si:kw@ns}{$(a_n)$ \newline \skFolge}
\eintrag{258}{die}{Nullfolge}{null sequence}{"n2l "si:kw@ns}{$\limn a_n = 0$}
\eintrag{259}{die}{arithmetische Folge}{arithmetic sequence}{""\ae rIT"metIk "si:kw@ns}{$a_{n+1} = a_n + d$}
\eintrag{260}{die}{geometrische Folge}{geometric sequence}{""dZi:@"metrIk "si:kw@ns}{$a_{n+1} = a_n \cdot q$}
\eintrag{261}{das}{Anfangsglied}{first term}{"f3:st t3:m}{$a_1$}
\eintrag{262}{}{konvergent}{convergent}{k@n"v3:dZ@nt}{}
\eintrag{263}{}{divergent}{divergent}{daI"v3:dZ@nt}{}
\eintrag{264}{der}{Grenzwert}{limit}{"lImIt}{$\limn a_n = g$}
\eintrag{265}{die}{Reihe}{series}{"sI@ri:z}{$\displaystyle\sum_{k=1}^{n} a_k$}
\eintrag{266}{die}{unendliche Reihe}{infinite series}{"InfIn@t "sI@ri:z}{$\displaystyle\sum_{k=1}^{\infty} a_k$}
\eintrag{267}{die}{Potenzreihe}{power series}{"paU@ ""sI@ri:z}{$\displaystyle\sum_{k=0}^{\infty} a_k x^k$}
\eintrag{268}{}{stetig}{continuous}{k@n"tInju@s}{$\limxto{x_0} f(x) = f(x_0)$}
\eintrag{269}{}{unstetig}{discontinuous}{""dIsk@n"tInju@s}{}
\eintrag{270}{das}{Intervall}{interval}{"Int@v\textsyllabic{l}}{$[a;\,b]$}
\eintrag{271}{die}{Asymptote}{asymptote}{"\ae sImpt@Ut}{\skAsymptote}
\eintrag{272}{der}{Differenzenquotient}{difference quotient}{"dIfr@ns "kw@US\textsyllabic{n}t}{$\dfrac{f(x_0 + h) - f(x_0)}{h}$}
\eintrag{273}{die}{Ableitung}{derivative}{dI"rIv@tIv}{$f'(x)$}
\eintrag{274}{die}{Differenzialrechnung}{differential calculus}{""dIf@"renS\textsyllabic{l} "k\ae lkj@l@s}{}
\eintrag{275}{der}{Differenzialquotient}{differential quotient}{""dIf@"renS\textsyllabic{l} "kw@US\textsyllabic{n}t}{$\dfrac{\mathrm{d}y}{\mathrm{d}x}$}
\eintrag{276}{das}{Differenzial}{differential}{""dIf@"renS\textsyllabic{l}}{$\mathrm{d}y$}
\eintrag{277}{die}{partielle Ableitung}{partial derivative}{"pA:S\textsyllabic{l} dI"rIv@tIv}{$\dfrac{\partial z}{\partial x}$}
\eintrag{278}{}{ableiten}{differentiate}{""dIf@"renSieIt}{$\dfrac{\mathrm{d}f}{\mathrm{d}x}$}
\eintrag{279}{die}{Tangente}{tangent}{"t\ae ndZ@nt}{$m_t = f'(x_0)$ \newline \skTangenteGraph}
\eintrag{280}{die}{Funktionsuntersuchung}{curve sketching}{"k3:v ""sketSIN}{}
\eintrag{281}{der}{Extrempunkt}{extreme point}{Ik"stri:m ""pOInt}{$f'(x_E) = 0$}
\eintrag{282}{der}{Hochpunkt}{maximum}{"m\ae ksIm@m}{$f''(x_E) < 0$ \newline \skHochpunkt}
\eintrag{283}{der}{Tiefpunkt}{minimum}{"mInIm@m}{$f''(x_E) > 0$ \newline \skTiefpunkt}
\eintrag{284}{der}{Wendepunkt}{point of inflection}{"pOInt @v In"flekS\textsyllabic{n}}{$f''(x_W) = 0$ \newline \skWendepunkt}
\eintrag{285}{die}{Stammfunktion}{antiderivative}{""\ae ntidI"rIv@tIv}{$F'(x) = f(x)$}
\eintrag{286}{die}{Integralrechnung}{integral calculus}{"IntIgr@l "k\ae lkj@l@s}{}
\eintrag{287}{das}{unbestimmte Integral}{indefinite integral}{In"defIn@t "IntIgr@l}{$\displaystyle\int f(x)\dx$}
\eintrag{288}{das}{bestimmte Integral}{definite integral}{"defIn@t "IntIgr@l}{$\displaystyle\int_a^b f(x)\dx$}
\eintrag{289}{das}{Integralzeichen}{integral sign}{"IntIgr@l saIn}{$\displaystyle\int$}
\eintrag{290}{der}{Integrand}{integrand}{"IntIgr\ae nd}{$\displaystyle\int \frot{f(x)}\dx$}
\eintrag{291}{die}{Integrationsvariable}{variable of integration}{"ve@ri@b\textsyllabic{l} @v ""IntI"greIS\textsyllabic{n}}{$\displaystyle\int f(\frot{x})\,\mathrm{d}\frot{x}$}
\eintrag{292}{die}{Integrationskonstante}{constant of integration}{"k6nst@nt @v ""IntI"greIS\textsyllabic{n}}{$F(x) + \frot{C}$}
\eintrag{293}{die}{Obersumme}{upper sum}{"2p@ s2m}{\skObersumme}
\eintrag{294}{die}{Untersumme}{lower sum}{"l@U@ s2m}{\skUntersumme}
\eintrag{295}{das}{Integral}{integral}{"IntIgr@l}{$\displaystyle\int_a^b f(x)\dx$ \newline \skIntegral}
\eintrag{296}{}{integrieren}{integrate}{"IntIgreIt}{$\displaystyle\int f(x)\dx = F(x) + C$}
\eintrag{297}{die}{Integrationsgrenze}{limit of integration}{"lImIt @v ""IntI"greIS\textsyllabic{n}}{$\displaystyle\int_{\frot{a}}^{\frot{b}} f(x)\dx$}
\eintrag{298}{die}{eulersche Zahl}{Euler's number}{"OIl@z "n2mb@}{$e \approx 2{,}71828$}
\eintrag{299}{die}{natürliche Exponentialfunktion}{natural exponential function}{"n\ae tSr@l ""eksp@"nenS\textsyllabic{l} "f2NkS\textsyllabic{n}}{$f(x) = e^x$}
\eintrag{300}{der}{natürliche Logarithmus}{natural logarithm}{"n\ae tSr@l "l6g@rID@m}{$\ln x$}
\end{woerter}

\platz\subsection*{Beispiele}

\begin{lesen}
\lies{$f(x) = x^3 - 3x$ \newline $f'(x) = 3x^2 - 3$}{Differentiate $f$: the derivative is three $x$ squared minus three.}{Leite $f$ ab: Die Ableitung ist drei $x$ Quadrat minus drei.}
\lies{$f'(x) = 0 \ \Rightarrow\ x = \pm 1$}{The derivative is zero at $x$ equals plus or minus one.}{Die Ableitung ist null für $x$ gleich plus oder minus eins.}
\lies{$H(-1 \mid 2),\ T(1 \mid -2)$ \newline $W(0 \mid 0)$}{The graph has a maximum at minus one two, a minimum at one minus two and a point of inflection at the origin.}{Der Graph hat den Hochpunkt $H$, den Tiefpunkt $T$ und den Wendepunkt $W$ im Ursprung.}
\lies{$F(x) = x^2$ \newline $f(x) = 2x$}{Capital $F$ is an antiderivative of $f$.}{Groß $F$ ist eine Stammfunktion von $f$.}
\lies{$\displaystyle\int_0^3 2x \dx = \big[x^2\big]_0^3 = 9$}{The integral from nought to three of two $x$ d $x$ equals nine.}{Das Integral von null bis drei über zwei $x$ d $x$ ist neun.}
\lies{$\limn \dfrac{2n + 1}{n} = 2$}{The sequence converges; its limit is two.}{Die Folge konvergiert; ihr Grenzwert ist zwei.}
\end{lesen}


%%%%  VEKTOREN UND ANALYTISCHE GEOMETRIE  %%%%
\newpage
\section{Vektoren und analytische Geometrie}

\begin{woerter}
\eintrag{301}{der}{Vektor}{vector}{"vekt@}{$\vec{a} = \vektor{a_1}{a_2}{a_3}$ \newline \skVektor}
\eintrag{302}{der}{Ortsvektor}{position vector}{p@"zIS\textsyllabic{n} "vekt@}{$\overrightarrow{OP}$ \newline \skOrtsvektor}
\eintrag{303}{der}{Einheitsvektor}{unit vector}{"ju:nIt ""vekt@}{$|\vec{e}\,| = 1$}
\eintrag{304}{der}{Nullvektor}{zero vector}{"zI@r@U ""vekt@}{$\vec{0}$}
\eintrag{305}{der}{Betrag eines Vektors}{magnitude}{"m\ae gnItju:d}{$|\vec{a}| = \sqrt{a_1^2 + a_2^2 + a_3^2}$}
\eintrag{306}{das}{Skalarprodukt}{dot product, scalar product}{"d6t ""pr6d2kt, "skeIl@ ""pr6d2kt}{$\vec{a} \cdot \vec{b}$}
\eintrag{307}{das}{Kreuzprodukt}{cross product}{"kr6s ""pr6d2kt}{$\vec{a} \times \vec{b}$ \newline \skKreuzprodukt}
\eintrag{308}{die}{Ebene}{plane}{pleIn}{\skEbene}
\eintrag{309}{die}{Geradengleichung}{equation of a straight line}{I"kweIZ\textsyllabic{n} @v @ ""streIt "laIn}{$y = mx + n$}
\eintrag{310}{der}{Steigungswinkel}{angle of inclination}{"\ae Ng\textsyllabic{l} @v ""InklI"neIS\textsyllabic{n}}{$m = \tan\alpha$}
\eintrag{311}{die}{Polarkoordinaten}{polar coordinates}{"p@Ul@ k@U"O:dIn@ts}{$P(r;\,\varphi)$ \newline \skPolar}
\eintrag{312}{der}{Kegelschnitt}{conic section}{"k6nIk "sekS\textsyllabic{n}}{}
\eintrag{313}{die}{Kreisgleichung}{equation of a circle}{I"kweIZ\textsyllabic{n} @v @ "s3:k\textsyllabic{l}}{$x^2 + y^2 = r^2$ \newline \skKreisgleichung}
\eintrag{314}{die}{Ellipse}{ellipse}{I"lIps}{$\frac{x^2}{a^2} + \frac{y^2}{b^2} = 1$ \newline \skEllipse}
\eintrag{315}{die}{Hyperbel}{hyperbola}{haI"p3:b@l@}{$\frac{x^2}{a^2} - \frac{y^2}{b^2} = 1$ \newline \skHyperbel}
\eintrag{316}{der}{Brennpunkt}{focus}{"f@Uk@s}{\skBrennpunkt}
\eintrag{317}{die}{Halbachse}{semi-axis}{""semi"\ae ksIs}{\skHalbachse}
\eintrag{318}{die}{Leitlinie}{directrix}{d@"rektrIks}{\skLeitlinie}
\eintrag{319}{}{linear abhängig}{linearly dependent}{"lIni@li dI"pend@nt}{$\vec{b} = r \cdot \vec{a}$}
\eintrag{320}{die}{Matrix}{matrix}{"meItrIks}{$\begin{pmatrix} a & b\\ c & d\end{pmatrix}$}
\eintrag{321}{die}{quadratische Matrix}{square matrix}{"skwe@ "meItrIks}{$n \times n$}
\eintrag{322}{die}{Zeile}{row}{r@U}{}
\eintrag{323}{die}{Spalte}{column}{"k6l@m}{}
\eintrag{324}{die}{Hauptdiagonale}{main diagonal}{"meIn daI"\ae g@n\textsyllabic{l}}{$\begin{pmatrix} \frot{a} & b\\ c & \frot{d}\end{pmatrix}$}
\eintrag{325}{die}{Determinante}{determinant}{dI"t3:mIn@nt}{$\begin{vmatrix} a & b\\ c & d\end{vmatrix} = ad - bc$}
\end{woerter}

\platz\subsection*{Beispiele}

\begin{lesen}
\lies{$\vec{a} = \vektor{3}{4}{0}$ \newline $|\vec{a}\,| = \sqrt{9 + 16 + 0} = 5$}{The magnitude of vector $a$ is five.}{Der Betrag des Vektors $a$ ist fünf.}
\lies{$A(1 \mid 2 \mid 3),\ B(4 \mid 6 \mid 3)$ \newline $\overrightarrow{AB} = \vektor{3}{4}{0}$}{The vector from $A$ to $B$ is three, four, nought.}{Der Vektor von $A$ nach $B$ ist drei, vier, null.}
\lies{$g\colon\ \vec{x} = \vektor{1}{2}{3} + t \cdot \vektor{3}{4}{0}$}{The equation of the line $g$: a position vector plus $t$ times a direction vector.}{Die Gleichung der Geraden $g$: Stützvektor plus $t$ mal Richtungsvektor.}
\lies{$\vec{a} \cdot \vec{b} = 0$}{The vectors $a$ and $b$ are perpendicular (orthogonal).}{Die Vektoren $a$ und $b$ sind zueinander senkrecht (orthogonal).}
\lies{$\vec{b} = 2 \cdot \vec{a}$}{The vectors $a$ and $b$ are linearly dependent (parallel).}{Die Vektoren $a$ und $b$ sind linear abhängig (parallel).}
\end{lesen}


%%%%  STOCHASTIK  %%%%
\newpage
\section{Stochastik}

\begin{woerter}
\eintrag{326}{das}{Zufallsexperiment}{random experiment}{"r\ae nd@m Ik"sperIm@nt}{}
\eintrag{327}{das}{Ergebnis}{outcome}{"aUtk2m}{$\omega \in \Omega$}
\eintrag{328}{das}{Ereignis}{event}{I"vent}{$A \subseteq \Omega$}
\eintrag{329}{das}{Gegenereignis}{complementary event}{""k6mplI"mentri I"vent}{$\overline{A}$}
\eintrag{330}{die}{Wahrscheinlichkeit}{probability}{""pr6b@"bIl@ti}{$P(A)$}
\eintrag{331}{die}{relative Häufigkeit}{relative frequency}{"rel@tIv "fri:kw@nsi}{$h_n(A)$}
\eintrag{332}{das}{Baumdiagramm}{tree diagram}{"tri: ""daI@gr\ae m}{\skBaum}
\eintrag{333}{die}{Vierfeldertafel}{two-way table}{"tu: weI "teIb\textsyllabic{l}}{\skVierfeldertafel}
\eintrag{334}{die}{Zufallsgröße}{random variable}{"r\ae nd@m "ve@ri@b\textsyllabic{l}}{$X$}
\eintrag{335}{der}{Erwartungswert}{expected value}{Ik"spektId "v\ae lju:}{$\E(X) = \mu$}
\eintrag{336}{die}{Varianz}{variance}{"ve@ri@ns}{$\V(X)$}
\eintrag{337}{die}{Standardabweichung}{standard deviation}{"st\ae nd@d ""di:vi"eIS\textsyllabic{n}}{$\sigma = \sqrt{\V(X)}$}
\eintrag{338}{der}{Mittelwert}{mean}{mi:n}{$\bar{x}$}
\eintrag{339}{der}{Median}{median}{"mi:di@n}{$\tilde{x}$}
\eintrag{340}{die}{Fakultät}{factorial}{f\ae k"tO:ri@l}{$n! = 1 \cdot 2 \cdot \ldots \cdot n$}
\eintrag{341}{die}{Permutation}{permutation}{""p3:mju"teIS\textsyllabic{n}}{$n!$}
\eintrag{342}{die}{Variation}{variation, arrangement}{""ve@ri"eIS\textsyllabic{n}, @"reIndZm@nt}{$\dfrac{n!}{(n-k)!}$}
\eintrag{343}{die}{Kombination}{combination}{""k6mbI"neIS\textsyllabic{n}}{$\dbinom{n}{k}$}
\eintrag{344}{der}{Binomialkoeffizient}{binomial coefficient}{baI"n@Umi@l ""k@UI"fIS\textsyllabic{n}t}{$\dbinom{n}{k}$}
\eintrag{345}{die}{Binomialverteilung}{binomial distribution}{baI"n@Umi@l ""dIstrI"bju:S\textsyllabic{n}}{$B_{n;p}(k)$ \newline \skSaeulen}
\end{woerter}

\platz\subsection*{Beispiele}

\begin{lesen}
\lies{$P(\text{„6“}) = \dfrac16$}{The probability of rolling a six with a fair die is one sixth.}{Die Wahrscheinlichkeit, mit einem fairen Würfel eine Sechs zu würfeln, beträgt ein Sechstel.}
\lies{$P(\text{„zweimal Kopf“}) = \dfrac12 \cdot \dfrac12 = \dfrac14$}{A fair coin is tossed twice. The probability of getting two heads is one quarter.}{Eine faire Münze wird zweimal geworfen. Die Wahrscheinlichkeit für zweimal Kopf beträgt ein Viertel.}
\lies{$P(\overline{A}) = 1 - P(A)$}{The probability of the complementary event equals one minus the probability of $A$.}{Die Wahrscheinlichkeit des Gegenereignisses ist eins minus die Wahrscheinlichkeit von $A$.}
\lies{$2;\ 4;\ 6;\ 8$ \newline $\bar{x} = 5$}{The mean of the four values is five.}{Der Mittelwert der vier Werte ist fünf.}
\lies{$\E(X) = 3{,}5$}{The expected value of the number shown by a fair die is three point five.}{Der Erwartungswert der Augenzahl eines fairen Würfels ist drei Komma fünf.}
\end{lesen}


%%%%  LOGIK UND BEWEIS  %%%%
\newpage
\section{Logik und Beweis}

\begin{woerter}
\eintrag{346}{die}{Definition}{definition}{""defI"nIS\textsyllabic{n}}{}
\eintrag{347}{der}{Satz}{theorem}{"TI@r@m}{}
\eintrag{348}{die}{Voraussetzung}{assumption, hypothesis}{@"s2mpS\textsyllabic{n}, haI"p6T@sIs}{}
\eintrag{349}{der}{Beweis}{proof}{pru:f}{$\square$}
\eintrag{350}{}{daraus folgt}{implies}{Im"plaIz}{$A \Rightarrow B$}
\eintrag{351}{}{genau dann, wenn}{if and only if}{If @nd "@Unli If}{$A \Leftrightarrow B$}
\eintrag{352}{}{für alle}{for all}{f@r "O:l}{$\forall$}
\eintrag{353}{}{es gibt}{there exists}{De@r Ig"zIsts}{$\exists$}
\end{woerter}

\platz\subsection*{Beispiele}

\begin{lesen}
\lies{$x > 2 \ \Rightarrow\ x^2 > 4$}{If $x$ is greater than two, then $x$ squared is greater than four.}{Wenn $x$ größer als zwei ist, dann ist $x$ Quadrat größer als vier.}
\lies{$n$ gerade $\Leftrightarrow$ $n^2$ gerade}{$n$ is even if and only if $n$ squared is even.}{$n$ ist genau dann gerade, wenn $n$ Quadrat gerade ist.}
\lies{$\forall\, x \in \R\colon\ x^2 \ge 0$}{For all real numbers $x$, $x$ squared is greater than or equal to zero.}{Für alle reellen Zahlen $x$ gilt: $x$ Quadrat ist größer oder gleich null.}
\lies{$\exists\, x \in \R\colon\ x^2 = 2$}{There exists a real number $x$ such that $x$ squared equals two.}{Es gibt eine reelle Zahl $x$, für die $x$ Quadrat gleich zwei ist.}
\lies{Annahme}{Assume that the square root of two is rational. \ldots\ This is a contradiction. (proof by contradiction)}{Angenommen, die Wurzel aus zwei ist rational. \ldots\ Das ist ein Widerspruch. (Beweis durch Widerspruch)}
\end{lesen}



%%%%  TEIL B: MATHEMATIK LESEN UND SPRECHEN  %%%%
\newpage
\part{Mathematik lesen und sprechen}

In diesem Teil steht, wie man Zeichen, Formeln und Rechenwege auf Englisch (und zum Vergleich auf
Deutsch) vorliest. Links steht der Ausdruck, rechts die englische Lesart und darunter grau die
deutsche. Wo es mehrere übliche Lesarten gibt, sind sie durch Schrägstriche getrennt. Die Beispiele
sind bewusst einfach gehalten, damit sie sich direkt im Unterricht verwenden lassen.

%--------------------------------------------------------------------------------------------------
\section{Zahlen, Zeichen und Schreibweisen}

\platz\subsection*{Ordnungszahlen, Indizes und Aufzählungen}

\begin{lesen}
\lies{1.\quad 2.\quad 3.\quad 21.}{1st, 2nd, 3rd, 21st -- first, second, third, twenty-first}{erstens, zweitens, drittens -- der Erste, der Zweite, \ldots}
\lies{1.\ \ldots, 2.\ \ldots}{firstly \ldots, secondly \ldots}{erstens \ldots, zweitens \ldots}
\lies{$x_1,\ x_2$}{$x$ one, $x$ two \;/\; $x$ sub one, $x$ sub two}{$x$ eins, $x$ zwei \;/\; $x$ Index eins}
\lies{$a_n$}{$a$ sub $n$ \;/\; $a$ $n$}{$a$ Index $n$ \;/\; $a$ $n$}
\lies{$1;\ 2;\ 3;\ \ldots$}{one, two, three, and so on (to infinity)}{eins, zwei, drei und so weiter (unbegrenzt)}
\lies{$1;\ 2;\ \ldots;\ 25$}{one, two, and so on up to twenty-five}{eins, zwei und so weiter bis fünfundzwanzig}
\end{lesen}

\platz\subsection*{Dezimalzahlen, große Zahlen und Brüche}

Im Englischen trennt ein \textbf{Punkt} die Nachkommastellen ab, das Komma gliedert große Zahlen in
Dreiergruppen -- im Deutschen ist es genau umgekehrt. Das führt leicht zu Verwechslungen.

\begin{lesen}
\lies{$3{,}14$ \newline \footnotesize engl. 3.14}{three point one four}{drei Komma eins vier}
\lies{$0{,}3$ \newline \footnotesize engl. 0.3 oder .3}{nought point three \ipa{nO:t} \;/\; zero point three \;/\; point three}{null Komma drei}
\lies{$30\,000$ \newline \footnotesize engl. 30,000}{thirty thousand}{dreißigtausend}
\lies{$2{,}5 \cdot 10^6$ \newline \footnotesize engl. $2.5 \times 10^6$}{two point five times ten to the (power of) six}{zwei Komma fünf mal zehn hoch sechs}
\lies{$\frac12;\ \frac13;\ \frac14$}{a half, a third, a quarter}{ein Halb, ein Drittel, ein Viertel}
\lies{$\frac34;\ \frac58$}{three quarters (AE: three fourths), five eighths}{drei Viertel, fünf Achtel}
\lies{$\frac{1}{100}$}{one hundredth}{ein Hundertstel}
\lies{$\frac{17}{25}$}{seventeen over twenty-five \;/\; seventeen twenty-fifths}{siebzehn Fünfundzwanzigstel}
\end{lesen}

\platz\subsection*{Britisches und amerikanisches Englisch}

Die Lautschrift im Wörterbuch folgt dem britischen Englisch. Im amerikanischen Englisch heißen
einige Dinge anders:

\begin{lesen}
\lies{\textbf{Mathematik}}{maths (BE) \;/\; math (AE)}{}
\lies{$0$}{nought, zero (BE) \;/\; zero (AE)}{}
\lies{$(\ \ )\quad[\ \ ]$}{brackets, square brackets (BE) \;/\; parentheses, brackets (AE)}{}
\lies{Trapez}{trapezium (BE) \;/\; trapezoid (AE)}{}
\lies{$\dfrac34$}{three quarters (BE) \;/\; three fourths (AE)}{}
\lies{Schreibung}{factorise, centre (BE) \;/\; factorize, center (AE)}{}
\end{lesen}

\anmerkung{In englischsprachigen Texten steht zwischen zwei Zahlen meist das Kreuz: $3 \times 5$. Ein
mittiger Punkt kann in älteren britischen Texten als Dezimalpunkt gemeint sein ($3{\cdot}5 = 3{,}5$).
Zwischen Variablen ist der Malpunkt üblich oder wird ganz weggelassen: $a \cdot b = ab$.}

%--------------------------------------------------------------------------------------------------

\platz\subsection*{Beispiele}

\begin{lesen}
\lies{$7{,}25$}{seven point two five}{sieben Komma zwei fünf}
\lies{$0{,}05$}{nought point nought five}{null Komma null fünf}
\lies{$1\,250\,000$ \newline \footnotesize engl. 1,250,000}{one million two hundred and fifty thousand}{eine Million zweihundertfünfzigtausend}
\end{lesen}

\newpage
\section{Relationszeichen}

\begin{lesen}
\lies{$p = q$}{$p$ equals $q$ \;/\; $p$ is equal to $q$}{$p$ (ist) gleich $q$}
\lies{$a \neq b$}{$a$ is not equal to $b$ \;/\; $a$ does not equal $b$}{$a$ (ist) ungleich $b$ \;/\; $a$ nicht gleich $b$}
\lies{$\pi \approx 3{,}14$}{pi is approximately equal to three point one four}{pi (ist) ungefähr gleich drei Komma eins vier}
\lies{$f(x) \equiv 0$}{$f$ of $x$ is identically equal to zero}{$f$ von $x$ ist identisch gleich null}
\lies{$a < b$}{$a$ is less than $b$}{$a$ (ist) kleiner als $b$}
\lies{$a > b$}{$a$ is greater than $b$}{$a$ (ist) größer als $b$}
\lies{$a \le b$}{$a$ is less than or equal to $b$ \;/\; $a$ is not greater than $b$}{$a$ (ist) kleiner oder gleich $b$ \;/\; höchstens gleich $b$}
\lies{$a \ge b$}{$a$ is greater than or equal to $b$ \;/\; $a$ is not less than $b$}{$a$ (ist) größer oder gleich $b$ \;/\; mindestens gleich $b$}
\lies{$a \ll b$}{$a$ is much less than $b$}{$a$ (ist) viel kleiner als $b$}
\lies{$a \gg b$}{$a$ is much greater than $b$}{$a$ (ist) viel größer als $b$}
\lies{$y \sim x$ \newline \footnotesize engl. $y \propto x$}{$y$ is (directly) proportional to $x$}{$y$ ist (direkt) proportional zu $x$}
\lies{$a : b = c : d$}{$a$ is to $b$ as $c$ is to $d$}{$a$ verhält sich zu $b$ wie $c$ zu $d$}
\lies{$2 \le x \le 5$}{$x$ is greater than or equal to two and less than or equal to five \;/\; $x$ lies between two and five inclusive}{$x$ ist mindestens zwei und höchstens fünf \;/\; $x$ liegt zwischen zwei und fünf (einschließlich)}
\end{lesen}

\anmerkung{Das Zeichen $\sim$ bedeutet im Deutschen „proportional“, in der Geometrie aber „ähnlich“
($\triangle ABC \sim \triangle DEF$, \textit{similar}). Im Englischen schreibt man für
„proportional“ meist $\propto$.}

%--------------------------------------------------------------------------------------------------

\platz\subsection*{Beispiele}

\begin{lesen}
\lies{$5 < 8$}{five is less than eight}{fünf ist kleiner als acht}
\lies{$x \ge 0$}{$x$ is greater than or equal to zero: $x$ is non-negative}{$x$ ist größer oder gleich null: $x$ ist nicht negativ}
\lies{$7 \neq 9$}{seven is not equal to nine}{sieben ist ungleich neun}
\lies{$\sqrt{2} \approx 1{,}41$}{the square root of two is approximately one point four one}{die Wurzel aus zwei ist ungefähr eins Komma vier eins}
\lies{$-2 < x < 2$}{$x$ lies between minus two and two}{$x$ liegt zwischen minus zwei und zwei}
\end{lesen}

\newpage
\section{Grundrechenarten}

Die Namen der Rechengrößen auf einen Blick:
\[
\ub{a}{summand}{Summand} + \ub{b}{summand}{Summand} = \ub{s}{sum}{Summe}
\qquad\qquad
\ub{a}{minuend}{Minuend} - \ub{b}{subtrahend}{Subtrahend} = \ub{d}{difference}{Differenz}
\]
\[
\ub{a}{factor}{Faktor} \cdot \ub{b}{factor}{Faktor} = \ub{p}{product}{Produkt}
\qquad\qquad
\ub{a}{dividend}{Dividend} : \ub{b}{divisor}{Divisor} = \ub{q}{quotient}{Quotient}
\]

\begin{lesen}
\lies{$a_1 + a_2 = s$}{$a$ one plus $a$ two equals $s$}{$a$ eins plus $a$ zwei gleich $s$}
\lies{$L - l = d$}{capital $L$ minus small $l$ equals $d$}{groß $L$ minus klein $l$ gleich $d$}
\lies{$a \cdot b = c$ \newline \footnotesize engl. auch $a \times b$}{$a$ times $b$ equals $c$ \;/\; $a$ multiplied by $b$ equals $c$}{$a$ mal $b$ gleich $c$}
\lies{$ab = c$}{$a$ $b$ equals $c$}{$a$ $b$ gleich $c$}
\lies{$a : b = c$ \newline \footnotesize engl. auch $a \div b$}{$a$ divided by $b$ equals $c$}{$a$ geteilt durch $b$ gleich $c$ \;/\; $a$ durch $b$ gleich $c$}
\lies{$\dfrac{a}{b} = a/b$}{$a$ over $b$}{$a$ durch $b$ \;/\; $a$ $b$-tel}
\end{lesen}

\platz\subsection*{Beispiele}

\begin{lesen}
\lies{$17 + 8 = 25$}{seventeen plus eight equals twenty-five}{siebzehn plus acht gleich fünfundzwanzig}
\lies{$42 - 15 = 27$}{forty-two minus fifteen is twenty-seven}{zweiundvierzig minus fünfzehn ist siebenundzwanzig}
\lies{$6 \cdot 7 = 42$}{six times seven is forty-two \;/\; six sevens are forty-two}{sechs mal sieben ist zweiundvierzig}
\lies{$56 : 8 = 7$}{fifty-six divided by eight is seven}{sechsundfünfzig geteilt durch acht ist sieben}
\lies{$(3 + 4) \cdot 5 = 35$}{the sum of three and four, multiplied by five, equals thirty-five \;/\; open bracket three plus four close bracket times five}{die Summe aus drei und vier, mal fünf, ist fünfunddreißig \;/\; Klammer auf drei plus vier Klammer zu mal fünf}
\end{lesen}

%--------------------------------------------------------------------------------------------------
\newpage
\section{Rechnen mit Brüchen}

\begin{lesen}
\lies{$\dfrac{12}{18} = \dfrac{2}{3}$}{twelve eighteenths equals two thirds. \newline \textit{Rule:} To reduce a fraction to its lowest terms, divide the numerator and the denominator by their highest common factor (here 6).}{zwölf Achtzehntel gleich zwei Drittel. \newline \textit{Regel:} Man kürzt einen Bruch, indem man Zähler und Nenner durch dieselbe Zahl dividiert (hier 6).}
\lies{$\dfrac{3}{4} = \dfrac{15}{20}$}{three quarters equals fifteen twentieths. \newline \textit{Rule:} To expand a fraction, multiply the numerator and the denominator by the same number (here 5).}{drei Viertel gleich fünfzehn Zwanzigstel. \newline \textit{Regel:} Man erweitert einen Bruch, indem man Zähler und Nenner mit derselben Zahl multipliziert (hier 5).}
\lies{$\dfrac{2}{3} + \dfrac{1}{4} = \dfrac{8}{12} + \dfrac{3}{12} = \dfrac{11}{12}$}{two thirds plus one quarter equals eight twelfths plus three twelfths equals eleven twelfths. \newline \textit{Rule:} To add (or subtract) unlike fractions, change them into like fractions with the lowest common denominator and then add (or subtract) the numerators.}{zwei Drittel plus ein Viertel gleich acht Zwölftel plus drei Zwölftel gleich elf Zwölftel. \newline \textit{Regel:} Ungleichnamige Brüche macht man zuerst gleichnamig (Hauptnenner) und addiert bzw.\ subtrahiert dann die Zähler.}
\lies{$\dfrac{2}{3} \cdot \dfrac{9}{10} = \dfrac{18}{30} = \dfrac{3}{5}$}{two thirds times nine tenths equals eighteen thirtieths equals three fifths. \newline \textit{Rule:} To multiply two fractions, multiply the numerators together and the denominators together.}{zwei Drittel mal neun Zehntel gleich achtzehn Dreißigstel gleich drei Fünftel. \newline \textit{Regel:} Man multipliziert Zähler mit Zähler und Nenner mit Nenner.}
\lies{$\dfrac{3}{4} : \dfrac{9}{8} = \dfrac{3}{4} \cdot \dfrac{8}{9} = \dfrac{2}{3}$}{three quarters divided by nine eighths equals three quarters times eight ninths equals two thirds. \newline \textit{Rule:} To divide by a fraction, multiply by its reciprocal.}{drei Viertel geteilt durch neun Achtel gleich drei Viertel mal acht Neuntel gleich zwei Drittel. \newline \textit{Regel:} Man dividiert durch einen Bruch, indem man mit seinem Kehrwert multipliziert.}
\lies{$\dfrac{11}{4} = 2\tfrac{3}{4} = 2 + \dfrac{3}{4}$}{eleven quarters equals two and three quarters. \newline \textit{Rule:} To convert an improper fraction into a mixed number, split it into the sum of a whole number and a proper fraction.}{elf Viertel gleich zwei drei Viertel. \newline \textit{Regel:} Einen unechten Bruch schreibt man als gemischte Zahl, indem man ihn in eine ganze Zahl und einen echten Bruch zerlegt.}
\end{lesen}

\anmerkung{„Cancel“ heißt wörtlich „streichen“: \textit{Cancel 6 out of the fraction} -- „Kürze den
Bruch durch 6“. Bei gemischten Zahlen sagt man im Englischen immer \textit{and}: $2\tfrac34$ --
\textit{two and three quarters}.}

%--------------------------------------------------------------------------------------------------

\platz\subsection*{Beispiele}

\begin{lesen}
\lies{$\dfrac56 - \dfrac14 = \dfrac{10}{12} - \dfrac{3}{12} = \dfrac{7}{12}$}{five sixths minus one quarter equals ten twelfths minus three twelfths, which is seven twelfths}{fünf Sechstel minus ein Viertel ist gleich zehn Zwölftel minus drei Zwölftel, also sieben Zwölftel}
\lies{$\dfrac45 \cdot 10 = 8$}{four fifths of ten is eight}{vier Fünftel von zehn sind acht}
\lies{$2\tfrac12 = \dfrac52$}{two and a half equals five halves}{zweieinhalb ist gleich fünf Halbe}
\lies{$\dfrac38 = 0{,}375 = 37{,}5\,\%$}{three eighths equals nought point three seven five or thirty-seven point five per cent}{drei Achtel ist gleich null Komma drei sieben fünf oder siebenunddreißig Komma fünf Prozent}
\end{lesen}

\newpage
\section{Potenzen, Wurzeln und Logarithmen}

Aus der Gleichung $b^n = c$ lassen sich drei Fragen stellen -- je nachdem, welche Größe gesucht ist:

\begin{center}
\begin{tikzpicture}[font=\small]
	\node[font=\huge] (f) at (0.6,0) {$b^{\,n} = c$};
	\coordinate (b) at ($(f.west)+(0.3,-0.15)$);
	\coordinate (n) at ($(f.north west)+(0.62,-0.2)$);
	\coordinate (c) at ($(f.east)+(-0.2,0.05)$);
	% Namen
	\node[align=center, blue!70!black] (nb) at (-2.4,1.3) {base\\[-2pt]\color{black!60}Basis};
	\node[align=center, blue!70!black] (nn) at (0.2,2.2) {exponent\\[-2pt]\color{black!60}Exponent};
	\node[align=center, blue!70!black] (nc) at (3.9,1.3) {power (value of the power)\\[-2pt]\color{black!60}Potenz(wert)};
	\draw[-{Stealth}, blue!70!black, shorten >=4pt] (nb) -- (b);
	\draw[-{Stealth}, blue!70!black, shorten >=4pt] (nn) -- (n);
	\draw[-{Stealth}, blue!70!black, shorten >=4pt] (nc) -- (c);
	% drei Rechenarten
	\node[align=center] (p) at (-5,-2.2) {\textbf{$c$ gesucht:} $c = b^n$\\ raising to a power\\\color{black!60}Potenzieren};
	\node[align=center] (r) at (0.6,-2.6) {\textbf{$b$ gesucht:} $b = \sqrt[n]{c}$\\ extracting a root\\\color{black!60}Radizieren};
	\node[align=center] (l) at (6,-2.2) {\textbf{$n$ gesucht:} $n = \log_b c$\\ taking the logarithm\\\color{black!60}Logarithmieren};
	\draw[red, thick, -{Stealth}] (0.6,-0.5) -- (p.north east);
	\draw[red, thick, -{Stealth}] (0.6,-0.5) -- (r.north);
	\draw[red, thick, -{Stealth}] (0.6,-0.5) -- (l.north west);
\end{tikzpicture}
\end{center}

\platz\subsection*{Potenzen}
\begin{lesen}
\lies{$b^n$}{$b$ to the power of $n$ \;/\; $b$ to the $n$th (power) \;/\; the $n$th power of $b$}{$b$ hoch $n$}
\lies{$5^2 = 25$}{five squared is twenty-five}{fünf Quadrat \;/\; fünf hoch zwei ist fünfundzwanzig}
\lies{$3^3 = 27$}{three cubed is twenty-seven}{drei hoch drei ist siebenundzwanzig}
\lies{$2^5 = 32$}{two to the power of five equals thirty-two \;/\; two to the fifth equals thirty-two}{zwei hoch fünf gleich zweiunddreißig}
\lies{$b^{-n} = \dfrac{1}{b^n}$}{$b$ to the power of minus $n$ equals one over $b$ to the $n$}{$b$ hoch minus $n$ gleich eins durch $b$ hoch $n$}
\lies{$10^{-3} = 0{,}001$}{ten to the minus three equals nought point nought nought one}{zehn hoch minus drei gleich null Komma null null eins}
\end{lesen}

\platz\subsection*{Wurzeln}
\begin{lesen}
\lies{$\sqrt[n]{c} = b$}{the $n$th root of $c$ is (equal to) $b$}{die $n$-te Wurzel aus $c$ ist gleich $b$}
\lies{$\sqrt{49} = 7$}{the square root of forty-nine is seven}{die (Quadrat-)Wurzel aus neunundvierzig ist sieben}
\lies{$\sqrt[3]{8} = 2$}{the cube root of eight is two}{die dritte Wurzel aus acht ist zwei}
\lies{$\sqrt[4]{81} = 3$}{the fourth root of eighty-one is three}{die vierte Wurzel aus einundachtzig ist drei}
\lies{$\sqrt[n]{c^m} = c^{\frac{m}{n}}$}{the $n$th root of $c$ to the power of $m$ equals $c$ to the power of $m$ over $n$}{die $n$-te Wurzel aus $c$ hoch $m$ gleich $c$ hoch $m$ $n$-tel}
\lies{$\sqrt{a^2 + b^2}$}{the square root of $a$ squared plus $b$ squared}{die Wurzel aus $a$ Quadrat plus $b$ Quadrat}
\end{lesen}

\platz\subsection*{Logarithmen}
\begin{lesen}
\lies{$\log_b c = n$}{the logarithm of $c$ to the base $b$ is $n$ \;/\; log to the base $b$ of $c$ equals $n$}{der Logarithmus von $c$ zur Basis $b$ ist $n$}
\lies{$\log_2 32 = 5$}{log to the base two of thirty-two equals five}{Logarithmus von zweiunddreißig zur Basis zwei gleich fünf}
\lies{$\lg 1000 = 3$}{the common logarithm of one thousand is three \;/\; log one thousand equals three}{der Zehnerlogarithmus (dekadische Logarithmus) von tausend ist drei}
\lies{$\ln \mathrm{e} = 1$}{the natural logarithm of $e$ is one \;/\; log $e$ of $e$ equals one}{der natürliche Logarithmus von e ist eins}
\lies{$\log_b(x \cdot y) = \log_b x + \log_b y$}{the log of a product equals the sum of the logs}{der Logarithmus eines Produkts ist gleich der Summe der Logarithmen}
\end{lesen}

\anmerkung{$\ln x$ liest man im Englischen meist \textit{log $x$} oder \textit{ell en $x$}, $\lg x$ ist
im Englischen unüblich -- dort schreibt man $\log_{10} x$ oder nur $\log x$.}

%--------------------------------------------------------------------------------------------------

\platz\subsection*{Beispiele}

\begin{lesen}
\lies{$2^{10} = 1024$}{two to the power of ten equals one thousand and twenty-four}{zwei hoch zehn ist eintausendvierundzwanzig}
\lies{$a^m \cdot a^n = a^{m+n}$}{$a$ to the $m$ times $a$ to the $n$ equals $a$ to the $m$ plus $n$}{$a$ hoch $m$ mal $a$ hoch $n$ ist gleich $a$ hoch $m$ plus $n$}
\lies{$9^{\frac12} = \sqrt{9} = 3$}{nine to the power of one half equals the square root of nine, which is three}{neun hoch ein Halb ist gleich der Wurzel aus neun, also drei}
\lies{$\sqrt[3]{27} = 3$}{the cube root of twenty-seven is three}{die dritte Wurzel aus siebenundzwanzig ist drei}
\end{lesen}

\newpage
\section{Zahlbereiche}

\begin{center}
\begin{tikzpicture}[font=\small]
	\foreach \k/\col in {4/blue!6, 3/blue!12, 2/blue!18, 1/blue!24, 0/blue!32} {
		\filldraw[fill=\col, draw=blue, thick] ({0.8*\k},0) ellipse[x radius={1.15+1.2*\k}, y radius={0.62+0.42*\k}];
	}
	\node[font=\large] at (-0.35,0.2) {$\N$};       \node at (0.4,-0.2) {$7$};
	\node[font=\large] at (2.15,0.4) {$\Z$};        \node at (2.15,-0.3) {$-5$};
	\node[font=\large] at (4.15,0.4) {$\Q$};        \node at (4.15,-0.3) {$\frac12;\ 0{,}\overline{3}$};
	\node[font=\large] at (6.15,0.4) {$\R$};        \node at (6.15,-0.3) {$\sqrt2;\ \pi$};
	\node[font=\large] at (8.1,0.4) {$\C$};         \node at (8.1,-0.3) {$\mathrm{i};\ 2 + 3\mathrm{i}$};
\end{tikzpicture}
\end{center}

\begin{lesen}
\lies{$\N = \{0;\ 1;\ 2;\ \ldots\}$}{the set of natural numbers}{die Menge der natürlichen Zahlen}
\lies{$\Z = \{\ldots;\ -1;\ 0;\ 1;\ \ldots\}$}{the set of integers: the union of the positive integers, the negative integers and zero}{die Menge der ganzen Zahlen: positive und negative ganze Zahlen und die Null}
\lies{$\Q = \left\{\dfrac{p}{q} \;\middle|\; p, q \in \Z;\ q \neq 0\right\}$}{the set of all fractions $p$ over $q$, where $p$ and $q$ are integers and $q$ is not zero}{die Menge aller Brüche $p$ durch $q$, wobei $p$ und $q$ ganze Zahlen sind und $q$ ungleich null ist}
\lies{$\R$}{the set of real numbers: all terminating and recurring decimals (rational numbers) and all non-terminating, non-recurring decimals (irrational numbers)}{die Menge der reellen Zahlen: alle abbrechenden und periodischen Dezimalbrüche (rationale Zahlen) und alle nichtabbrechenden, nichtperiodischen Dezimalbrüche (irrationale Zahlen)}
\lies{$\C = \{a + b\,\mathrm{i} \mid a, b \in \R\}$ \newline mit $\mathrm{i}^2 = -1$}{the set of complex numbers: $a$ plus $b$ i, where $a$ and $b$ are real numbers and i is the imaginary unit}{die Menge der komplexen Zahlen: $a$ plus $b$ i, wobei $a$ und $b$ reelle Zahlen sind und i die imaginäre Einheit ist}
\lies{$\N \subset \Z \subset \Q \subset \R \subset \C$}{each number set is a proper subset of the next one}{jeder Zahlbereich ist eine echte Teilmenge des nächsten}
\end{lesen}

\anmerkung{Ob die Null zu den natürlichen Zahlen gehört, wird unterschiedlich gehandhabt. Im Englischen
ist man deshalb oft genauer: \textit{positive integers} ($1;\,2;\,3;\,\ldots$) bzw.\
\textit{non-negative integers} ($0;\,1;\,2;\,\ldots$). Ältere Bücher schreiben noch $N$, $I$, $F$, $R$, $C$ statt
der Doppelstrich-Buchstaben.}

%--------------------------------------------------------------------------------------------------

\platz\subsection*{Beispiele}

\begin{lesen}
\lies{$-3 \in \Z,\ -3 \notin \N$}{minus three is an integer, but it is not a natural number}{minus drei ist eine ganze Zahl, aber keine natürliche Zahl}
\lies{$\frac34 \in \Q$}{three quarters is a rational number}{drei Viertel ist eine rationale Zahl}
\lies{$\sqrt{2} \notin \Q$}{the square root of two is irrational}{die Wurzel aus zwei ist irrational}
\lies{$0{,}\overline{3} = \frac13$}{nought point three recurring equals one third}{null Komma Periode drei ist gleich ein Drittel}
\end{lesen}

\newpage
\section{Mengen}

\begin{lesen}
\lies{$a \in M$}{$a$ is an element of $M$ \;/\; $a$ is in $M$ \;/\; $a$ belongs to $M$}{$a$ ist Element von $M$}
\lies{$a \notin M$}{$a$ is not an element of $M$}{$a$ ist nicht Element von $M$}
\lies{$M = \{2;\ 4;\ 6\}$ \newline \footnotesize engl. $\{2, 4, 6\}$}{$M$ is the set with the elements two, four and six}{$M$ ist die Menge mit den Elementen zwei, vier und sechs}
\lies{$M = \emptyset$}{$M$ is the empty set \;/\; $M$ is a null set}{$M$ ist die leere Menge}
\lies{$A \subseteq B$}{$A$ is a subset of $B$}{$A$ ist (eine) Teilmenge von $B$}
\lies{$A \subset B$}{$A$ is a proper subset of $B$}{$A$ ist (eine) echte Teilmenge von $B$}
\lies{$A \cup B$}{$A$ union $B$ \;/\; the union of $A$ and $B$}{$A$ vereinigt $B$ \;/\; die Vereinigungsmenge von $A$ und $B$}
\lies{$A \cap B$}{$A$ intersection $B$ \;/\; the intersection of $A$ and $B$}{$A$ geschnitten $B$ \;/\; die Schnittmenge von $A$ und $B$}
\lies{$A \setminus B$}{$A$ minus $B$ \;/\; the difference of $A$ and $B$}{$A$ ohne $B$ \;/\; die Differenzmenge von $A$ und $B$}
\lies{$A \times B$}{$A$ cross $B$ \;/\; the Cartesian product of $A$ and $B$}{$A$ kreuz $B$ \;/\; das kartesische Produkt von $A$ und $B$}
\lies{$\{x \in \R \mid x < 4\}$}{the set of all real numbers $x$ such that $x$ is less than four}{die Menge aller reellen Zahlen $x$, für die gilt: $x$ ist kleiner als vier}
\end{lesen}

%--------------------------------------------------------------------------------------------------

\platz\subsection*{Beispiele}

\begin{lesen}
\lies{$A = \{1;\,2;\,3;\,4\}$ \newline $B = \{3;\,4;\,5\}$}{$A$ is the set with the elements one, two, three and four; $B$ is the set with the elements three, four and five}{$A$ ist die Menge mit den Elementen eins, zwei, drei und vier; $B$ ist die Menge mit den Elementen drei, vier und fünf}
\lies{$A \cap B = \{3;\,4\}$}{the intersection of $A$ and $B$ is the set containing three and four}{die Schnittmenge von $A$ und $B$ ist die Menge mit den Elementen drei und vier}
\lies{$A \cup B = \{1;\,2;\,3;\,4;\,5\}$}{the union of $A$ and $B$ contains the numbers from one to five}{die Vereinigungsmenge von $A$ und $B$ enthält die Zahlen von eins bis fünf}
\end{lesen}

\newpage
\section{Gleichungen und Ungleichungen lösen}

\platz\subsection*{Beispiel: eine lineare Gleichung}

\begin{lesen}
\lies{$\dfrac{2(x - 1)}{3} + \dfrac{1}{2} = \dfrac{5}{2}$}{\textit{Solve the equation.}}{\textit{Löse die Gleichung.}}
\lies{$4(x - 1) + 3 = 15$}{Multiply both sides by the lowest common denominator, six.}{Multipliziere beide Seiten mit dem Hauptnenner sechs.}
\lies{$4x - 4 + 3 = 15$}{Expand the brackets \;/\; remove the parentheses.}{Löse die Klammern auf (multipliziere aus).}
\lies{$4x - 1 = 15$}{Simplify \;/\; collect like terms.}{Fasse zusammen.}
\lies{$4x = 16$}{Add one to both sides.}{Addiere eins auf beiden Seiten.}
\lies{$x = 4$}{Divide both sides by the coefficient of $x$, four.}{Dividiere beide Seiten durch den Koeffizienten von $x$, also durch vier.}
\lies{$L = \{4\}$}{The solution set contains only one element, four. \;/\; The solution is $x$ equals four.}{Die Lösungsmenge enthält nur das Element vier. \;/\; Die Lösung ist $x$ gleich vier.}
\lies{\textit{Probe:} \newline $\frac{2 \cdot 3}{3} + \frac12 = \frac52$ \;\checkmark}{\textit{Check:} Substitute four for $x$: two plus a half equals five halves. \checkmark}{\textit{Probe:} Setze vier für $x$ ein: zwei plus ein Halb gleich fünf Halbe. \checkmark}
\end{lesen}

\platz\subsection*{Beispiel: eine lineare Ungleichung}

\begin{lesen}
\lies{$3 - 2x < 9$}{\textit{Solve the inequality.}}{\textit{Löse die Ungleichung.}}
\lies{$-2x < 6$}{Subtract three from both sides.}{Subtrahiere drei auf beiden Seiten.}
\lies{$x > -3$}{Divide both sides by minus two. \textbf{Remember:} when you multiply or divide by a negative number, the inequality sign is reversed.}{Dividiere beide Seiten durch minus zwei. \textbf{Achtung:} Beim Multiplizieren oder Dividieren mit einer negativen Zahl dreht sich das Relationszeichen um.}
\lies{$L = \{x \in \R \mid x > -3\}$}{The solution set contains all real numbers greater than minus three.}{Die Lösungsmenge enthält alle reellen Zahlen, die größer als minus drei sind.}
\end{lesen}

\begin{center}
\begin{tikzpicture}[scale=0.8, font=\small]
	\draw[achse] (-6.4,0) -- (3.6,0);
	\foreach \x in {-6,...,3} {\draw (\x,0.1) -- (\x,-0.1) node[below] {$\x$};}
	\draw[blue, very thick, -{Stealth}] (-3,0) -- (3.4,0);
	\filldraw[fill=white, draw=blue, very thick] (-3,0) circle (4pt);
	\node[above, blue] at (0.2,0.15) {$x > -3$};
\end{tikzpicture}
\end{center}

\platz\subsection*{Gleichungsarten}

\begin{lesen}
\lies{$3x - 7 = 11$}{a linear equation (an equation of the first degree): three $x$ minus seven equals eleven}{eine lineare Gleichung (Gleichung ersten Grades): drei $x$ minus sieben gleich elf}
\lies{$x^2 - 5x + 6 = 0$}{a quadratic equation (an equation of the second degree): $x$ squared minus five $x$ plus six equals zero}{eine quadratische Gleichung (Gleichung zweiten Grades): $x$ Quadrat minus fünf $x$ plus sechs gleich null}
\lies{$\sqrt{x + 3} = 5$}{a radical equation: the square root of $x$ plus three (the square root of the sum $x$ plus three) equals five}{eine Wurzelgleichung: die Wurzel aus $x$ plus drei (aus der Summe) gleich fünf}
\lies{$2^x = 64$}{an exponential equation: two to the power of $x$ equals sixty-four}{eine Exponentialgleichung: zwei hoch $x$ gleich vierundsechzig}
\lies{$\log_3 x = 4$}{a logarithmic equation: log to the base three of $x$ equals four}{eine logarithmische Gleichung: Logarithmus von $x$ zur Basis drei gleich vier}
\lies{$2 \sin x = 1$}{a trigonometric equation: two sine $x$ equals one}{eine trigonometrische Gleichung: zwei mal Sinus $x$ gleich eins}
\lies{$\dfrac{x - 3}{2x + 1} = \dfrac16$}{$x$ minus three, all over two $x$ plus one, equals one sixth}{$x$ minus drei, alles geteilt durch zwei $x$ plus eins, gleich ein Sechstel}
\end{lesen}

\anmerkung{Bei längeren Brüchen hilft im Englischen \textit{all over}, um zu zeigen, dass der ganze Term
im Zähler steht. Klammern liest man bei Bedarf mit: \textit{open bracket} \ldots\ \textit{close bracket}.}

%--------------------------------------------------------------------------------------------------

\platz\subsection*{Beispiele}

\begin{lesen}
\lies{$x^2 - 5x + 6 = 0$ \newline $(x - 2)(x - 3) = 0$}{Factorise the left-hand side.}{Zerlege die linke Seite in Faktoren.}
\lies{$x_1 = 2,\ x_2 = 3$}{A product is zero if one of its factors is zero. The solutions are two and three.}{Ein Produkt ist null, wenn einer der Faktoren null ist. Die Lösungen sind zwei und drei.}
\end{lesen}

\newpage
\section{Funktionen}

\begin{lesen}
\lies{$y = f(x)$}{$y$ equals $f$ of $x$ \;/\; $y$ is a function of $x$ ($y$ is the value of the function corresponding to $x$)}{$y$ gleich $f$ von $x$ \;/\; $y$ ist der Funktionswert zu $x$}
\lies{$x$}{the independent variable}{die unabhängige Variable}
\lies{$y$}{the dependent variable}{die abhängige Variable}
\lies{$f\colon x \mapsto x^2$}{$f$ maps $x$ to $x$ squared}{$f$ ordnet $x$ das Quadrat von $x$ zu \;/\; $f$: $x$ wird abgebildet auf $x$ Quadrat}
\lies{$f(x) = 2x^2 - 3$}{$f$ of $x$ equals two $x$ squared minus three}{$f$ von $x$ gleich zwei $x$ Quadrat minus drei}
\lies{$f(3) = 15$}{$f$ of three equals fifteen \;/\; the value of $f$ at three is fifteen}{$f$ von drei gleich fünfzehn \;/\; der Funktionswert an der Stelle drei ist fünfzehn}
\lies{$D_f,\ W_f$}{the domain (of definition) and the range of $f$}{der Definitionsbereich und der Wertebereich von $f$}
\end{lesen}

\platz\subsection*{Der Graph einer Funktion}

\begin{center}
\begin{tikzpicture}[font=\small, scale=1.05]
	\draw[achse] (-0.7,0) -- (8.6,0) node[below left] {$x$};
	\draw[achse] (0,-1.1) -- (0,5.4) node[below left] {$y$};
	\begin{scope}
		\clip (-0.7,-1.1) rectangle (8.3,5.2);
		\draw[blue, very thick, domain=-0.35:7.75, samples=80] plot (\x, {0.1*\x*\x*\x - 1.1*\x*\x + 3.2*\x + 0.3});
	\end{scope}
	\coordinate (P) at (6.5,2.0825);
	\draw[red, thick, dashed] (P) -- (6.5,0) (P) -- (0,2.0825);
	\fill[red] (P) circle (2.5pt) node[above left] {$P(x_P \mid y_P)$};
	\node[below] at (6.5,0) {$x_P$}; \node[left] at (0,2.0825) {$y_P$};
	\node[below left] at (0,0) {$O$};
	% Beschriftungen
	\node[align=left, anchor=west] (gr) at (8.0,4.6) {graph of $f$\\\color{black!60}Graph von $f$};
	\draw[-{Stealth}, thin] (gr.west) -- (7.6,4.45);
	\node[align=center] (ab) at (5.1,-1.3) {abscissa \ipa{\ae b\textprimstress sIs@}\\\color{black!60}Abszisse};
	\draw[-{Stealth}, thin] (ab.east) -- (6.35,-0.35);
	\node[align=center] (or) at (-2.0,3.7) {ordinate \ipa{\textprimstress O:dIn@t}\\\color{black!60}Ordinate};
	\draw[-{Stealth}, thin] (or.south) -- (-0.15,2.55);
	\node[align=center] (ur) at (1.6,-1.3) {origin \ipa{\textprimstress 6rIdZIn}\\\color{black!60}Ursprung};
	\draw[-{Stealth}, thin] (ur.north west) -- (-0.1,-0.2);
	\node[align=center] (xa) at (8.6,0.9) {$x$-axis\\\color{black!60}$x$-Achse};
	\node[align=center] (ya) at (1.3,5.3) {$y$-axis\\\color{black!60}$y$-Achse};
\end{tikzpicture}
\end{center}

\begin{lesen}
\lies{$P(x_P \mid y_P)$ \newline \footnotesize engl. $P(x_P, y_P)$}{the point $P$ with the coordinates $x_P$ and $y_P$ in the rectangular (Cartesian) coordinate system}{der Punkt $P$ mit den Koordinaten $x_P$ und $y_P$ im rechtwinkligen (kartesischen) Koordinatensystem}
\lies{$x_P$}{the abscissa (the $x$-coordinate) of $P$}{die Abszisse ($x$-Koordinate) von $P$}
\lies{$y_P$}{the ordinate (the $y$-coordinate) of $P$}{die Ordinate ($y$-Koordinate) von $P$}
\end{lesen}

\platz\subsection*{Einteilung der elementaren Funktionen}

\begin{center}
\begin{tikzpicture}[font=\footnotesize, every node/.style={align=left, anchor=west}]
	\node (root) at (0,0)      {\textbf{elementary functions}\\\color{black!60}elementare Funktionen};
	\node (alg) at (4.4,2.7)  {\textbf{algebraic functions}\\\color{black!60}algebraische Funktionen};
	\node (tra) at (4.4,-2.95){\textbf{transcendental functions}\\\color{black!60}transzendente Funktionen};
	\node (rat) at (9.2,3.7)  {rational functions\\\color{black!60}rationale Funktionen};
	\node (irr) at (9.2,1.7)  {irrational functions\\\color{black!60}irrationale Funktionen\\ $y = \sqrt{x}$};
	\node (pol) at (13.0,4.75) {polynomial functions\\\color{black!60}ganzrationale Funktionen\\ $y = x^3 - 2x$};
	\node (fra) at (13.0,2.75) {rational functions\\\color{black!60}gebrochenrationale F.\\ $y = \frac{x + 1}{x - 2}$};
	\node (exp) at (9.2,-0.2) {exponential functions\\\color{black!60}Exponentialfunktionen\\ $y = \mathrm{e}^x$};
	\node (log) at (9.2,-2.05) {logarithmic functions\\\color{black!60}Logarithmusfunktionen\\ $y = \ln x$};
	\node (tri) at (9.2,-3.9) {trigonometric functions\\\color{black!60}Winkelfunktionen\\ $y = \sin x$};
	\node (arc) at (9.2,-5.75) {inverse trig.\ functions\\\color{black!60}Arkusfunktionen\\ $y = \arcsin x$};
	\foreach \a/\b in {root/alg, root/tra, alg/rat, alg/irr, rat/pol, rat/fra, tra/exp, tra/log, tra/tri, tra/arc}
		{\draw[blue, thick] (\a.east) -- ++(0.35,0) |- (\b.west);}
\end{tikzpicture}
\end{center}

\begin{lesen}
\lies{$y = \displaystyle\sum_{k=0}^{4} a_k x^k$}{$y$ equals the sum of $a$ sub $k$ times $x$ to the power of $k$, for $k$ from zero to four}{$y$ gleich Summe über $a$ Index $k$ mal $x$ hoch $k$ für $k$ von null bis vier}
\lies{$y = \dfrac{7 + x}{2 - x^2}$}{$y$ equals the fraction with the numerator seven plus $x$ and the denominator two minus $x$ squared}{$y$ ist gleich dem Bruch mit dem Zähler sieben plus $x$ und dem Nenner zwei minus $x$ Quadrat}
\lies{$y = -\sqrt{r^2 - x^2}$}{$y$ equals the negative square root of $r$ squared minus $x$ squared}{$y$ gleich minus die Wurzel aus $r$ Quadrat minus $x$ Quadrat}
\lies{$y = \mathrm{e}^x$}{$y$ equals $e$ to the (power of) $x$}{$y$ gleich e hoch $x$}
\lies{$y = \ln x$}{$y$ equals log $x$ \;/\; $y$ equals ell en $x$}{$y$ gleich ln $x$ \;/\; $y$ gleich natürlicher Logarithmus von $x$}
\end{lesen}

%--------------------------------------------------------------------------------------------------

\platz\subsection*{Beispiele}

\begin{lesen}
\lies{$f(x) = 2x^2 - 3$ \newline $f(2) = 5$}{$f$ of two equals five: the point $P(2 \mid 5)$ lies on the graph of $f$}{$f$ von zwei ist fünf: Der Punkt $P(2 \mid 5)$ liegt auf dem Graphen von $f$}
\lies{$f(x) = \frac12 x + 1$}{a linear function with the gradient one half and the $y$-intercept one}{eine lineare Funktion mit der Steigung ein Halb und dem $y$-Achsenabschnitt eins}
\lies{$f(x) = 0$}{Find the zeros of $f$: for which values of $x$ is $f$ of $x$ equal to zero?}{Bestimme die Nullstellen von $f$: Für welche $x$ ist $f$ von $x$ gleich null?}
\lies{$f(x) = g(x)$}{Find the points of intersection of the two graphs.}{Bestimme die Schnittpunkte der beiden Graphen.}
\end{lesen}

\newpage
\section{Winkelfunktionen}

\begin{center}
\begin{tikzpicture}[scale=2.1, font=\small]
	\draw[achse] (-1.25,0) -- (1.5,0) node[below left] {$x$};
	\draw[achse] (0,-1.2) -- (0,1.35) node[below left] {$y$};
	\draw[blue, thick] (0,0) circle (1);
	\coordinate (P) at (35:1);
	\draw[mass] (0.3,0) arc[start angle=0, end angle=35, radius=0.3]; \node at (17:0.42) {$\alpha$};
	\draw[blue, very thick] (0,0) -- (P);
	\draw[red, very thick] (P) -- (P |- 0,0) node[midway, left, fill=white, inner sep=1pt] {$\sin\alpha$};
	\draw[green!50!black, very thick] (0,0) -- (P |- 0,0) node[midway, below] {$\cos\alpha$};
	\draw[violet, very thick] (1,0) -- (1,{tan(35)}) node[midway, right] {$\tan\alpha$};
	\draw[violet, dashed] (P) -- (1,{tan(35)});
	\fill[blue] (P) circle (1pt) node[above left] {$P$};
	\node[below] at (1,0) {$1$};
	\node[align=left, anchor=west, font=\footnotesize] at (1.55,0.9) {unit circle\\\color{black!60}Einheitskreis};
\end{tikzpicture}
\end{center}

\begin{lesen}
\lies{$\sin x$}{the sine \ipa{saIn} of $x$ \;/\; sine $x$}{Sinus $x$}
\lies{$\cos x$}{the cosine \ipa{"k@UsaIn} of $x$}{Kosinus $x$}
\lies{$\tan x$}{the tangent \ipa{"t\ae ndZ@nt} of $x$ \;/\; tan $x$}{Tangens $x$}
\lies{$\cot x$}{the cotangent \ipa{""k@U"t\ae ndZ@nt} of $x$}{Kotangens $x$}
\lies{$\sec x,\ \csc x$}{the secant \ipa{"si:k@nt} and the cosecant \ipa{""k@U"si:k@nt} of $x$ (in English textbooks also cosec $x$)}{Sekans und Kosekans von $x$ (im deutschen Unterricht kaum verwendet)}
\lies{$\arcsin x$ \newline \footnotesize engl. auch $\sin^{-1} x$}{the inverse sine of $x$ \;/\; arc sine $x$ \;/\; the angle whose sine is $x$}{Arkussinus $x$ \;/\; der Winkel, dessen Sinus $x$ ist}
\lies{$\sin 30\degree = \dfrac12$}{the sine of thirty degrees is one half}{der Sinus von dreißig Grad ist ein Halb}
\lies{$\tan 45\degree = 1$}{the tangent of forty-five degrees equals one}{Tangens fünfundvierzig Grad gleich eins}
\lies{$\sin^2 x + \cos^2 x = 1$}{sine squared $x$ plus cosine squared $x$ equals one}{Sinus Quadrat $x$ plus Kosinus Quadrat $x$ gleich eins}
\lies{$5\sin^2 x - \cos 2x = 0$}{five times the sine squared of $x$, minus the cosine of two $x$, equals zero}{fünf mal Sinus Quadrat $x$ minus Kosinus zwei $x$ ist gleich null}
\end{lesen}

\anmerkung{Vorsicht bei $\sin^{-1} x$: Gemeint ist die Umkehrfunktion $\arcsin x$, nicht der Kehrwert
$\frac{1}{\sin x}$. Auf Taschenrechnern steht dafür ebenfalls oft \textsf{SIN$^{-1}$}.}

%--------------------------------------------------------------------------------------------------

\platz\subsection*{Beispiele}

\begin{lesen}
\lies{$\cos 60\degree = \dfrac12$}{the cosine of sixty degrees is one half}{der Kosinus von sechzig Grad ist ein Halb}
\lies{$\sin x = 0{,}5$ \newline $x_1 = 30\degree,\ x_2 = 150\degree$}{For angles between nought and three hundred and sixty degrees the equation has two solutions: thirty degrees and one hundred and fifty degrees.}{Für Winkel zwischen null und dreihundertsechzig Grad hat die Gleichung zwei Lösungen: dreißig Grad und hundertfünfzig Grad.}
\lies{$\sin(x + 360\degree) = \sin x$}{The sine function is periodic with a period of three hundred and sixty degrees.}{Die Sinusfunktion ist periodisch mit der Periode dreihundertsechzig Grad.}
\lies{$-1 \le \sin x \le 1$}{The values of the sine function lie between minus one and one.}{Die Werte der Sinusfunktion liegen zwischen minus eins und eins.}
\end{lesen}

\newpage
\section{Intervalle und Grenzwerte}

\begin{lesen}
\lies{\intervall{o}{o} \newline $]a;\,b[$ \quad engl. $(a, b)$}{the open interval from $a$ to $b$: $a < x < b$}{das offene Intervall von $a$ bis $b$: $a < x < b$}
\lies{\intervall{g}{g} \newline $[a;\,b]$ \quad engl. $[a, b]$}{the closed interval from $a$ to $b$: $a \le x \le b$}{das abgeschlossene Intervall von $a$ bis $b$: $a \le x \le b$}
\lies{\intervall{o}{g} \newline $]a;\,b]$ \quad engl. $(a, b]$}{the half-open interval from $a$ to $b$, open on the left and closed on the right: $a < x \le b$}{das halboffene Intervall von $a$ bis $b$, links offen und rechts abgeschlossen: $a < x \le b$}
\lies{\intervallR \newline $[a;\,\infty[$ \quad engl. $[a, \infty)$}{the interval from $a$ to infinity: $x \ge a$}{das Intervall von $a$ bis unendlich: $x \ge a$}
\lies{$\R = \ ]{-\infty};\,\infty[$}{the set of real numbers is the open interval from minus infinity to plus infinity}{die Menge der reellen Zahlen ist das offene Intervall von minus unendlich bis plus unendlich}
\end{lesen}

\anmerkung{Im Deutschen sind auch runde Klammern für offene Intervalle verbreitet: $(a;\,b)$. Für
Monotonie-Angaben schreiben wir die Intervalle als Ungleichungskette, z.\,B.\ \textit{fallend:}
$-\infty < x \le 2$.}

\begin{lesen}
\lies{$x \to x_0$}{$x$ tends to $x$ nought \;/\; $x$ approaches $x$ nought}{$x$ gegen $x$ null \;/\; $x$ strebt gegen $x$ null}
\lies{$\limxto{2} x^2 = 4$}{the limit of $x$ squared as $x$ tends to two is four}{der Grenzwert (Limes) von $x$ Quadrat für $x$ gegen zwei ist vier}
\lies{$\limxto{x_1} f(x) = L$}{the limit of $f$ of $x$ as $x$ tends to $x$ one is capital $L$ -- $f(x)$ approaches $L$ as $x$ tends to $x_1$ in any way}{Limes $f$ von $x$ für $x$ gegen $x$ eins ist groß $L$ -- $f(x)$ nähert sich $L$, wenn $x$ beliebig gegen $x_1$ strebt}
\lies{$\limxto{x_0} f(x) \neq f(x_0)$}{the limit of $f$ of $x$ as $x$ tends to $x$ nought is not equal to $f$ of $x$ nought: the function $f$ is discontinuous at $x_0$}{der Grenzwert von $f$ für $x$ gegen $x$ null ist ungleich $f$ von $x$ null: $f$ ist an der Stelle $x_0$ unstetig}
\lies{$\limn \dfrac{1}{n} = 0$}{the limit of one over $n$ as $n$ tends to infinity is zero: $\left(\frac1n\right)$ is a null sequence}{der Grenzwert von eins durch $n$ für $n$ gegen unendlich ist null: $\left(\frac1n\right)$ ist eine Nullfolge}
\lies{$\limn a_n = g$}{the sequence $(a_n)$ converges to $g$ \;/\; the sequence is convergent with the limit $g$}{die Folge $(a_n)$ konvergiert gegen $g$ \;/\; sie ist konvergent mit dem Grenzwert $g$}
\end{lesen}

%--------------------------------------------------------------------------------------------------

\platz\subsection*{Beispiele}

\begin{lesen}
\lies{$x \in [2;\,5]$}{$x$ is an element of the closed interval from two to five}{$x$ ist Element des abgeschlossenen Intervalls von zwei bis fünf}
\lies{$\lim\limits_{x \to \infty} \dfrac{1}{x} = 0$}{the limit of one over $x$ as $x$ tends to infinity is zero}{der Grenzwert von eins durch $x$ für $x$ gegen unendlich ist null}
\lies{$\limn \dfrac{2n + 1}{n} = 2$}{the limit of two $n$ plus one, all over $n$, as $n$ tends to infinity is two}{der Grenzwert von zwei $n$ plus eins, geteilt durch $n$, für $n$ gegen unendlich ist zwei}
\lies{steigend: $2 \le x < \infty$}{the function is increasing for all $x$ greater than or equal to two}{die Funktion ist steigend für alle $x$, die größer oder gleich zwei sind}
\end{lesen}

\newpage
\section{Differenzialrechnung}

\begin{center}
\begin{tikzpicture}[scale=1.25, font=\small]
	\draw[achse] (-0.4,0) -- (5,0) node[below left] {$x$};
	\draw[achse] (0,-0.4) -- (0,4.4) node[below left] {$y$};
	\begin{scope}
		\clip (-0.4,-0.4) rectangle (4.8,4.2);
		\draw[blue, very thick, domain=0:4.6, samples=60] plot (\x, {0.2*\x*\x + 0.4});
		% Sekante durch P(1|0,6) und Q(3,5|2,85), Steigung 0,9
		\draw[red, thick, dashed, domain=0:4.8] plot (\x, {0.9*(\x-1) + 0.6});
		% Tangente in P, Steigung 0,4
		\draw[red, very thick, domain=0:4.8] plot (\x, {0.4*(\x-1) + 0.6});
	\end{scope}
	\fill (1,0.6) circle (2pt) node[above left] {$P$};
	\fill (3.5,2.85) circle (2pt) node[above left] {$Q$};
	\draw[black!60] (1,0.6) -- node[below] {$h = \Delta x$} (3.5,0.6) -- node[right] {$\Delta y$} (3.5,2.85);
	\draw[black!40, dotted] (1,0.6) -- (1,0) node[below, black] {$x_0$};
	\draw[black!40, dotted] (3.5,0.6) -- (3.5,0) node[below, black] {$x_0 + h$};
	\node[red, align=left, anchor=west] at (4.85,3.8) {secant \ipa{\textprimstress si:k@nt} -- Sekante};
	\node[red, align=left, anchor=west] at (4.85,2.2) {tangent -- Tangente};
\end{tikzpicture}
\end{center}

\begin{lesen}
\lies{$\dfrac{f(x_0 + h) - f(x_0)}{h}$}{the difference quotient: the slope of the secant through $P$ and $Q$}{der Differenzenquotient: die Steigung der Sekante durch $P$ und $Q$}
\lies{$f'(x_0) = {}$ \newline $\displaystyle\lim_{h \to 0} \frac{f(x_0 + h) - f(x_0)}{h}$}{$f$ prime of $x$ nought is the limit, as $h$ tends to zero, of $f$ of $x$ nought plus $h$ minus $f$ of $x$ nought, all over $h$. Geometrically it is the slope of the tangent at $P$.}{$f$ Strich von $x$ null ist der Grenzwert für $h$ gegen null von $f$ von $x$ null plus $h$ minus $f$ von $x$ null, alles durch $h$. Geometrisch ist das die Steigung der Tangente in $P$.}
\lies{$f'(x),\ f''(x),\ f^{(n)}(x)$}{the first, second, $n$th derivative of $f$ of $x$: $f$ prime of $x$ (or: $f$ dash of $x$), $f$ double prime of $x$, $f$ $n$ prime of $x$}{die erste, zweite, $n$-te Ableitung von $f$: $f$ Strich von $x$, $f$ zwei Strich von $x$, $f$ $n$ Strich von $x$}
\lies{$\dfrac{\mathrm{d}y}{\mathrm{d}x}$}{d $y$ by d $x$ \;/\; d $y$ d $x$ -- the differential quotient of $y$}{d $y$ nach d $x$ -- der Differenzialquotient von $y$}
\lies{$\dfrac{\mathrm{d}^2 y}{\mathrm{d}x^2}$}{d two $y$ by d $x$ squared}{d zwei $y$ nach d $x$ Quadrat}
\lies{$\dfrac{\mathrm{d}^n y}{\mathrm{d}x^n}$}{d $n$ $y$ by d $x$ to the $n$th}{d $n$ $y$ nach d $x$ hoch $n$}
\lies{$\dfrac{\partial z}{\partial x}$}{partial d $z$ by d $x$ \;/\; the partial derivative of $z$ with respect to $x$}{d $z$ partiell nach d $x$}
\lies{$\dot{s} = \dfrac{\mathrm{d}s}{\mathrm{d}t}$}{$s$ dot equals d $s$ by d $t$ (derivative with respect to time)}{$s$ Punkt gleich d $s$ nach d $t$ (Ableitung nach der Zeit)}
\end{lesen}

\platz\subsection*{Beispiele}

\begin{lesen}
\lies{$(x^3)' = 3x^2$}{the derivative of $x$ cubed is three $x$ squared}{die Ableitung von $x$ hoch drei ist drei $x$ Quadrat}
\lies{$(\sin x)' = \cos x$}{the (first) derivative of the sine of $x$ equals the cosine of $x$}{die (erste) Ableitung von Sinus $x$ ist Kosinus $x$}
\lies{$f(x) = x^2 - 4x$ \newline $f'(x) = 2x - 4$}{Differentiate $f$: $f$ prime of $x$ equals two $x$ minus four.}{Leite $f$ ab: $f$ Strich von $x$ gleich zwei $x$ minus vier.}
\lies{$2x - 4 = 0 \ \Rightarrow\ x_E = 2$}{Set the first derivative equal to zero and solve for $x$.}{Setze die erste Ableitung gleich null und löse nach $x$ auf.}
\lies{$f''(2) = 2 > 0$}{The second derivative at two is positive, so the graph has a minimum (a turning point) there: $T(2 \mid -4)$.}{Die zweite Ableitung an der Stelle zwei ist positiv, also hat der Graph dort einen Tiefpunkt: $T(2 \mid -4)$.}
\end{lesen}

%--------------------------------------------------------------------------------------------------
\newpage
\section{Integralrechnung}

\begin{center}
\begin{tikzpicture}[scale=1.15, font=\small]
	% Untersumme (grau) und Obersumme (rot gestrichelt) für f(x) = 0,15x^2 + 0,6 auf [1; 5], n = 8
	\foreach \i in {0,...,7} {
		\pgfmathsetmacro{\xl}{1+0.5*\i} \pgfmathsetmacro{\xr}{1.5+0.5*\i}
		\pgfmathsetmacro{\yl}{0.15*\xl*\xl+0.6} \pgfmathsetmacro{\yr}{0.15*\xr*\xr+0.6}
		\filldraw[fill=gray!35, draw=gray] (\xl,0) rectangle (\xr,\yl);
		\draw[red, dashed] (\xl,\yl) -- (\xl,\yr) -- (\xr,\yr) -- (\xr,\yl);
	}
	\draw[achse] (-0.4,0) -- (6.2,0) node[below left] {$x$};
	\draw[achse] (0,-0.4) -- (0,5) node[below left] {$y$};
	\draw[blue, very thick, domain=0:5.6, samples=60] plot (\x, {0.15*\x*\x + 0.6}) node[right] {$f$};
	\draw[black!60] (1,0) -- (1,0.75); \draw[black!60] (5,0) -- (5,4.35);
	\node[below] at (1,0) {$a$}; \node[below] at (5,0) {$b$};
	\draw[|<->|, black!70] (3,-0.5) -- node[below] {$\Delta x$} (3.5,-0.5);
	\node[align=left, anchor=west] at (6.4,1.2) {\textcolor{gray}{\rule{8pt}{8pt}} lower sum -- Untersumme\\[2pt] \textcolor{red}{\rule[2pt]{8pt}{0.8pt}} upper sum -- Obersumme};
\end{tikzpicture}
\end{center}

\begin{lesen}
\lies{$\displaystyle\int f(x)\dx = F(x) + C$, \newline wenn $F'(x) = f(x)$}{the integral of $f$ of $x$ d $x$ equals capital $F$ of $x$ plus capital $C$, if capital $F$ prime of $x$ is equal to $f$ of $x$}{Integral $f$ von $x$ d $x$ ist gleich groß $F$ von $x$ plus groß $C$, wenn groß $F$ Strich von $x$ gleich $f$ von $x$ ist}
\lies{$\displaystyle\int f(x)\dx$}{the indefinite integral of $f$; $F$ is an antiderivative of $f$, $C$ is the constant of integration}{das unbestimmte Integral von $f$; $F$ ist eine Stammfunktion von $f$, $C$ die Integrationskonstante}
\lies{$\displaystyle\int_a^b f(x)\dx$}{the integral from $a$ to $b$ of $f$ of $x$ d $x$ -- a definite integral with the lower limit $a$ and the upper limit $b$}{Integral von $a$ bis $b$ über $f$ von $x$ d $x$ -- ein bestimmtes Integral mit der unteren Grenze $a$ und der oberen Grenze $b$}
\lies{$\big[F(x)\big]_a^b = F(b) - F(a)$}{capital $F$ of $x$ between the limits $a$ and $b$ \;/\; $F$ of $x$ evaluated from $a$ to $b$ equals $F$ of $b$ minus $F$ of $a$}{groß $F$ von $x$ in den Grenzen von $a$ bis $b$ ist gleich $F$ von $b$ minus $F$ von $a$}
\lies{$\displaystyle\lim_{\Delta x \to 0} \sum_{k} f(x_k)\,\Delta x$}{the limit, for delta $x$ tending to zero, of the sum of $f$ of $x$ sub $k$ times delta $x$}{der Grenzwert für Delta $x$ gegen null der Summe über $f$ von $x$ Index $k$ mal Delta $x$}
\lies{$\displaystyle\int_a^b f(x)\dx = A$}{for $f(x) \ge 0$: the area bounded by the graph of $f$, the $x$-axis and the lines $x = a$ and $x = b$}{für $f(x) \ge 0$: der Inhalt der Fläche zwischen dem Graphen von $f$, der $x$-Achse und den Geraden $x = a$ und $x = b$}
\end{lesen}

\platz\subsection*{Beispiele}

\begin{lesen}
\lies{$\displaystyle\int \cos x \dx = \sin x + C$}{the integral of cosine $x$ d $x$ equals sine $x$ plus $C$}{Integral Kosinus $x$ d $x$ gleich Sinus $x$ plus $C$}
\lies{$\displaystyle\int_0^2 x^2 \dx = \left[\tfrac13 x^3\right]_0^2 = \tfrac83$}{the integral from nought to two of $x$ squared d $x$ equals one third $x$ cubed between nought and two, which is eight thirds}{Integral von null bis zwei über $x$ Quadrat d $x$ gleich ein Drittel $x$ hoch drei in den Grenzen von null bis zwei, also acht Drittel}
\lies{$\displaystyle\int_1^{\mathrm{e}} \frac{1}{x} \dx = \big[\ln x\big]_1^{\mathrm{e}} = 1$}{the integral from one to $e$ of one over $x$ d $x$ equals log $x$ from one to $e$, which is one}{Integral von eins bis e über eins durch $x$ d $x$ gleich ln $x$ in den Grenzen von eins bis e, also eins}
\end{lesen}

\anmerkung{Wichtige Anwendungen der Integralrechnung in der Geometrie -- \textit{important applications of
integral calculus in geometry}: Flächeninhalte (\textit{areas}), Volumina (\textit{volumes}
\ipa{\textprimstress v6lju:mz}) und Bogenlängen (\textit{arc lengths} \ipa{A:k}).}

%--------------------------------------------------------------------------------------------------
\newpage
\section{Geraden und Winkel}

\begin{center}
\begin{tikzpicture}[font=\small]
	\draw[mass] (1.3,0) arc[start angle=0, end angle=40, radius=1.3]; \node at (20:0.9) {$\alpha$};
	\draw[blue, very thick] (4.6,0) -- (0,0) -- (40:4.6);
	\fill (0,0) circle (2pt) node[above left] {$S$}; \fill (3.2,0) circle (2pt) node[below] {$A$}; \fill (40:3.2) circle (2pt) node[above left] {$B$};
	\node[align=center] (v) at (-2.2,-0.9) {vertex\\\color{black!60}Scheitel};
	\draw[-{Stealth}, thin] (v.east) -- (-0.15,-0.1);
	\node[align=center] (l) at (6.6,0.9) {arms (legs)\\\color{black!60}Schenkel};
	\draw[-{Stealth}, thin] (l.west) -- (4.3,0.1); \draw[-{Stealth}, thin] (l.west) -- ($(40:4.3)+(0.15,-0.1)$);
\end{tikzpicture}
\end{center}

\begin{lesen}
\lies{$g$}{the (straight) line $g$}{die Gerade $g$}
\lies{$\overline{AB}$}{the line segment $A$ $B$ \;/\; the distance from $A$ to $B$}{die Strecke $A$ $B$}
\lies{$g \parallel h$}{$g$ is parallel to $h$}{$g$ ist parallel zu $h$}
\lies{$g \perp h$}{$g$ is perpendicular to $h$ \;/\; $g$ is at right angles to $h$}{$g$ ist senkrecht zu $h$ \;/\; $g$ steht senkrecht auf $h$}
\lies{$\sphericalangle ASB$ \newline \footnotesize engl. $\angle ASB$}{angle $A$ $S$ $B$ (the vertex is named in the middle)}{Winkel $A$ $S$ $B$ (der Scheitel steht in der Mitte)}
\lies{$\alpha = 40\degree$}{alpha equals forty degrees: an acute angle}{Alpha gleich vierzig Grad: ein spitzer Winkel}
\lies{$\alpha + \beta = 90\degree$}{alpha and beta are complementary angles: they add up to ninety degrees}{Alpha und Beta sind Komplementwinkel: Sie ergänzen sich zu neunzig Grad}
\lies{$\alpha + \beta = 180\degree$}{alpha and beta are supplementary angles: they add up to one hundred and eighty degrees}{Alpha und Beta sind Supplementwinkel: Sie ergänzen sich zu hundertachtzig Grad}
\end{lesen}

\platz\subsection*{Winkel an geschnittenen Parallelen}

\begin{center}
\begin{tikzpicture}[font=\small, scale=1.1]
	\coordinate (S1) at (0.6,2.2); \coordinate (S2) at (-0.67,0);
	\draw[red, thin] (S1) circle (0.42); \draw[red, thin] (S2) circle (0.42);
	\draw[blue, very thick] (-4,2.2) -- (4,2.2) node[right] {$g_1$};
	\draw[blue, very thick] (-4,0) -- (4,0) node[right] {$g_2$};
	\draw[blue, very thick] ($(S2)+(240:1.5)$) -- ($(S1)+(60:1.5)$) node[right] {$t$};
	\foreach \w/\n in {30/\alpha, 120/\beta, 210/\gamma, 300/\delta} {
		\node at ($(S1)+(\w:0.78)$) {$\n_1$}; \node at ($(S2)+(\w:0.78)$) {$\n_2$};}
	\node[align=left, anchor=west] at (5,1.1) {$g_1 \parallel g_2$\\ $t$: transversal\\\color{black!60}Transversale (schneidende Gerade)};
\end{tikzpicture}
\end{center}

\begin{lesen}
\lies{$\alpha_1$ und $\alpha_2$}{corresponding angles -- equal if the lines are parallel}{Stufenwinkel -- an Parallelen gleich groß}
\lies{$\gamma_1$ und $\alpha_2$}{alternate angles -- equal if the lines are parallel}{Wechselwinkel -- an Parallelen gleich groß}
\lies{$\delta_1$ und $\alpha_2$}{co-interior (allied) angles -- they add up to $180\degree$ if the lines are parallel}{entgegengesetzt liegende Winkel -- ergänzen sich an Parallelen zu $180\degree$}
\lies{$\alpha_1$ und $\gamma_1$}{vertically opposite angles -- always equal}{Scheitelwinkel -- immer gleich groß}
\lies{$\alpha_1$ und $\beta_1$}{adjacent angles on a straight line -- they add up to $180\degree$}{Nebenwinkel -- ergänzen sich zu $180\degree$}
\end{lesen}

\platz\subsection*{Beispiele}

\begin{lesen}
\lies{$\alpha_1 = 65\degree$ \newline $\beta_1 = 180\degree - 65\degree = 115\degree$}{Alpha one is sixty-five degrees. Beta one is its adjacent angle on the straight line, so beta one is one hundred and fifteen degrees.}{Alpha eins beträgt fünfundsechzig Grad. Beta eins ist der Nebenwinkel, also hundertfünfzehn Grad.}
\lies{$\gamma_1 = \alpha_1 = 65\degree$}{Gamma one and alpha one are vertically opposite angles, so they are equal.}{Gamma eins und Alpha eins sind Scheitelwinkel, also gleich groß.}
\lies{$\alpha_2 = \alpha_1 = 65\degree$}{Alpha two and alpha one are corresponding angles on parallel lines, so they are equal.}{Alpha zwei und Alpha eins sind Stufenwinkel an Parallelen, also gleich groß.}
\lies{$\alpha = 90\degree$}{The lines $g$ and $h$ intersect at right angles.}{Die Geraden $g$ und $h$ schneiden einander rechtwinklig.}
\end{lesen}

%--------------------------------------------------------------------------------------------------
\newpage
\section{Dreiecke, Vierecke und Vielecke}

\begin{center}
\begin{tikzpicture}[font=\small, scale=1.1]
	\draw[mass] (0.7,0) arc[start angle=0, end angle=59.04, radius=0.7]; \node at (28:0.45) {$\alpha$};
	\draw[mass] ($(5,0)+(136.85:0.7)$) arc[start angle=136.85, end angle=180, radius=0.7]; \node at ($(5,0)+(160:1.0)$) {$\beta$};
	\draw[mass] ($(1.8,3)+(239.04:0.6)$) arc[start angle=239.04, end angle=316.85, radius=0.6]; \node at ($(1.8,3)+(280:0.9)$) {$\gamma$};
	\draw[mass] (5.9,0) arc[start angle=0, end angle=136.85, radius=0.9]; \node at ($(5,0)+(60:1.2)$) {$\delta$};
	\draw[black!60, thick, dashed] (1.8,3) -- node[right, pos=0.6] {$h_c$} (1.8,0);
	\draw[black!60, thick] (2.15,0) arc[start angle=0, end angle=90, radius=0.35]; \fill[black!60] (1.95,0.15) circle (0.8pt);
	\draw[blue, very thick] (0,0) -- node[below] {$c$} (5,0) -- node[above right] {$a$} (1.8,3) -- node[above left] {$b$} cycle;
	\draw[blue, thick, dashed] (5,0) -- (6.6,0);
	\node[below left] at (0,0) {$A$}; \node[below] at (5,0) {$B$}; \node[above] at (1.8,3) {$C$};
	\node[align=left, anchor=west] at (7.4,2.2) {$a,\ b,\ c$: sides -- \color{black!60}Seiten};
	\node[align=left, anchor=west] at (7.4,1.6) {$\alpha,\ \beta,\ \gamma$: interior angles -- \color{black!60}Innenwinkel};
	\node[align=left, anchor=west] at (7.4,1.0) {$\delta$: exterior angle -- \color{black!60}Außenwinkel};
	\node[align=left, anchor=west] at (7.4,0.4) {$h_c$: height (altitude) -- \color{black!60}Höhe};
\end{tikzpicture}
\end{center}

\begin{lesen}
\lies{$\triangle ABC$}{triangle $A$ $B$ $C$}{Dreieck $A$ $B$ $C$}
\lies{$\alpha + \beta + \gamma = 180\degree$}{the interior angles of a triangle add up to one hundred and eighty degrees}{die Innenwinkel eines Dreiecks ergeben zusammen hundertachtzig Grad}
\lies{$\delta = \alpha + \gamma$}{an exterior angle equals the sum of the two opposite interior angles}{ein Außenwinkel ist so groß wie die beiden nicht anliegenden Innenwinkel zusammen}
\lies{$A = \dfrac{1}{2} \cdot c \cdot h_c$}{the area equals half the base times the height}{der Flächeninhalt ist gleich ein Halb mal Grundseite mal Höhe}
\lies{$a^2 + b^2 = c^2$}{$a$ squared plus $b$ squared equals $c$ squared -- Pythagoras' theorem: in a right-angled triangle the square on the hypotenuse equals the sum of the squares on the other two sides}{$a$ Quadrat plus $b$ Quadrat gleich $c$ Quadrat -- Satz des Pythagoras: Im rechtwinkligen Dreieck ist das Hypotenusenquadrat so groß wie die beiden Kathetenquadrate zusammen}
\lies{$\triangle ABC \cong \triangle DEF$}{triangle $A$ $B$ $C$ is congruent to triangle $D$ $E$ $F$}{Dreieck $A$ $B$ $C$ ist kongruent zu Dreieck $D$ $E$ $F$}
\lies{$\triangle ABC \sim \triangle PQR$}{triangle $A$ $B$ $C$ is similar to triangle $P$ $Q$ $R$}{Dreieck $A$ $B$ $C$ ist ähnlich zu Dreieck $P$ $Q$ $R$}
\end{lesen}

\platz\subsection*{Besondere Dreiecke}

\begin{lesen}
\lies{\skGleichschenklig}{an isosceles triangle: two equal sides (the legs), the third side is the base; the two base angles are equal}{ein gleichschenkliges Dreieck: zwei gleich lange Seiten (Schenkel), die dritte Seite ist die Basis; die Basiswinkel sind gleich groß}
\lies{\skRechtwinklig}{a right-angled triangle (AE: right triangle): the two legs enclose the right angle, the hypotenuse lies opposite it}{ein rechtwinkliges Dreieck: Die Katheten schließen den rechten Winkel ein, die Hypotenuse liegt ihm gegenüber}
\end{lesen}

\begin{center}
\begin{tikzpicture}[font=\small, scale=1.0]
	\coordinate (A) at (0,0); \coordinate (B) at (3.6,0); \coordinate (C) at (1.8,3.118); \coordinate (M) at (1.8,1.039);
	\draw[red, thick] (M) circle (1.039); \draw[red, thick] (M) circle (2.078);
	\foreach \p/\q in {A/{2.7,1.559}, B/{0.9,1.559}, C/{1.8,0}} {\draw[black!60, dash dot] ($(\p)!-0.25!(\q)$) -- ($(\p)!1.45!(\q)$);}
	\draw[blue, very thick] (A) -- (B) -- (C) -- cycle; \fill (M) circle (2pt) node[right=2pt] {$M$};
	\node[align=left, anchor=west] at (5.2,3.1) {an equilateral triangle with its three axes of symmetry\\\color{black!60}ein gleichseitiges Dreieck mit seinen drei Symmetrieachsen};
	\node[align=left, anchor=west] at (5.2,1.9) {\textcolor{red}{incircle} and \textcolor{red}{circumcircle} --
		\color{black!60}Inkreis und Umkreis};
	\node[align=left, anchor=west, text width=10cm] at (5.2,0.2) {$M$ is the point of intersection of the altitudes, the perpendicular bisectors, the medians and the angle bisectors.\\\color{black!60}$M$ ist der Schnittpunkt der Höhen, der Mittelsenkrechten, der Seitenhalbierenden und der Winkelhalbierenden.};
\end{tikzpicture}
\end{center}

\platz\subsection*{Vierecke und regelmäßige Vielecke}

\begin{lesen}
\lies{\skTrapez}{a trapezium (AE: trapezoid): a quadrilateral with one pair of parallel sides}{ein Trapez: ein Viereck mit einem Paar paralleler Seiten}
\lies{\skParallelogramm}{a parallelogram: a quadrilateral whose opposite sides are parallel}{ein Parallelogramm: ein Viereck, dessen gegenüberliegende Seiten parallel sind}
\lies{\skRaute}{a rhombus: a parallelogram with four equal sides}{eine Raute (ein Rhombus): ein Parallelogramm mit vier gleich langen Seiten}
\lies{\skRechteck}{a rectangle: a quadrilateral with four right angles}{ein Rechteck: ein Viereck mit vier rechten Winkeln}
\lies{\skQuadrat}{a square: a rectangle with four equal sides}{ein Quadrat: ein Rechteck mit vier gleich langen Seiten}
\lies{\skNeck{5} \skNeck{6} \skNeck{8}}{regular polygons: a pentagon, a hexagon and an octagon (all sides and all angles are equal)}{regelmäßige Vielecke: ein Fünfeck, ein Sechseck und ein Achteck (alle Seiten und alle Winkel sind gleich groß)}
\end{lesen}

\platz\subsection*{Beispiele}

\begin{lesen}
\lies{$a = 3\,\mathrm{cm},\ b = 4\,\mathrm{cm}$ \newline $c = \sqrt{3^2 + 4^2}\,\mathrm{cm} = 5\,\mathrm{cm}$}{In a right-angled triangle the legs are three and four centimetres long. The hypotenuse is the square root of three squared plus four squared, that is five centimetres.}{In einem rechtwinkligen Dreieck sind die Katheten drei und vier Zentimeter lang. Die Hypotenuse ist die Wurzel aus drei Quadrat plus vier Quadrat, also fünf Zentimeter.}
\lies{$A = \frac12 \cdot 8\,\mathrm{cm} \cdot 5\,\mathrm{cm} = 20\,\mathrm{cm}^2$}{The area of the triangle is half of eight centimetres times five centimetres, which is twenty square centimetres.}{Der Flächeninhalt des Dreiecks ist ein Halb mal acht Zentimeter mal fünf Zentimeter, also zwanzig Quadratzentimeter.}
\lies{$u = 2 \cdot 6\,\mathrm{cm} + 2 \cdot 4\,\mathrm{cm} = 20\,\mathrm{cm}$}{The perimeter of a rectangle six centimetres long and four centimetres wide is twenty centimetres.}{Der Umfang eines Rechtecks mit sechs Zentimetern Länge und vier Zentimetern Breite beträgt zwanzig Zentimeter.}
\lies{$(6 - 2) \cdot 180\degree = 720\degree$}{The interior angles of a hexagon add up to seven hundred and twenty degrees.}{Die Innenwinkel eines Sechsecks ergeben zusammen siebenhundertzwanzig Grad.}
\end{lesen}

%--------------------------------------------------------------------------------------------------
\newpage
\section{Der Kreis und seine Teile}

\begin{center}
\begin{tikzpicture}[font=\small, scale=1.2]
	\fill[gray!40] (120:2) arc[start angle=120, end angle=200, radius=2] -- cycle;
	\fill[blue!15] (0,0) -- (300:2) arc[start angle=300, end angle=345, radius=2] -- cycle;
	\draw[blue, very thick] (0,0) circle (2);
	\draw[blue, thick] (120:2) -- node[right, black, inner sep=2pt, pos=0.2] {$s$} (200:2);
	\draw[blue, thick] (300:2) -- (0,0) -- (345:2);
	\draw[black!70, thick] (90:2) -- node[left, pos=0.75] {$d$} (270:2);
	\draw[black!70, thick] (0,0) -- node[above] {$r$} (15:2);
	\draw[red, line width=2pt] (35:2) arc[start angle=35, end angle=75, radius=2]; \node[red] at (55:2.3) {$b$};
	\draw[red, thick] (-2.7,2) -- (2.7,2) node[right] {$t$};
	\fill[red] (0,2) circle (2pt) node[above] {$B$};
	\draw[red, thick] ($(215:2)!-0.5!(255:2)$) -- ($(215:2)!1.6!(255:2)$) node[right] {$g$};
	\fill (0,0) circle (2pt) node[left] {$M$};
	\node at (160:1.75) {$A_1$}; \node at (322:1.3) {$A_2$};
\end{tikzpicture}
\end{center}

\begin{lesen}
\lies{$M$}{the centre (AE: center)}{der Mittelpunkt}
\lies{$r$}{the radius \ipa{"reIdi@s} -- plural: radii \ipa{"reIdiaI}}{der Radius -- Plural: Radien}
\lies{$d = 2r$}{the diameter: the diameter is twice the radius}{der Durchmesser: Er ist doppelt so groß wie der Radius}
\lies{$s$}{a chord: a line segment joining two points on the circle}{eine Sehne: eine Strecke zwischen zwei Punkten des Kreises}
\lies{$g$}{a secant: a line cutting the circle in two points}{eine Sekante: eine Gerade, die den Kreis in zwei Punkten schneidet}
\lies{$t$, $B$}{a tangent: a line touching the circle at one point, the point of tangency (point of contact) $B$; the tangent is perpendicular to the radius at $B$}{eine Tangente: eine Gerade, die den Kreis in einem Punkt berührt, dem Berührungspunkt $B$; sie steht senkrecht auf dem Berührungsradius}
\lies{$b$}{an arc (of the circle)}{ein Kreisbogen}
\lies{$A_1$}{a segment: the region between a chord and an arc}{ein Kreisabschnitt (Kreissegment): die Fläche zwischen Sehne und Bogen}
\lies{$A_2$}{a sector: the region between two radii and an arc}{ein Kreisausschnitt (Kreissektor): die Fläche zwischen zwei Radien und einem Bogen}
\lies{$u = 2\pi r$}{the circumference equals two pi $r$}{der Umfang ist gleich zwei pi $r$}
\lies{$A = \pi r^2$}{the area equals pi $r$ squared}{der Flächeninhalt ist gleich pi $r$ Quadrat}
\end{lesen}

\platz\subsection*{Beispiele}

\begin{lesen}
\lies{$r = 5\,\mathrm{cm}$ \newline $u = 2\pi \cdot 5\,\mathrm{cm} \approx 31{,}4\,\mathrm{cm}$}{The radius is five centimetres, so the circumference is two pi times five centimetres, which is approximately thirty-one point four centimetres.}{Der Radius beträgt fünf Zentimeter, also ist der Umfang zwei pi mal fünf Zentimeter, ungefähr einunddreißig Komma vier Zentimeter.}
\lies{$A = \pi \cdot (5\,\mathrm{cm})^2 \approx 78{,}5\,\mathrm{cm}^2$}{The area is pi times five centimetres squared, approximately seventy-eight point five square centimetres.}{Der Flächeninhalt ist pi mal fünf Zentimeter zum Quadrat, ungefähr achtundsiebzig Komma fünf Quadratzentimeter.}
\lies{$b = \dfrac{60\degree}{360\degree} \cdot 2\pi \cdot 5\,\mathrm{cm} \approx 5{,}24\,\mathrm{cm}$}{The arc belonging to an angle of sixty degrees is one sixth of the circumference, about five point two four centimetres.}{Der Bogen zu einem Winkel von sechzig Grad ist ein Sechstel des Umfangs, etwa fünf Komma zwei vier Zentimeter.}
\end{lesen}

%--------------------------------------------------------------------------------------------------
\newpage
\section{Körper berechnen}

\begin{center}
\begin{tikzpicture}[font=\small]
	\def\dx{1.2}\def\dy{0.9}
	\draw[blue, thick, dashed] (0,0) -- (\dx,\dy) -- (4+\dx,\dy) (\dx,\dy) -- (\dx,2.5+\dy);
	\draw[red, thick, dashed] (0,0) -- (4+\dx,\dy);
	\draw[red, very thick] (0,0) -- (4+\dx,2.5+\dy);
	\draw[blue, very thick] (0,0) rectangle (4,2.5) (4,0) -- (4+\dx,\dy) -- (4+\dx,2.5+\dy) -- (\dx,2.5+\dy) -- (0,2.5) (4,2.5) -- (4+\dx,2.5+\dy);
	\node[below] at (2,0) {$a$}; \node[below right] at (4.55,0.4) {$b$}; \node[left] at (0,1.25) {$c$};
	\node[red, above left] at (2.4,1.6) {$d$}; \node[red, below] at (3.1,0.5) {$e$};
	\node[align=center] at (2.6,-0.9) {cuboid -- \color{black!60}Quader};
	\begin{scope}[shift={(9,0.2)}]
		\draw[blue, thick, dashed] (1.2,0) arc[start angle=0, end angle=180, x radius=1.2, y radius=0.35];
		\draw[blue, very thick] (-1.2,0) arc[start angle=180, end angle=360, x radius=1.2, y radius=0.35];
		\draw[blue, very thick] (0,2.8) ellipse[x radius=1.2, y radius=0.35];
		\draw[blue, very thick] (-1.2,0) -- (-1.2,2.8) (1.2,0) -- (1.2,2.8);
		\draw[red, thick] (0,2.8) -- node[above, inner sep=1.5pt] {$r$} (1.2,2.8); \fill (0,2.8) circle (1.5pt);
		\draw[red, thick, |<->|] (1.6,0) -- node[right] {$h$} (1.6,2.8);
		\node[align=center] at (0,-1.1) {cylinder -- \color{black!60}Zylinder};
	\end{scope}
\end{tikzpicture}
\end{center}

\begin{lesen}
\lies{Quader}{A cuboid has eight vertices (corners), twelve edges and six faces.}{Ein Quader hat acht Ecken, zwölf Kanten und sechs Flächen.}
\lies{$e$, $d$}{$e$ is a face diagonal, $d$ is a space diagonal (body diagonal)}{$e$ ist eine Flächendiagonale, $d$ ist eine Raumdiagonale}
\lies{$V = a \cdot b \cdot c$}{the volume equals length times width times height}{das Volumen ist gleich Länge mal Breite mal Höhe}
\lies{$O = 2(ab + ac + bc)$}{the surface area equals two times the sum of $a$ $b$, $a$ $c$ and $b$ $c$}{der Oberflächeninhalt ist gleich zwei mal Klammer auf $a$ $b$ plus $a$ $c$ plus $b$ $c$ Klammer zu}
\lies{$V = a^3$}{the volume of a cube equals $a$ cubed}{das Volumen eines Würfels ist gleich $a$ hoch drei}
\lies{$V = \pi r^2 h$}{the volume of a cylinder equals pi $r$ squared $h$ (base area times height)}{das Volumen eines Zylinders ist gleich pi $r$ Quadrat $h$ (Grundfläche mal Höhe)}
\lies{$M = 2\pi r h$}{the curved surface area of a cylinder equals two pi $r$ $h$}{der Mantelflächeninhalt eines Zylinders ist gleich zwei pi $r$ $h$}
\lies{$V = \dfrac13 \cdot G \cdot h$}{the volume of a pyramid or a cone equals one third of the base area times the height}{das Volumen einer Pyramide oder eines Kegels ist gleich ein Drittel mal Grundfläche mal Höhe}
\lies{$V = \dfrac43 \pi r^3$}{the volume of a sphere equals four thirds pi $r$ cubed}{das Volumen einer Kugel ist gleich vier Drittel pi $r$ hoch drei}
\lies{$O = 4\pi r^2$}{the surface area of a sphere equals four pi $r$ squared}{der Oberflächeninhalt einer Kugel ist gleich vier pi $r$ Quadrat}
\end{lesen}

\platz\subsection*{Besondere Körper}

\begin{lesen}
\lies{\skPyramidenstumpf \; \skKegelstumpf}{a frustum of a pyramid and a frustum of a cone (a truncated cone): the top has been cut off parallel to the base}{ein Pyramidenstumpf und ein Kegelstumpf: Die Spitze wurde parallel zur Grundfläche abgeschnitten}
\lies{\skTetraeder}{a tetrahedron: a regular polyhedron whose four faces are equilateral triangles}{ein (regelmäßiges) Tetraeder: ein Polyeder, dessen vier Flächen gleichseitige Dreiecke sind}
\lies{\skRotation}{a solid of revolution: it is generated by rotating a plane region about an axis -- here a semicircle generates a sphere}{ein Rotationskörper: Er entsteht, wenn eine ebene Fläche um eine Achse rotiert -- hier erzeugt ein Halbkreis eine Kugel}
\end{lesen}

\platz\subsection*{Beispiele}

\begin{lesen}
\lies{$a = 4\,\mathrm{cm},\ b = 3\,\mathrm{cm},\ c = 2\,\mathrm{cm}$ \newline $V = 24\,\mathrm{cm}^3$}{The volume of the cuboid is four times three times two, that is twenty-four cubic centimetres.}{Das Volumen des Quaders ist vier mal drei mal zwei, also vierundzwanzig Kubikzentimeter.}
\lies{$O = 2 \cdot (12 + 8 + 6)\,\mathrm{cm}^2 = 52\,\mathrm{cm}^2$}{Its surface area is fifty-two square centimetres.}{Sein Oberflächeninhalt beträgt zweiundfünfzig Quadratzentimeter.}
\lies{$d = \sqrt{4^2 + 3^2 + 2^2}\,\mathrm{cm} = \sqrt{29}\,\mathrm{cm} \approx 5{,}39\,\mathrm{cm}$}{The space diagonal is the square root of twenty-nine, approximately five point three nine centimetres.}{Die Raumdiagonale ist die Wurzel aus neunundzwanzig, ungefähr fünf Komma drei neun Zentimeter.}
\lies{$r = 3\,\mathrm{cm},\ h = 10\,\mathrm{cm}$ \newline $V = \pi \cdot 9 \cdot 10\,\mathrm{cm}^3 \approx 282{,}7\,\mathrm{cm}^3$}{The volume of the cylinder is ninety pi, approximately two hundred and eighty-two point seven cubic centimetres.}{Das Volumen des Zylinders ist neunzig pi, ungefähr zweihundertzweiundachtzig Komma sieben Kubikzentimeter.}
\lies{$r = 3\,\mathrm{cm}$ \newline $V = \frac43 \pi \cdot 27\,\mathrm{cm}^3 \approx 113{,}1\,\mathrm{cm}^3$}{The volume of the sphere is thirty-six pi, approximately one hundred and thirteen point one cubic centimetres.}{Das Volumen der Kugel ist sechsunddreißig pi, ungefähr hundertdreizehn Komma eins Kubikzentimeter.}
\end{lesen}

%--------------------------------------------------------------------------------------------------
\newpage
\section{Trigonometrie im Dreieck}

\begin{center}
\begin{tikzpicture}[font=\small, scale=1.1]
	\draw[mass] (1.3,0) arc[start angle=0, end angle=30.96, radius=1.3]; \node at (14:0.95) {$\alpha$};
	\draw[black!60, thick] (5,0.45) arc[start angle=90, end angle=180, radius=0.45]; \fill[black!60] (4.82,0.18) circle (1pt);
	\draw[mass] ($(5,3)+(210.96:0.8)$) arc[start angle=210.96, end angle=270, radius=0.8]; \node at ($(5,3)+(242:0.5)$) {$\beta$};
	\draw[blue, very thick] (0,0) -- (5,0) -- (5,3) -- cycle;
	\node[below left] at (0,0) {$A$}; \node[below right] at (5,0) {$C$}; \node[above right] at (5,3) {$B$};
	\node[align=center, above left] at (2.4,1.6) {$c$: hypotenuse\\\color{black!60}Hypotenuse};
	\node[align=left, right] at (5.15,1.5) {$a$: side opposite $\alpha$\\\color{black!60}Gegenkathete von $\alpha$};
	\node[align=center, below] at (2.5,-0.1) {$b$: side adjacent to $\alpha$\\\color{black!60}Ankathete von $\alpha$};
\end{tikzpicture}
\end{center}

\begin{lesen}
\lies{$\sin\alpha = \dfrac{a}{c}$}{sine alpha equals opposite over hypotenuse}{Sinus Alpha gleich Gegenkathete durch Hypotenuse}
\lies{$\cos\alpha = \dfrac{b}{c}$}{cosine alpha equals adjacent over hypotenuse}{Kosinus Alpha gleich Ankathete durch Hypotenuse}
\lies{$\tan\alpha = \dfrac{a}{b}$}{tan alpha equals opposite over adjacent}{Tangens Alpha gleich Gegenkathete durch Ankathete}
\lies{$\cot\alpha = \dfrac{b}{a}$}{cot alpha equals adjacent over opposite}{Kotangens Alpha gleich Ankathete durch Gegenkathete}
\lies{$\dfrac{a}{\sin\alpha} = \dfrac{b}{\sin\beta} = \dfrac{c}{\sin\gamma}$}{the sine rule (law of sines): $a$ over sine alpha equals $b$ over sine beta equals $c$ over sine gamma -- it holds in every triangle}{der Sinussatz: $a$ durch Sinus Alpha gleich $b$ durch Sinus Beta gleich $c$ durch Sinus Gamma -- er gilt in jedem Dreieck}
\lies{$c^2 = a^2 + b^2 - 2ab\cos\gamma$}{the cosine rule (law of cosines): $c$ squared equals $a$ squared plus $b$ squared minus two $a$ $b$ cosine gamma}{der Kosinussatz: $c$ Quadrat gleich $a$ Quadrat plus $b$ Quadrat minus zwei $a$ $b$ Kosinus Gamma}
\lies{$180\degree \mathrel{\widehat{=}} \pi$}{one hundred and eighty degrees corresponds to pi radians}{hundertachtzig Grad entsprechen pi im Bogenmaß}
\end{lesen}

\anmerkung{Englischsprachige Schülerinnen und Schüler merken sich die Definitionen mit dem Kunstwort
\textbf{SOH-CAH-TOA}: \textbf{S}ine = \textbf{O}pposite over \textbf{H}ypotenuse, \textbf{C}osine =
\textbf{A}djacent over \textbf{H}ypotenuse, \textbf{T}angent = \textbf{O}pposite over \textbf{A}djacent.
Die Tafelwerke mit Winkelfunktions- und Logarithmentafeln aus älteren Büchern ersetzt heute der
Taschenrechner (\textit{calculator}).}

\platz\subsection*{Beispiele}

\begin{lesen}
\lies{$a = 3\,\mathrm{cm},\ c = 5\,\mathrm{cm}$ \newline $\sin\alpha = \frac35 = 0{,}6$ \newline $\alpha \approx 36{,}9\degree$}{Sine alpha equals three over five, which is nought point six. So alpha is the inverse sine of nought point six, approximately thirty-six point nine degrees.}{Sinus Alpha ist drei durch fünf, also null Komma sechs. Damit ist Alpha der Arkussinus von null Komma sechs, ungefähr sechsunddreißig Komma neun Grad.}
\lies{$\sin 30{,}2\degree \approx 0{,}5030$}{the sine of thirty point two degrees is approximately nought point five nought three nought}{der Sinus von dreißig Komma zwei Grad ist ungefähr null Komma fünf null drei null}
\lies{$h = 20\,\mathrm{m} \cdot \tan 35\degree \approx 14{,}0\,\mathrm{m}$}{From a point twenty metres away the top of a tree is seen at an angle of elevation of thirty-five degrees. The tree is about fourteen metres high.}{Aus zwanzig Metern Entfernung sieht man die Spitze eines Baumes unter einem Höhenwinkel von fünfunddreißig Grad. Der Baum ist etwa vierzehn Meter hoch.}
\lies{$a = 5,\ b = 7,\ \gamma = 60\degree$ \newline $c^2 = 25 + 49 - 70 \cdot \frac12 = 39$ \newline $c \approx 6{,}24$}{By the cosine rule, $c$ squared equals twenty-five plus forty-nine minus seventy times one half, which is thirty-nine. So $c$ is approximately six point two four.}{Nach dem Kosinussatz ist $c$ Quadrat gleich fünfundzwanzig plus neunundvierzig minus siebzig mal ein Halb, also neununddreißig. Damit ist $c$ ungefähr sechs Komma zwei vier.}
\lies{$a = 8,\ \alpha = 30\degree,\ \beta = 45\degree$ \newline $b = \dfrac{8 \cdot \sin 45\degree}{\sin 30\degree} \approx 11{,}31$}{By the sine rule, $b$ equals eight times sine forty-five degrees over sine thirty degrees, approximately eleven point three one.}{Nach dem Sinussatz ist $b$ gleich acht mal Sinus fünfundvierzig Grad durch Sinus dreißig Grad, ungefähr elf Komma drei eins.}
\end{lesen}

%--------------------------------------------------------------------------------------------------
\newpage
\section{Koordinaten und Geraden}

\begin{center}
\begin{tikzpicture}[font=\small, scale=0.95]
	\draw[achse] (-0.7,0) -- (5.6,0) node[below left] {$x$}; \draw[achse] (0,-0.7) -- (0,4.2) node[below left] {$y$};
	\draw[mass] (1.2,0) arc[start angle=0, end angle=36.87, radius=1.2]; \node at (17:1.6) {$\varphi$};
	\draw[black!60, dashed] (4,0) node[below, black] {$x_P$} -- (4,3) -- (0,3) node[left, black] {$y_P$};
	\draw[red, very thick] (0,0) -- node[above left, inner sep=2pt] {$r$} (4,3);
	\fill (4,3) circle (2pt) node[above right] {$P$}; \node[below left] at (0,0) {$O$};
	\begin{scope}[shift={(9.6,0.3)}]
		\draw[achse] (-3.2,0) -- (6.2,0) node[below left] {$x$}; \draw[achse] (0,-0.9) -- (0,4.2) node[below left] {$y$};
		\draw[mass] (-1,0) arc[start angle=0, end angle=26.57, radius=1]; \node at ($(-2,0)+(12:1.4)$) {$\alpha$};
		\draw[blue, very thick] (-3,-0.5) -- (5.6,3.8) node[right] {$g$};
		\draw[black!60, thick] (2,2) -- node[below] {$\Delta x$} (4,2) -- node[right] {$\Delta y$} (4,3);
		\draw[red, thick, |<->|] (-0.3,0) -- node[left] {$n$} (-0.3,1);
		\fill[red] (0,1) circle (2pt);
	\end{scope}
\end{tikzpicture}
\end{center}

\begin{lesen}
\lies{$P(4 \mid 3)$ \newline \footnotesize engl. $P(4, 3)$}{the point $P$ four three: its $x$-coordinate (abscissa) is four, its $y$-coordinate (ordinate) is three}{der Punkt $P$ vier drei: Seine $x$-Koordinate (Abszisse) ist vier, seine $y$-Koordinate (Ordinate) ist drei}
\lies{$P(r;\,\varphi)$}{the polar coordinates of $P$: the distance $r$ from the origin (the pole) and the polar angle phi}{die Polarkoordinaten von $P$: der Abstand $r$ vom Ursprung (Pol) und der Polarwinkel Phi}
\lies{$x = r\cos\varphi,\ y = r\sin\varphi$}{$x$ equals $r$ cosine phi, $y$ equals $r$ sine phi}{$x$ gleich $r$ Kosinus Phi, $y$ gleich $r$ Sinus Phi}
\lies{$y = mx + n$}{the equation of a straight line in slope-intercept form: $y$ equals $m$ $x$ plus $n$}{die Normalform der Geradengleichung: $y$ gleich $m$ $x$ plus $n$}
\lies{$m = \dfrac{\Delta y}{\Delta x} = \tan\alpha$}{the gradient (slope) $m$ equals delta $y$ over delta $x$ (\glqq rise over run\grqq); it is the tangent of the angle of inclination alpha}{die Steigung $m$ ist gleich Delta $y$ durch Delta $x$; sie ist der Tangens des Steigungswinkels Alpha}
\lies{$n$}{the $y$-intercept: the line cuts the $y$-axis at the point $(0 \mid n)$}{der $y$-Achsenabschnitt: Die Gerade schneidet die $y$-Achse im Punkt $(0 \mid n)$}
\lies{$m_1 = m_2$}{parallel lines have the same gradient}{parallele Geraden haben die gleiche Steigung}
\lies{$m_1 \cdot m_2 = -1$}{two lines are perpendicular if the product of their gradients is minus one}{zwei Geraden sind zueinander senkrecht, wenn das Produkt ihrer Steigungen minus eins ist}
\end{lesen}

\platz\subsection*{Beispiele}

\begin{lesen}
\lies{$A(1 \mid 2),\ B(3 \mid 6)$ \newline $m = \dfrac{6 - 2}{3 - 1} = 2$}{Find the equation of the line through $A$ and $B$. The gradient is six minus two over three minus one, which is two.}{Bestimme die Gleichung der Geraden durch $A$ und $B$. Die Steigung ist sechs minus zwei durch drei minus eins, also zwei.}
\lies{$2 = 2 \cdot 1 + n \ \Rightarrow\ n = 0$ \newline $y = 2x$}{Substitute the coordinates of $A$: $n$ equals zero. The equation of the line is $y$ equals two $x$.}{Setze die Koordinaten von $A$ ein: $n$ ist null. Die Geradengleichung lautet $y$ gleich zwei $x$.}
\lies{$\overline{AB} = \sqrt{2^2 + 4^2} = \sqrt{20} \approx 4{,}47$}{The distance between $A$ and $B$ is the square root of twenty, approximately four point four seven.}{Der Abstand von $A$ und $B$ ist die Wurzel aus zwanzig, ungefähr vier Komma vier sieben.}
\lies{$2x = -x + 6 \ \Rightarrow\ x = 2$ \newline $S(2 \mid 4)$}{The lines $y$ equals two $x$ and $y$ equals minus $x$ plus six intersect at the point $S$ two four.}{Die Geraden $y$ gleich zwei $x$ und $y$ gleich minus $x$ plus sechs schneiden einander im Punkt $S$ zwei vier.}
\end{lesen}

%--------------------------------------------------------------------------------------------------
\newpage
\section{Kegelschnitte}

\begin{center}
\begin{tikzpicture}[font=\footnotesize, scale=0.95]
	\begin{scope}
		\draw[achse] (-1.7,0) -- (1.9,0) node[below] {$x$}; \draw[achse] (0,-1.7) -- (0,1.9) node[left] {$y$};
		\draw[blue, very thick] (0,0) circle (1.3);
		\draw[red, thick] (0,0) -- node[above left, inner sep=1pt] {$r$} (40:1.3); \fill (40:1.3) circle (1.5pt) node[above right] {$P$};
		\node[align=center] at (0,-2.3) {circle\\\color{black!60}Kreis};
	\end{scope}
	\begin{scope}[shift={(4.3,0)}]
		\draw[achse] (-2.1,0) -- (2.3,0) node[below] {$x$}; \draw[achse] (0,-1.7) -- (0,1.9) node[left] {$y$};
		\draw[blue, very thick] (0,0) ellipse[x radius=1.7, y radius=1.1];
		\coordinate (P) at ({1.7*cos(65)},{1.1*sin(65)});
		\draw[red, thick] (-1.296,0) -- (P) -- (1.296,0);
		\fill[red] (-1.296,0) circle (1.5pt) node[below] {$F_1$} (1.296,0) circle (1.5pt) node[below] {$F_2$};
		\fill (P) circle (1.5pt) node[above right] {$P$};
		\node[align=center] at (0,-2.3) {ellipse\\\color{black!60}Ellipse};
	\end{scope}
	\begin{scope}[shift={(8.9,0)}]
		\draw[achse] (-2.1,0) -- (2.3,0) node[below] {$x$}; \draw[achse] (0,-1.7) -- (0,1.9) node[left] {$y$};
		\draw[black!50, dashed] (-1.6,-1.6) -- (1.6,1.6) (-1.6,1.6) -- (1.6,-1.6);
		\draw[blue, very thick, domain=-1.5:1.5, samples=40, variable=\t] plot ({0.7*cosh(\t)}, {0.7*sinh(\t)});
		\draw[blue, very thick, domain=-1.5:1.5, samples=40, variable=\t] plot ({-0.7*cosh(\t)}, {0.7*sinh(\t)});
		\fill[red] (-0.99,0) circle (1.5pt) node[below left, inner sep=1pt] {$F_1$} (0.99,0) circle (1.5pt) node[below right, inner sep=1pt] {$F_2$};
		\node[align=center] at (0,-2.3) {hyperbola\\\color{black!60}Hyperbel};
	\end{scope}
	\begin{scope}[shift={(13.2,0)}]
		\draw[achse] (-1.2,0) -- (2.3,0) node[below] {$x$}; \draw[achse] (0,-1.7) -- (0,1.9) node[left] {$y$};
		\draw[blue, very thick, domain=-1.6:1.6, samples=40, variable=\t] plot ({\t*\t/1.6}, {\t});
		\draw[red, thick, dashed] (-0.4,-1.6) -- (-0.4,1.7) node[left] {$l$};
		\fill[red] (0.4,0) circle (1.5pt) node[below right, inner sep=1pt] {$F$};
		\node[align=center] at (0.5,-2.3) {parabola\\\color{black!60}Parabel};
	\end{scope}
\end{tikzpicture}
\end{center}

\begin{lesen}
\lies{$x^2 + y^2 = r^2$}{the equation of a circle with its centre at the origin and the radius $r$: $x$ squared plus $y$ squared equals $r$ squared}{die Gleichung eines Kreises um den Ursprung mit dem Radius $r$: $x$ Quadrat plus $y$ Quadrat gleich $r$ Quadrat}
\lies{$\dfrac{x^2}{a^2} + \dfrac{y^2}{b^2} = 1$}{the equation of an ellipse \ipa{I"lIps}: $x$ squared over $a$ squared plus $y$ squared over $b$ squared equals one; $a$ is the semi-major axis, $b$ the semi-minor axis}{die Gleichung einer Ellipse: $x$ Quadrat durch $a$ Quadrat plus $y$ Quadrat durch $b$ Quadrat gleich eins; $a$ ist die große, $b$ die kleine Halbachse}
\lies{$\overline{PF_1} + \overline{PF_2} = 2a$}{$F_1$ and $F_2$ are the foci \ipa{"f@UsaI} (singular: focus): for every point $P$ on the ellipse the sum of the distances from the foci is constant}{$F_1$ und $F_2$ sind die Brennpunkte: Für jeden Ellipsenpunkt $P$ ist die Summe der Abstände zu den Brennpunkten konstant}
\lies{$\dfrac{x^2}{a^2} - \dfrac{y^2}{b^2} = 1$}{the equation of a hyperbola \ipa{haI"p3:b@l@}: $x$ squared over $a$ squared minus $y$ squared over $b$ squared equals one; the dashed lines are its asymptotes}{die Gleichung einer Hyperbel: $x$ Quadrat durch $a$ Quadrat minus $y$ Quadrat durch $b$ Quadrat gleich eins; die gestrichelten Geraden sind ihre Asymptoten}
\lies{$y^2 = 2px$}{the equation of a parabola with its vertex at the origin: $y$ squared equals two $p$ $x$; $F$ is the focus, $l$ the directrix \ipa{d@"rektrIks}}{die Scheitelgleichung einer Parabel: $y$ Quadrat gleich zwei $p$ $x$; $F$ ist der Brennpunkt, $l$ die Leitlinie}
\lies{$\overline{PF} = \overline{Pl}$}{every point of the parabola is equidistant from the focus and the directrix}{jeder Parabelpunkt hat von Brennpunkt und Leitlinie den gleichen Abstand}
\end{lesen}

\anmerkung{Kreis, Ellipse, Parabel und Hyperbel heißen Kegelschnitte (\textit{conic sections}, kurz
\textit{conics}), weil sie entstehen, wenn eine Ebene einen Doppelkegel schneidet.}

\platz\subsection*{Beispiele}

\begin{lesen}
\lies{$x^2 + y^2 = 25$ \newline $P(3 \mid 4)$: $9 + 16 = 25$ \checkmark}{The point $P$ three four lies on the circle, because three squared plus four squared equals twenty-five.}{Der Punkt $P$ drei vier liegt auf dem Kreis, denn drei Quadrat plus vier Quadrat ist fünfundzwanzig.}
\lies{$\dfrac{x^2}{25} + \dfrac{y^2}{9} = 1$ \newline $e = \sqrt{25 - 9} = 4$}{The semi-axes of the ellipse are five and three. The foci are at the points four nought and minus four nought.}{Die Halbachsen der Ellipse sind fünf und drei. Die Brennpunkte liegen bei $(4 \mid 0)$ und $(-4 \mid 0)$.}
\lies{$y^2 = 8x$ \newline $p = 4$}{The focus of the parabola is the point two nought; the directrix is the line $x$ equals minus two.}{Der Brennpunkt der Parabel ist $(2 \mid 0)$; die Leitlinie ist die Gerade $x$ gleich minus zwei.}
\end{lesen}

%--------------------------------------------------------------------------------------------------
\newpage
\section{Folgen, Vektoren, Matrizen und Kombinatorik}

\platz\subsection*{Folgen und Reihen}
\begin{lesen}
\lies{$a_n = a_1 + (n - 1) \cdot d$}{the $n$th term of an arithmetic sequence: $a$ sub $n$ equals $a$ one plus $n$ minus one, times $d$ ($d$ is the common difference)}{das $n$-te Glied einer arithmetischen Folge: $a$ $n$ gleich $a$ eins plus $n$ minus eins mal $d$}
\lies{$a_n = a_1 \cdot q^{\,n-1}$}{the $n$th term of a geometric sequence: $a$ sub $n$ equals $a$ one times $q$ to the power of $n$ minus one ($q$ is the common ratio)}{das $n$-te Glied einer geometrischen Folge: $a$ $n$ gleich $a$ eins mal $q$ hoch $n$ minus eins}
\lies{$\displaystyle\sum_{k=1}^{n} k = \frac{n(n + 1)}{2}$}{the sum from $k$ equals one to $n$ of $k$ equals $n$ times $n$ plus one, all over two}{die Summe über $k$ für $k$ von eins bis $n$ ist gleich $n$ mal $n$ plus eins, geteilt durch zwei}
\lies{$\displaystyle\sum_{k=0}^{\infty} q^k = \frac{1}{1 - q}$ \newline für $|q| < 1$}{the infinite geometric series converges if the absolute value of $q$ is less than one; otherwise it diverges}{die unendliche geometrische Reihe konvergiert, wenn der Betrag von $q$ kleiner als eins ist; sonst divergiert sie}
\end{lesen}

\platz\subsection*{Vektoren und Matrizen}
\begin{lesen}
\lies{$\vec{a},\ |\vec{a}\,|$}{vector $a$; the magnitude (modulus) of vector $a$}{Vektor $a$; der Betrag von Vektor $a$}
\lies{$\vec{a} \cdot \vec{b}$}{$a$ dot $b$ -- the scalar product (dot product)}{$a$ mal $b$ -- das Skalarprodukt}
\lies{$\vec{a} \times \vec{b}$}{$a$ cross $b$ -- the vector product (cross product)}{$a$ kreuz $b$ -- das Vektorprodukt (Kreuzprodukt)}
\lies{$\vec{a} \perp \vec{b} \ \Leftrightarrow\ \vec{a} \cdot \vec{b} = 0$}{two vectors are perpendicular if and only if their scalar product is zero}{zwei Vektoren sind genau dann zueinander senkrecht, wenn ihr Skalarprodukt null ist}
\lies{$A = \begin{pmatrix} 1 & 2 & 3\\ 4 & 5 & 6 \end{pmatrix}$}{$A$ is a two by three matrix: it has two rows and three columns (plural of matrix: matrices \ipa{"meItrIsi:z})}{$A$ ist eine Zwei-Kreuz-Drei-Matrix: Sie hat zwei Zeilen und drei Spalten}
\lies{$a_{ij}$}{$a$ $i$ $j$: the element in row $i$ and column $j$}{$a$ $i$ $j$: das Element in Zeile $i$ und Spalte $j$}
\lies{$\det A$}{the determinant of $A$}{die Determinante von $A$}
\end{lesen}

\platz\subsection*{Kombinatorik}
\begin{lesen}
\lies{$n!$}{$n$ factorial: the number of permutations of $n$ different elements}{$n$ Fakultät: die Anzahl der Permutationen von $n$ verschiedenen Elementen}
\lies{$\dbinom{n}{k}$}{$n$ choose $k$: the number of combinations of $k$ elements chosen from $n$ (without repetition)}{$n$ über $k$: die Anzahl der Kombinationen von $k$ aus $n$ Elementen (ohne Wiederholung)}
\lies{$\dfrac{n!}{(n - k)!}$}{the number of arrangements of $k$ elements taken from $n$, where the order matters}{die Anzahl der Variationen von $k$ aus $n$ Elementen (mit Beachtung der Reihenfolge)}
\end{lesen}

\platz\subsection*{Beispiele}

\begin{lesen}
\lies{$3;\ 7;\ 11;\ 15;\ \ldots$ \newline $a_{10} = 3 + 9 \cdot 4 = 39$}{an arithmetic sequence with the first term three and the common difference four; the tenth term is thirty-nine}{eine arithmetische Folge mit dem Anfangsglied drei und der Differenz vier; das zehnte Glied ist neununddreißig}
\lies{$\vektor{1}{2}{3} \cdot \vektor{4}{-1}{2} = 4 - 2 + 6 = 8$}{the scalar product of the two vectors is eight}{das Skalarprodukt der beiden Vektoren ist acht}
\lies{$\begin{vmatrix} 1 & 2\\ 3 & 4 \end{vmatrix} = 1 \cdot 4 - 2 \cdot 3 = -2$}{the determinant equals one times four minus two times three, which is minus two}{die Determinante ist eins mal vier minus zwei mal drei, also minus zwei}
\lies{$4! = 4 \cdot 3 \cdot 2 \cdot 1 = 24$}{Four factorial is twenty-four: four books can be arranged on a shelf in twenty-four different ways.}{Vier Fakultät ist vierundzwanzig: Vier Bücher lassen sich auf vierundzwanzig Arten ins Regal stellen.}
\lies{$\dbinom{5}{2} = \dfrac{5 \cdot 4}{2} = 10$}{Five choose two is ten: there are ten ways of choosing two pupils from a group of five.}{Fünf über zwei ist zehn: Es gibt zehn Möglichkeiten, aus fünf Schülern zwei auszuwählen.}
\end{lesen}

%--------------------------------------------------------------------------------------------------
\newpage
\section{Das griechische Alphabet}

Griechische Buchstaben werden im Englischen anders ausgesprochen als im Deutschen -- besonders
auffällig bei $\pi$, $\varphi$, $\chi$ und $\psi$.

\begin{center}
{\small\renewcommand{\arraystretch}{1.25}
\begin{tabular}{@{}c c l l @{\hspace{1.3cm}} c c l l@{}}
	\toprule
	\multicolumn{2}{@{}l}{\textbf{Zeichen}} & \textbf{Englisch} & \textbf{Deutsch} & \multicolumn{2}{l}{\textbf{Zeichen}} & \textbf{Englisch} & \textbf{Deutsch} \\ \midrule
	$\mathrm{A}$ & $\alpha$ & alpha \ipa{"\ae lf@} & Alpha & $\mathrm{N}$ & $\nu$ & nu \ipa{nju:} & Ny \\
	$\mathrm{B}$ & $\beta$ & beta \ipa{"bi:t@} & Beta & $\Xi$ & $\xi$ & xi \ipa{ksaI} & Xi \\
	$\Gamma$ & $\gamma$ & gamma \ipa{"g\ae m@} & Gamma & $\mathrm{O}$ & $o$ & omicron \ipa{@U"maIkr6n} & Omikron \\
	$\Delta$ & $\delta$ & delta \ipa{"delt@} & Delta & $\Pi$ & $\pi$ & pi \ipa{paI} & Pi \\
	$\mathrm{E}$ & $\varepsilon$ & epsilon \ipa{"epsIl6n} & Epsilon & $\mathrm{P}$ & $\varrho$ & rho \ipa{r@U} & Rho \\
	$\mathrm{Z}$ & $\zeta$ & zeta \ipa{"zi:t@} & Zeta & $\Sigma$ & $\sigma$ & sigma \ipa{"sIgm@} & Sigma \\
	$\mathrm{H}$ & $\eta$ & eta \ipa{"i:t@} & Eta & $\mathrm{T}$ & $\tau$ & tau \ipa{taU} & Tau \\
	$\Theta$ & $\vartheta$ & theta \ipa{"Ti:t@} & Theta & $\Upsilon$ & $\upsilon$ & upsilon \ipa{ju:p"saIl@n} & Ypsilon \\
	$\mathrm{I}$ & $\iota$ & iota \ipa{aI"@Ut@} & Jota & $\Phi$ & $\varphi$ & phi \ipa{faI} & Phi \\
	$\mathrm{K}$ & $\kappa$ & kappa \ipa{"k\ae p@} & Kappa & $\mathrm{X}$ & $\chi$ & chi \ipa{kaI} & Chi \\
	$\Lambda$ & $\lambda$ & lambda \ipa{"l\ae md@} & Lambda & $\Psi$ & $\psi$ & psi \ipa{psaI} & Psi \\
	$\mathrm{M}$ & $\mu$ & mu \ipa{mju:} & My & $\Omega$ & $\omega$ & omega \ipa{"@UmIg@} & Omega \\
	\bottomrule
\end{tabular}}
\end{center}

\platz\subsection*{Beispiele}

\begin{lesen}
\lies{$\alpha,\ \beta,\ \gamma$}{alpha, beta and gamma are the interior angles of the triangle}{Alpha, Beta und Gamma sind die Innenwinkel des Dreiecks}
\lies{$\Delta x$}{delta $x$: the difference between two $x$-values}{Delta $x$: die Differenz zweier $x$-Werte}
\lies{$\varepsilon > 0$}{for every epsilon greater than zero}{für jedes Epsilon größer als null}
\lies{$\vec{x} = \vec{p} + \lambda \cdot \vec{u}$}{vector $x$ equals vector $p$ plus lambda times vector $u$}{Vektor $x$ gleich Vektor $p$ plus Lambda mal Vektor $u$}
\lies{$\mu = 3{,}5;\ \sigma \approx 1{,}71$}{the mean mu is three point five, the standard deviation sigma is approximately one point seven one}{der Erwartungswert My ist drei Komma fünf, die Standardabweichung Sigma ist ungefähr eins Komma sieben eins}
\lies{$\displaystyle\sum_{k=1}^{n} a_k$}{capital sigma: the sum of $a$ sub $k$ for $k$ from one to $n$}{groß Sigma: die Summe über $a$ $k$ für $k$ von eins bis $n$}
\lies{$\pi \approx 3{,}14$}{pi \ipa{paI} is approximately three point one four}{Pi ist ungefähr drei Komma eins vier}
\end{lesen}


\newpage
%%%%  REGISTER  %%%%
\part{Register}
\section{Register Deutsch -- Englisch}
\begin{multicols}{2}\small\raggedright

\medskip{\large\textbf{\color{cornflower}A}}\par\nopagebreak
\textbf{abhängige Variable} (die) -- dependent variable\dotfill\hyperref[e:107]{\pageref*{e:107}}\par
\textbf{ableiten} -- differentiate\dotfill\hyperref[e:278]{\pageref*{e:278}}\par
\textbf{Ableitung} (die) -- derivative\dotfill\hyperref[e:273]{\pageref*{e:273}}\par
\textbf{Abstand} (der) -- distance\dotfill\hyperref[e:152]{\pageref*{e:152}}\par
\textbf{Abszisse} (die) -- abscissa\dotfill\hyperref[e:117]{\pageref*{e:117}}\par
\textbf{achsensymmetrisch zur $y$-Achse} -- symmetric about the $y$-axis\dotfill\hyperref[e:144]{\pageref*{e:144}}\par
\textbf{Achteck} (das) -- octagon\dotfill\hyperref[e:207]{\pageref*{e:207}}\par
\textbf{addieren} -- add\dotfill\hyperref[e:14]{\pageref*{e:14}}\par
\textbf{ähnlich} -- similar\dotfill\hyperref[e:176]{\pageref*{e:176}}\par
\textbf{Anfangsglied} (das) -- first term\dotfill\hyperref[e:261]{\pageref*{e:261}}\par
\textbf{Ankathete} (die) -- adjacent side\dotfill\hyperref[e:188]{\pageref*{e:188}}\par
\textbf{arithmetische Folge} (die) -- arithmetic sequence\dotfill\hyperref[e:259]{\pageref*{e:259}}\par
\textbf{Arkusfunktion} (die) -- inverse trigonometric function\dotfill\hyperref[e:136]{\pageref*{e:136}}\par
\textbf{Asymptote} (die) -- asymptote\dotfill\hyperref[e:271]{\pageref*{e:271}}\par
\textbf{ausklammern} -- factorise\dotfill\hyperref[e:92]{\pageref*{e:92}}\par
\textbf{ausmultiplizieren} -- expand\dotfill\hyperref[e:91]{\pageref*{e:91}}\par
\textbf{Außenwinkel} (der) -- exterior angle\dotfill\hyperref[e:179]{\pageref*{e:179}}\par

\medskip{\large\textbf{\color{cornflower}B}}\par\nopagebreak
\textbf{Basis} (die) -- base\dotfill\hyperref[e:43]{\pageref*{e:43}}\par
\textbf{Basis (im Dreieck)} (die) -- base\dotfill\hyperref[e:183]{\pageref*{e:183}}\par
\textbf{Basiswinkel} (der) -- base angle\dotfill\hyperref[e:184]{\pageref*{e:184}}\par
\textbf{Baumdiagramm} (das) -- tree diagram\dotfill\hyperref[e:332]{\pageref*{e:332}}\par
\textbf{Berührungspunkt} (der) -- point of tangency\dotfill\hyperref[e:220]{\pageref*{e:220}}\par
\textbf{bestimmte Integral} (das) -- definite integral\dotfill\hyperref[e:288]{\pageref*{e:288}}\par
\textbf{Betrag} (der) -- absolute value\dotfill\hyperref[e:53]{\pageref*{e:53}}\par
\textbf{Betrag eines Vektors} (der) -- magnitude\dotfill\hyperref[e:305]{\pageref*{e:305}}\par
\textbf{Beweis} (der) -- proof\dotfill\hyperref[e:349]{\pageref*{e:349}}\par
\textbf{Binomialkoeffizient} (der) -- binomial coefficient\dotfill\hyperref[e:344]{\pageref*{e:344}}\par
\textbf{Binomialverteilung} (die) -- binomial distribution\dotfill\hyperref[e:345]{\pageref*{e:345}}\par
\textbf{binomische Formel} (die) -- binomial formula\dotfill\hyperref[e:93]{\pageref*{e:93}}\par
\textbf{Bogenmaß} (das) -- radian measure\dotfill\hyperref[e:139]{\pageref*{e:139}}\par
\textbf{Brennpunkt} (der) -- focus\dotfill\hyperref[e:316]{\pageref*{e:316}}\par
\textbf{Bruch} (der) -- fraction\dotfill\hyperref[e:28]{\pageref*{e:28}}\par

\medskip{\large\textbf{\color{cornflower}D}}\par\nopagebreak
\textbf{daraus folgt} -- implies\dotfill\hyperref[e:350]{\pageref*{e:350}}\par
\textbf{Deckfläche} (die) -- top face\dotfill\hyperref[e:249]{\pageref*{e:249}}\par
\textbf{Definition} (die) -- definition\dotfill\hyperref[e:346]{\pageref*{e:346}}\par
\textbf{Definitionsbereich} (der) -- domain\dotfill\hyperref[e:109]{\pageref*{e:109}}\par
\textbf{Determinante} (die) -- determinant\dotfill\hyperref[e:325]{\pageref*{e:325}}\par
\textbf{Dezimalzahl} (die) -- decimal number\dotfill\hyperref[e:13]{\pageref*{e:13}}\par
\textbf{Diagonale} (die) -- diagonal\dotfill\hyperref[e:197]{\pageref*{e:197}}\par
\textbf{Differenz} (die) -- difference\dotfill\hyperref[e:18]{\pageref*{e:18}}\par
\textbf{Differenzenquotient} (der) -- difference quotient\dotfill\hyperref[e:272]{\pageref*{e:272}}\par
\textbf{Differenzial} (das) -- differential\dotfill\hyperref[e:276]{\pageref*{e:276}}\par
\textbf{Differenzialquotient} (der) -- differential quotient\dotfill\hyperref[e:275]{\pageref*{e:275}}\par
\textbf{Differenzialrechnung} (die) -- differential calculus\dotfill\hyperref[e:274]{\pageref*{e:274}}\par
\textbf{Differenzmenge} (die) -- set difference\dotfill\hyperref[e:102]{\pageref*{e:102}}\par
\textbf{divergent} -- divergent\dotfill\hyperref[e:263]{\pageref*{e:263}}\par
\textbf{Dividend} (der) -- dividend\dotfill\hyperref[e:26]{\pageref*{e:26}}\par
\textbf{dividieren} -- divide\dotfill\hyperref[e:24]{\pageref*{e:24}}\par
\textbf{Divisor} (der) -- divisor\dotfill\hyperref[e:27]{\pageref*{e:27}}\par
\textbf{Drachenviereck} (das) -- kite\dotfill\hyperref[e:203]{\pageref*{e:203}}\par
\textbf{Drehung} (die) -- rotation\dotfill\hyperref[e:174]{\pageref*{e:174}}\par
\textbf{Dreieck} (das) -- triangle\dotfill\hyperref[e:177]{\pageref*{e:177}}\par
\textbf{Durchmesser} (der) -- diameter\dotfill\hyperref[e:214]{\pageref*{e:214}}\par

\medskip{\large\textbf{\color{cornflower}E}}\par\nopagebreak
\textbf{Ebene} (die) -- plane\dotfill\hyperref[e:308]{\pageref*{e:308}}\par
\textbf{echte Bruch} (der) -- proper fraction\dotfill\hyperref[e:31]{\pageref*{e:31}}\par
\textbf{echte Teilmenge} (die) -- proper subset\dotfill\hyperref[e:99]{\pageref*{e:99}}\par
\textbf{Ecke} (die) -- vertex\dotfill\hyperref[e:245]{\pageref*{e:245}}\par
\textbf{eckigen Klammern} (die) -- square brackets\dotfill\hyperref[e:89]{\pageref*{e:89}}\par
\textbf{Einheitskreis} (der) -- unit circle\dotfill\hyperref[e:137]{\pageref*{e:137}}\par
\textbf{Einheitsvektor} (der) -- unit vector\dotfill\hyperref[e:303]{\pageref*{e:303}}\par
\textbf{Element} (das) -- element\dotfill\hyperref[e:97]{\pageref*{e:97}}\par
\textbf{Ellipse} (die) -- ellipse\dotfill\hyperref[e:314]{\pageref*{e:314}}\par
\textbf{entgegengesetzt liegende Winkel} -- co-interior angles, allied angles\dotfill\hyperref[e:169]{\pageref*{e:169}}\par
\textbf{Ereignis} (das) -- event\dotfill\hyperref[e:328]{\pageref*{e:328}}\par
\textbf{Ergebnis} (das) -- outcome\dotfill\hyperref[e:327]{\pageref*{e:327}}\par
\textbf{Erwartungswert} (der) -- expected value\dotfill\hyperref[e:335]{\pageref*{e:335}}\par
\textbf{erweitern} -- expand (a fraction)\dotfill\hyperref[e:38]{\pageref*{e:38}}\par
\textbf{es gibt} -- there exists\dotfill\hyperref[e:353]{\pageref*{e:353}}\par
\textbf{eulersche Zahl} (die) -- Euler's number\dotfill\hyperref[e:298]{\pageref*{e:298}}\par
\textbf{Exponent} (der) -- exponent\dotfill\hyperref[e:44]{\pageref*{e:44}}\par
\textbf{Exponentialfunktion} (die) -- exponential function\dotfill\hyperref[e:130]{\pageref*{e:130}}\par
\textbf{Exponentialgleichung} (die) -- exponential equation\dotfill\hyperref[e:84]{\pageref*{e:84}}\par
\textbf{Extrempunkt} (der) -- extreme point\dotfill\hyperref[e:281]{\pageref*{e:281}}\par

\medskip{\large\textbf{\color{cornflower}F}}\par\nopagebreak
\textbf{Faktor} (der) -- factor\dotfill\hyperref[e:23]{\pageref*{e:23}}\par
\textbf{Fakultät} (die) -- factorial\dotfill\hyperref[e:340]{\pageref*{e:340}}\par
\textbf{Flächendiagonale} (die) -- face diagonal\dotfill\hyperref[e:251]{\pageref*{e:251}}\par
\textbf{flächengleich} -- equal in area\dotfill\hyperref[e:195]{\pageref*{e:195}}\par
\textbf{Flächeninhalt} (der) -- area\dotfill\hyperref[e:210]{\pageref*{e:210}}\par
\textbf{Folge} (die) -- sequence\dotfill\hyperref[e:257]{\pageref*{e:257}}\par
\textbf{Fünfeck} (das) -- pentagon\dotfill\hyperref[e:206]{\pageref*{e:206}}\par
\textbf{Funktion} (die) -- function\dotfill\hyperref[e:105]{\pageref*{e:105}}\par
\textbf{Funktionsgraph} (der) -- graph\dotfill\hyperref[e:111]{\pageref*{e:111}}\par
\textbf{Funktionsuntersuchung} (die) -- curve sketching\dotfill\hyperref[e:280]{\pageref*{e:280}}\par
\textbf{Funktionswert} (der) -- value of a function\dotfill\hyperref[e:108]{\pageref*{e:108}}\par
\textbf{für alle} -- for all\dotfill\hyperref[e:352]{\pageref*{e:352}}\par

\medskip{\large\textbf{\color{cornflower}G}}\par\nopagebreak
\textbf{ganze Zahl} (die) -- integer\dotfill\hyperref[e:6]{\pageref*{e:6}}\par
\textbf{ganzrationale Funktion} (die) -- polynomial function\dotfill\hyperref[e:126]{\pageref*{e:126}}\par
\textbf{gebrochenrationale Funktion} (die) -- rational function\dotfill\hyperref[e:127]{\pageref*{e:127}}\par
\textbf{Gegenereignis} (das) -- complementary event\dotfill\hyperref[e:329]{\pageref*{e:329}}\par
\textbf{Gegenkathete} (die) -- opposite side\dotfill\hyperref[e:189]{\pageref*{e:189}}\par
\textbf{gemischte Zahl} (die) -- mixed number\dotfill\hyperref[e:33]{\pageref*{e:33}}\par
\textbf{genau dann, wenn} -- if and only if\dotfill\hyperref[e:351]{\pageref*{e:351}}\par
\textbf{geometrische Folge} (die) -- geometric sequence\dotfill\hyperref[e:260]{\pageref*{e:260}}\par
\textbf{Gerade} (die) -- straight line\dotfill\hyperref[e:147]{\pageref*{e:147}}\par
\textbf{Geradengleichung} (die) -- equation of a straight line\dotfill\hyperref[e:309]{\pageref*{e:309}}\par
\textbf{geschweiften Klammern} (die) -- curly brackets, braces\dotfill\hyperref[e:90]{\pageref*{e:90}}\par
\textbf{gestreckte Winkel} (der) -- straight angle\dotfill\hyperref[e:160]{\pageref*{e:160}}\par
\textbf{gleich} -- equal\dotfill\hyperref[e:60]{\pageref*{e:60}}\par
\textbf{gleichnamige Brüche} -- like fractions\dotfill\hyperref[e:34]{\pageref*{e:34}}\par
\textbf{gleichschenklige Dreieck} (das) -- isosceles triangle\dotfill\hyperref[e:182]{\pageref*{e:182}}\par
\textbf{gleichseitige Dreieck} (das) -- equilateral triangle\dotfill\hyperref[e:185]{\pageref*{e:185}}\par
\textbf{Gleichung} (die) -- equation\dotfill\hyperref[e:76]{\pageref*{e:76}}\par
\textbf{Gleichungssystem} (das) -- system of equations\dotfill\hyperref[e:80]{\pageref*{e:80}}\par
\textbf{Grad} (der) -- degree\dotfill\hyperref[e:75]{\pageref*{e:75}}\par
\textbf{Gradmaß} (das) -- degree measure\dotfill\hyperref[e:138]{\pageref*{e:138}}\par
\textbf{Grenzwert} (der) -- limit\dotfill\hyperref[e:264]{\pageref*{e:264}}\par
\textbf{größer als} -- greater than\dotfill\hyperref[e:64]{\pageref*{e:64}}\par
\textbf{größer oder gleich} -- greater than or equal to\dotfill\hyperref[e:66]{\pageref*{e:66}}\par
\textbf{größte gemeinsame Teiler} (der) -- greatest common divisor\dotfill\hyperref[e:56]{\pageref*{e:56}}\par
\textbf{Grundfläche} (die) -- base\dotfill\hyperref[e:248]{\pageref*{e:248}}\par

\medskip{\large\textbf{\color{cornflower}H}}\par\nopagebreak
\textbf{Halbachse} (die) -- semi-axis\dotfill\hyperref[e:317]{\pageref*{e:317}}\par
\textbf{Hauptdiagonale} (die) -- main diagonal\dotfill\hyperref[e:324]{\pageref*{e:324}}\par
\textbf{Hauptnenner} (der) -- lowest common denominator\dotfill\hyperref[e:36]{\pageref*{e:36}}\par
\textbf{Hochpunkt} (der) -- maximum\dotfill\hyperref[e:282]{\pageref*{e:282}}\par
\textbf{Höhe} (die) -- height, altitude\dotfill\hyperref[e:191]{\pageref*{e:191}}\par
\textbf{Hohlzylinder} (der) -- hollow cylinder\dotfill\hyperref[e:237]{\pageref*{e:237}}\par
\textbf{Hyperbel} (die) -- hyperbola\dotfill\hyperref[e:315]{\pageref*{e:315}}\par
\textbf{Hypotenuse} (die) -- hypotenuse\dotfill\hyperref[e:186]{\pageref*{e:186}}\par

\medskip{\large\textbf{\color{cornflower}I}}\par\nopagebreak
\textbf{identisch gleich} -- identically equal\dotfill\hyperref[e:62]{\pageref*{e:62}}\par
\textbf{imaginäre Einheit} (die) -- imaginary unit\dotfill\hyperref[e:11]{\pageref*{e:11}}\par
\textbf{Index} (der) -- subscript, index\dotfill\hyperref[e:4]{\pageref*{e:4}}\par
\textbf{Inkreis} (der) -- incircle, inscribed circle\dotfill\hyperref[e:193]{\pageref*{e:193}}\par
\textbf{Innenwinkel} (der) -- interior angle\dotfill\hyperref[e:178]{\pageref*{e:178}}\par
\textbf{Integral} (das) -- integral\dotfill\hyperref[e:295]{\pageref*{e:295}}\par
\textbf{Integralrechnung} (die) -- integral calculus\dotfill\hyperref[e:286]{\pageref*{e:286}}\par
\textbf{Integralzeichen} (das) -- integral sign\dotfill\hyperref[e:289]{\pageref*{e:289}}\par
\textbf{Integrand} (der) -- integrand\dotfill\hyperref[e:290]{\pageref*{e:290}}\par
\textbf{Integrationsgrenze} (die) -- limit of integration\dotfill\hyperref[e:297]{\pageref*{e:297}}\par
\textbf{Integrationskonstante} (die) -- constant of integration\dotfill\hyperref[e:292]{\pageref*{e:292}}\par
\textbf{Integrationsvariable} (die) -- variable of integration\dotfill\hyperref[e:291]{\pageref*{e:291}}\par
\textbf{integrieren} -- integrate\dotfill\hyperref[e:296]{\pageref*{e:296}}\par
\textbf{Intervall} (das) -- interval\dotfill\hyperref[e:270]{\pageref*{e:270}}\par
\textbf{irrationale Zahl} (die) -- irrational number\dotfill\hyperref[e:9]{\pageref*{e:9}}\par

\medskip{\large\textbf{\color{cornflower}K}}\par\nopagebreak
\textbf{Kante} (die) -- edge\dotfill\hyperref[e:246]{\pageref*{e:246}}\par
\textbf{kartesische Produkt} (das) -- Cartesian product\dotfill\hyperref[e:103]{\pageref*{e:103}}\par
\textbf{Kathete} (die) -- leg\dotfill\hyperref[e:187]{\pageref*{e:187}}\par
\textbf{Kegel} (der) -- cone\dotfill\hyperref[e:231]{\pageref*{e:231}}\par
\textbf{Kegelschnitt} (der) -- conic section\dotfill\hyperref[e:312]{\pageref*{e:312}}\par
\textbf{Kegelstumpf} (der) -- frustum of a cone\dotfill\hyperref[e:236]{\pageref*{e:236}}\par
\textbf{Kehrwert} (der) -- reciprocal\dotfill\hyperref[e:39]{\pageref*{e:39}}\par
\textbf{Klammer} (die) -- bracket\dotfill\hyperref[e:87]{\pageref*{e:87}}\par
\textbf{kleiner als} -- less than\dotfill\hyperref[e:63]{\pageref*{e:63}}\par
\textbf{kleiner oder gleich} -- less than or equal to\dotfill\hyperref[e:65]{\pageref*{e:65}}\par
\textbf{kleinste gemeinsame Vielfache} (das) -- least common multiple\dotfill\hyperref[e:57]{\pageref*{e:57}}\par
\textbf{Koeffizient} (der) -- coefficient\dotfill\hyperref[e:74]{\pageref*{e:74}}\par
\textbf{Kombination} (die) -- combination\dotfill\hyperref[e:343]{\pageref*{e:343}}\par
\textbf{Komplementwinkel} (der) -- complementary angle\dotfill\hyperref[e:164]{\pageref*{e:164}}\par
\textbf{komplexe Zahl} (die) -- complex number\dotfill\hyperref[e:10]{\pageref*{e:10}}\par
\textbf{kongruent} -- congruent\dotfill\hyperref[e:175]{\pageref*{e:175}}\par
\textbf{konvergent} -- convergent\dotfill\hyperref[e:262]{\pageref*{e:262}}\par
\textbf{Koordinaten} (die) -- coordinates\dotfill\hyperref[e:116]{\pageref*{e:116}}\par
\textbf{Koordinatensystem} (das) -- coordinate system\dotfill\hyperref[e:113]{\pageref*{e:113}}\par
\textbf{Körper} (der) -- solid\dotfill\hyperref[e:225]{\pageref*{e:225}}\par
\textbf{Kosinus} (der) -- cosine\dotfill\hyperref[e:133]{\pageref*{e:133}}\par
\textbf{Kosinussatz} (der) -- cosine rule, law of cosines\dotfill\hyperref[e:141]{\pageref*{e:141}}\par
\textbf{Kotangens} (der) -- cotangent\dotfill\hyperref[e:135]{\pageref*{e:135}}\par
\textbf{Kreis} (der) -- circle\dotfill\hyperref[e:211]{\pageref*{e:211}}\par
\textbf{Kreisabschnitt} (der) -- segment (of a circle)\dotfill\hyperref[e:222]{\pageref*{e:222}}\par
\textbf{Kreisbogen} (der) -- arc\dotfill\hyperref[e:221]{\pageref*{e:221}}\par
\textbf{Kreisfläche} (die) -- area of a circle\dotfill\hyperref[e:223]{\pageref*{e:223}}\par
\textbf{Kreisgleichung} (die) -- equation of a circle\dotfill\hyperref[e:313]{\pageref*{e:313}}\par
\textbf{Kreissektor} (der) -- sector\dotfill\hyperref[e:217]{\pageref*{e:217}}\par
\textbf{Kreisumfang} (der) -- circumference\dotfill\hyperref[e:215]{\pageref*{e:215}}\par
\textbf{Kreiszahl} (die) -- pi\dotfill\hyperref[e:216]{\pageref*{e:216}}\par
\textbf{Kreuzprodukt} (das) -- cross product\dotfill\hyperref[e:307]{\pageref*{e:307}}\par
\textbf{Kubikwurzel} (die) -- cube root\dotfill\hyperref[e:46]{\pageref*{e:46}}\par
\textbf{Kugel} (die) -- sphere\dotfill\hyperref[e:232]{\pageref*{e:232}}\par
\textbf{Kugelsegment} (das) -- spherical cap, spherical segment\dotfill\hyperref[e:240]{\pageref*{e:240}}\par
\textbf{Kugelsektor} (der) -- spherical sector\dotfill\hyperref[e:241]{\pageref*{e:241}}\par
\textbf{kürzen} -- cancel, reduce\dotfill\hyperref[e:37]{\pageref*{e:37}}\par

\medskip{\large\textbf{\color{cornflower}L}}\par\nopagebreak
\textbf{leere Menge} (die) -- empty set\dotfill\hyperref[e:104]{\pageref*{e:104}}\par
\textbf{Leitlinie} (die) -- directrix\dotfill\hyperref[e:318]{\pageref*{e:318}}\par
\textbf{linear abhängig} -- linearly dependent\dotfill\hyperref[e:319]{\pageref*{e:319}}\par
\textbf{lineare Funktion} (die) -- linear function\dotfill\hyperref[e:119]{\pageref*{e:119}}\par
\textbf{lineare Gleichung} (die) -- linear equation\dotfill\hyperref[e:81]{\pageref*{e:81}}\par
\textbf{logarithmische Gleichung} (die) -- logarithmic equation\dotfill\hyperref[e:85]{\pageref*{e:85}}\par
\textbf{Logarithmus} (der) -- logarithm\dotfill\hyperref[e:95]{\pageref*{e:95}}\par
\textbf{Logarithmusfunktion} (die) -- logarithmic function\dotfill\hyperref[e:131]{\pageref*{e:131}}\par
\textbf{Lösung} (die) -- solution\dotfill\hyperref[e:78]{\pageref*{e:78}}\par
\textbf{Lösungsmenge} (die) -- solution set\dotfill\hyperref[e:79]{\pageref*{e:79}}\par

\medskip{\large\textbf{\color{cornflower}M}}\par\nopagebreak
\textbf{Mantelfläche} (die) -- lateral surface, curved surface\dotfill\hyperref[e:250]{\pageref*{e:250}}\par
\textbf{Matrix} (die) -- matrix\dotfill\hyperref[e:320]{\pageref*{e:320}}\par
\textbf{Median} (der) -- median\dotfill\hyperref[e:339]{\pageref*{e:339}}\par
\textbf{Menge} (die) -- set\dotfill\hyperref[e:96]{\pageref*{e:96}}\par
\textbf{Minuend} (der) -- minuend\dotfill\hyperref[e:19]{\pageref*{e:19}}\par
\textbf{Mittelpunkt} (der) -- centre\dotfill\hyperref[e:212]{\pageref*{e:212}}\par
\textbf{Mittelsenkrechte} (die) -- perpendicular bisector\dotfill\hyperref[e:170]{\pageref*{e:170}}\par
\textbf{Mittelwert} (der) -- mean\dotfill\hyperref[e:338]{\pageref*{e:338}}\par
\textbf{monoton fallend} -- decreasing\dotfill\hyperref[e:143]{\pageref*{e:143}}\par
\textbf{monoton steigend} -- increasing\dotfill\hyperref[e:142]{\pageref*{e:142}}\par
\textbf{multiplizieren} -- multiply\dotfill\hyperref[e:21]{\pageref*{e:21}}\par

\medskip{\large\textbf{\color{cornflower}N}}\par\nopagebreak
\textbf{$n$-te Wurzel} (die) -- $n$th root\dotfill\hyperref[e:47]{\pageref*{e:47}}\par
\textbf{natürliche Exponentialfunktion} (die) -- natural exponential function\dotfill\hyperref[e:299]{\pageref*{e:299}}\par
\textbf{natürliche Logarithmus} (der) -- natural logarithm\dotfill\hyperref[e:300]{\pageref*{e:300}}\par
\textbf{natürliche Zahl} (die) -- natural number\dotfill\hyperref[e:5]{\pageref*{e:5}}\par
\textbf{Nebenwinkel} (der) -- supplementary angle\dotfill\hyperref[e:163]{\pageref*{e:163}}\par
\textbf{Nenner} (der) -- denominator\dotfill\hyperref[e:30]{\pageref*{e:30}}\par
\textbf{Netz} (das) -- net\dotfill\hyperref[e:253]{\pageref*{e:253}}\par
\textbf{Nullfolge} (die) -- null sequence\dotfill\hyperref[e:258]{\pageref*{e:258}}\par
\textbf{Nullstelle} (die) -- zero, root\dotfill\hyperref[e:122]{\pageref*{e:122}}\par
\textbf{Nullvektor} (der) -- zero vector\dotfill\hyperref[e:304]{\pageref*{e:304}}\par

\medskip{\large\textbf{\color{cornflower}O}}\par\nopagebreak
\textbf{Oberflächeninhalt} (der) -- surface area\dotfill\hyperref[e:256]{\pageref*{e:256}}\par
\textbf{Obersumme} (die) -- upper sum\dotfill\hyperref[e:293]{\pageref*{e:293}}\par
\textbf{Ordinate} (die) -- ordinate\dotfill\hyperref[e:118]{\pageref*{e:118}}\par
\textbf{Ordnungszahl} (die) -- ordinal number\dotfill\hyperref[e:3]{\pageref*{e:3}}\par
\textbf{Ortsvektor} (der) -- position vector\dotfill\hyperref[e:302]{\pageref*{e:302}}\par

\medskip{\large\textbf{\color{cornflower}P}}\par\nopagebreak
\textbf{Parabel} (die) -- parabola\dotfill\hyperref[e:125]{\pageref*{e:125}}\par
\textbf{parallel} -- parallel\dotfill\hyperref[e:150]{\pageref*{e:150}}\par
\textbf{Parallelogramm} (das) -- parallelogram\dotfill\hyperref[e:200]{\pageref*{e:200}}\par
\textbf{partielle Ableitung} (die) -- partial derivative\dotfill\hyperref[e:277]{\pageref*{e:277}}\par
\textbf{Permutation} (die) -- permutation\dotfill\hyperref[e:341]{\pageref*{e:341}}\par
\textbf{Polarkoordinaten} (die) -- polar coordinates\dotfill\hyperref[e:311]{\pageref*{e:311}}\par
\textbf{Polyeder} (das) -- polyhedron\dotfill\hyperref[e:233]{\pageref*{e:233}}\par
\textbf{Potenz} (die) -- power\dotfill\hyperref[e:42]{\pageref*{e:42}}\par
\textbf{potenzieren} -- raise to a power\dotfill\hyperref[e:50]{\pageref*{e:50}}\par
\textbf{Potenzreihe} (die) -- power series\dotfill\hyperref[e:267]{\pageref*{e:267}}\par
\textbf{Primzahl} (die) -- prime number\dotfill\hyperref[e:12]{\pageref*{e:12}}\par
\textbf{Prisma} (das) -- prism\dotfill\hyperref[e:228]{\pageref*{e:228}}\par
\textbf{Produkt} (das) -- product\dotfill\hyperref[e:22]{\pageref*{e:22}}\par
\textbf{proportional zu} -- proportional to\dotfill\hyperref[e:68]{\pageref*{e:68}}\par
\textbf{Proportionalität} (die) -- proportionality\dotfill\hyperref[e:94]{\pageref*{e:94}}\par
\textbf{Prozentsatz} (der) -- percentage\dotfill\hyperref[e:58]{\pageref*{e:58}}\par
\textbf{Punkt} (der) -- point\dotfill\hyperref[e:146]{\pageref*{e:146}}\par
\textbf{punktsymmetrisch zum Ursprung} -- symmetric about the origin\dotfill\hyperref[e:145]{\pageref*{e:145}}\par
\textbf{Pyramide} (die) -- pyramid\dotfill\hyperref[e:230]{\pageref*{e:230}}\par
\textbf{Pyramidenstumpf} (der) -- frustum of a pyramid\dotfill\hyperref[e:235]{\pageref*{e:235}}\par

\medskip{\large\textbf{\color{cornflower}Q}}\par\nopagebreak
\textbf{Quader} (der) -- cuboid\dotfill\hyperref[e:227]{\pageref*{e:227}}\par
\textbf{Quadrat} (das) -- square\dotfill\hyperref[e:198]{\pageref*{e:198}}\par
\textbf{quadratische Funktion} (die) -- quadratic function\dotfill\hyperref[e:124]{\pageref*{e:124}}\par
\textbf{quadratische Gleichung} (die) -- quadratic equation\dotfill\hyperref[e:82]{\pageref*{e:82}}\par
\textbf{quadratische Matrix} (die) -- square matrix\dotfill\hyperref[e:321]{\pageref*{e:321}}\par
\textbf{Quadratwurzel} (die) -- square root\dotfill\hyperref[e:45]{\pageref*{e:45}}\par
\textbf{quadrieren} -- square\dotfill\hyperref[e:51]{\pageref*{e:51}}\par
\textbf{Quotient} (der) -- quotient\dotfill\hyperref[e:25]{\pageref*{e:25}}\par

\medskip{\large\textbf{\color{cornflower}R}}\par\nopagebreak
\textbf{Radikand} (der) -- radicand\dotfill\hyperref[e:48]{\pageref*{e:48}}\par
\textbf{Radius} (der) -- radius\dotfill\hyperref[e:213]{\pageref*{e:213}}\par
\textbf{radizieren, Wurzel ziehen} -- extract a root\dotfill\hyperref[e:52]{\pageref*{e:52}}\par
\textbf{rationale Zahl} (die) -- rational number\dotfill\hyperref[e:7]{\pageref*{e:7}}\par
\textbf{Raumdiagonale} (die) -- space diagonal\dotfill\hyperref[e:252]{\pageref*{e:252}}\par
\textbf{Raute} (die) -- rhombus\dotfill\hyperref[e:201]{\pageref*{e:201}}\par
\textbf{rechte Winkel} (der) -- right angle\dotfill\hyperref[e:158]{\pageref*{e:158}}\par
\textbf{Rechteck} (das) -- rectangle\dotfill\hyperref[e:199]{\pageref*{e:199}}\par
\textbf{rechtwinklige Dreieck} (das) -- right-angled triangle\dotfill\hyperref[e:181]{\pageref*{e:181}}\par
\textbf{reelle Zahl} (die) -- real number\dotfill\hyperref[e:8]{\pageref*{e:8}}\par
\textbf{regelmäßige Sechseck} (das) -- regular hexagon\dotfill\hyperref[e:205]{\pageref*{e:205}}\par
\textbf{Reihe} (die) -- series\dotfill\hyperref[e:265]{\pageref*{e:265}}\par
\textbf{relative Häufigkeit} (die) -- relative frequency\dotfill\hyperref[e:331]{\pageref*{e:331}}\par
\textbf{Rotationsellipsoid} (das) -- spheroid\dotfill\hyperref[e:242]{\pageref*{e:242}}\par
\textbf{Rotationshyperboloid} (das) -- hyperboloid of revolution\dotfill\hyperref[e:244]{\pageref*{e:244}}\par
\textbf{Rotationskörper} (der) -- solid of revolution\dotfill\hyperref[e:239]{\pageref*{e:239}}\par
\textbf{Rotationsparaboloid} (das) -- paraboloid of revolution\dotfill\hyperref[e:243]{\pageref*{e:243}}\par
\textbf{runden} -- round\dotfill\hyperref[e:59]{\pageref*{e:59}}\par
\textbf{runden Klammern} (die) -- round brackets, parentheses\dotfill\hyperref[e:88]{\pageref*{e:88}}\par

\medskip{\large\textbf{\color{cornflower}S}}\par\nopagebreak
\textbf{Satz} (der) -- theorem\dotfill\hyperref[e:347]{\pageref*{e:347}}\par
\textbf{Satz des Pythagoras} (der) -- Pythagorean theorem\dotfill\hyperref[e:190]{\pageref*{e:190}}\par
\textbf{Scheitel} (der) -- vertex\dotfill\hyperref[e:154]{\pageref*{e:154}}\par
\textbf{Scheitelpunkt} (der) -- vertex\dotfill\hyperref[e:129]{\pageref*{e:129}}\par
\textbf{Scheitelwinkel} (der) -- vertically opposite angle\dotfill\hyperref[e:166]{\pageref*{e:166}}\par
\textbf{Schenkel} (der) -- arm, leg\dotfill\hyperref[e:155]{\pageref*{e:155}}\par
\textbf{schiefe Zylinder} (der) -- oblique cylinder\dotfill\hyperref[e:238]{\pageref*{e:238}}\par
\textbf{Schnittmenge} (die) -- intersection\dotfill\hyperref[e:100]{\pageref*{e:100}}\par
\textbf{Schnittpunkt} (der) -- point of intersection\dotfill\hyperref[e:123]{\pageref*{e:123}}\par
\textbf{Schrägbild} (das) -- oblique drawing\dotfill\hyperref[e:254]{\pageref*{e:254}}\par
\textbf{Sehne} (die) -- chord\dotfill\hyperref[e:218]{\pageref*{e:218}}\par
\textbf{Seitenfläche} (die) -- face\dotfill\hyperref[e:247]{\pageref*{e:247}}\par
\textbf{Seitenhalbierende} (die) -- median\dotfill\hyperref[e:192]{\pageref*{e:192}}\par
\textbf{Sekante} (die) -- secant\dotfill\hyperref[e:219]{\pageref*{e:219}}\par
\textbf{senkrecht} -- perpendicular\dotfill\hyperref[e:151]{\pageref*{e:151}}\par
\textbf{Sinus} (der) -- sine\dotfill\hyperref[e:132]{\pageref*{e:132}}\par
\textbf{Sinussatz} (der) -- sine rule, law of sines\dotfill\hyperref[e:140]{\pageref*{e:140}}\par
\textbf{Skalarprodukt} (das) -- dot product, scalar product\dotfill\hyperref[e:306]{\pageref*{e:306}}\par
\textbf{Spalte} (die) -- column\dotfill\hyperref[e:323]{\pageref*{e:323}}\par
\textbf{Spiegelung} (die) -- reflection\dotfill\hyperref[e:172]{\pageref*{e:172}}\par
\textbf{spitze Winkel} (der) -- acute angle\dotfill\hyperref[e:157]{\pageref*{e:157}}\par
\textbf{Stammfunktion} (die) -- antiderivative\dotfill\hyperref[e:285]{\pageref*{e:285}}\par
\textbf{Standardabweichung} (die) -- standard deviation\dotfill\hyperref[e:337]{\pageref*{e:337}}\par
\textbf{Steigung} (die) -- gradient\dotfill\hyperref[e:120]{\pageref*{e:120}}\par
\textbf{Steigungswinkel} (der) -- angle of inclination\dotfill\hyperref[e:310]{\pageref*{e:310}}\par
\textbf{stetig} -- continuous\dotfill\hyperref[e:268]{\pageref*{e:268}}\par
\textbf{Strahl} (der) -- ray\dotfill\hyperref[e:149]{\pageref*{e:149}}\par
\textbf{Strecke} (die) -- line segment\dotfill\hyperref[e:148]{\pageref*{e:148}}\par
\textbf{Stufenwinkel} (der) -- corresponding angle\dotfill\hyperref[e:167]{\pageref*{e:167}}\par
\textbf{stumpfe Winkel} (der) -- obtuse angle\dotfill\hyperref[e:159]{\pageref*{e:159}}\par
\textbf{Subtrahend} (der) -- subtrahend\dotfill\hyperref[e:20]{\pageref*{e:20}}\par
\textbf{subtrahieren} -- subtract\dotfill\hyperref[e:17]{\pageref*{e:17}}\par
\textbf{Summand} (der) -- summand, addend\dotfill\hyperref[e:16]{\pageref*{e:16}}\par
\textbf{Summe} (die) -- sum\dotfill\hyperref[e:15]{\pageref*{e:15}}\par
\textbf{Supplementwinkel} (der) -- supplementary angle\dotfill\hyperref[e:165]{\pageref*{e:165}}\par
\textbf{Symmetrieachse} (die) -- axis of symmetry\dotfill\hyperref[e:171]{\pageref*{e:171}}\par

\medskip{\large\textbf{\color{cornflower}T}}\par\nopagebreak
\textbf{Tangens} (der) -- tangent\dotfill\hyperref[e:134]{\pageref*{e:134}}\par
\textbf{Tangente} (die) -- tangent\dotfill\hyperref[e:224]{\pageref*{e:224}}\par
\textbf{Tangente} (die) -- tangent\dotfill\hyperref[e:279]{\pageref*{e:279}}\par
\textbf{Teiler} (der) -- divisor\dotfill\hyperref[e:54]{\pageref*{e:54}}\par
\textbf{Teilmenge} (die) -- subset\dotfill\hyperref[e:98]{\pageref*{e:98}}\par
\textbf{Term} (der) -- expression\dotfill\hyperref[e:72]{\pageref*{e:72}}\par
\textbf{Tetraeder} (das) -- tetrahedron\dotfill\hyperref[e:234]{\pageref*{e:234}}\par
\textbf{Tiefpunkt} (der) -- minimum\dotfill\hyperref[e:283]{\pageref*{e:283}}\par
\textbf{Trapez} (das) -- trapezium\dotfill\hyperref[e:202]{\pageref*{e:202}}\par
\textbf{trigonometrische Gleichung} (die) -- trigonometric equation\dotfill\hyperref[e:86]{\pageref*{e:86}}\par

\medskip{\large\textbf{\color{cornflower}U}}\par\nopagebreak
\textbf{überstumpfe Winkel} (der) -- reflex angle\dotfill\hyperref[e:161]{\pageref*{e:161}}\par
\textbf{Umfang} (der) -- perimeter\dotfill\hyperref[e:209]{\pageref*{e:209}}\par
\textbf{Umkreis} (der) -- circumcircle\dotfill\hyperref[e:194]{\pageref*{e:194}}\par
\textbf{unabhängige Variable} (die) -- independent variable\dotfill\hyperref[e:106]{\pageref*{e:106}}\par
\textbf{Unbekannte} (die) -- unknown\dotfill\hyperref[e:73]{\pageref*{e:73}}\par
\textbf{unbestimmte Integral} (das) -- indefinite integral\dotfill\hyperref[e:287]{\pageref*{e:287}}\par
\textbf{unechte Bruch} (der) -- improper fraction\dotfill\hyperref[e:32]{\pageref*{e:32}}\par
\textbf{unendlich} -- infinite\dotfill\hyperref[e:70]{\pageref*{e:70}}\par
\textbf{unendliche Reihe} (die) -- infinite series\dotfill\hyperref[e:266]{\pageref*{e:266}}\par
\textbf{ungefähr} -- approximately\dotfill\hyperref[e:69]{\pageref*{e:69}}\par
\textbf{ungleich} -- not equal\dotfill\hyperref[e:61]{\pageref*{e:61}}\par
\textbf{ungleichnamige Brüche} -- unlike fractions\dotfill\hyperref[e:35]{\pageref*{e:35}}\par
\textbf{Ungleichung} (die) -- inequality\dotfill\hyperref[e:77]{\pageref*{e:77}}\par
\textbf{unregelmäßige Dreieck} (das) -- scalene triangle\dotfill\hyperref[e:180]{\pageref*{e:180}}\par
\textbf{unstetig} -- discontinuous\dotfill\hyperref[e:269]{\pageref*{e:269}}\par
\textbf{Untersumme} (die) -- lower sum\dotfill\hyperref[e:294]{\pageref*{e:294}}\par
\textbf{Ursprung} (der) -- origin\dotfill\hyperref[e:115]{\pageref*{e:115}}\par

\medskip{\large\textbf{\color{cornflower}V}}\par\nopagebreak
\textbf{Variable} (die) -- variable\dotfill\hyperref[e:71]{\pageref*{e:71}}\par
\textbf{Varianz} (die) -- variance\dotfill\hyperref[e:336]{\pageref*{e:336}}\par
\textbf{Variation} (die) -- variation, arrangement\dotfill\hyperref[e:342]{\pageref*{e:342}}\par
\textbf{Vektor} (der) -- vector\dotfill\hyperref[e:301]{\pageref*{e:301}}\par
\textbf{Vereinigungsmenge} (die) -- union\dotfill\hyperref[e:101]{\pageref*{e:101}}\par
\textbf{Verhältnis} (das) -- ratio\dotfill\hyperref[e:40]{\pageref*{e:40}}\par
\textbf{Verhältnisgleichung} (die) -- proportion\dotfill\hyperref[e:41]{\pageref*{e:41}}\par
\textbf{Verschiebung} (die) -- translation\dotfill\hyperref[e:173]{\pageref*{e:173}}\par
\textbf{viel kleiner als} -- much less than\dotfill\hyperref[e:67]{\pageref*{e:67}}\par
\textbf{Vieleck} (das) -- polygon\dotfill\hyperref[e:204]{\pageref*{e:204}}\par
\textbf{Vielfache} (das) -- multiple\dotfill\hyperref[e:55]{\pageref*{e:55}}\par
\textbf{Viereck} (das) -- quadrilateral\dotfill\hyperref[e:196]{\pageref*{e:196}}\par
\textbf{Vierfeldertafel} (die) -- two-way table\dotfill\hyperref[e:333]{\pageref*{e:333}}\par
\textbf{Vollwinkel} (der) -- full angle\dotfill\hyperref[e:162]{\pageref*{e:162}}\par
\textbf{Volumen} (das) -- volume\dotfill\hyperref[e:255]{\pageref*{e:255}}\par
\textbf{Voraussetzung} (die) -- assumption, hypothesis\dotfill\hyperref[e:348]{\pageref*{e:348}}\par

\medskip{\large\textbf{\color{cornflower}W}}\par\nopagebreak
\textbf{Wahrscheinlichkeit} (die) -- probability\dotfill\hyperref[e:330]{\pageref*{e:330}}\par
\textbf{Wechselwinkel} (der) -- alternate angle\dotfill\hyperref[e:168]{\pageref*{e:168}}\par
\textbf{Wendepunkt} (der) -- point of inflection\dotfill\hyperref[e:284]{\pageref*{e:284}}\par
\textbf{Wertebereich} (der) -- range\dotfill\hyperref[e:110]{\pageref*{e:110}}\par
\textbf{Wertetabelle} (die) -- table of values\dotfill\hyperref[e:112]{\pageref*{e:112}}\par
\textbf{Winkel} (der) -- angle\dotfill\hyperref[e:153]{\pageref*{e:153}}\par
\textbf{Winkelhalbierende} (die) -- angle bisector\dotfill\hyperref[e:156]{\pageref*{e:156}}\par
\textbf{Würfel} (der) -- cube\dotfill\hyperref[e:226]{\pageref*{e:226}}\par
\textbf{Wurzelexponent} (der) -- index (of a root)\dotfill\hyperref[e:49]{\pageref*{e:49}}\par
\textbf{Wurzelfunktion} (die) -- square root function\dotfill\hyperref[e:128]{\pageref*{e:128}}\par
\textbf{Wurzelgleichung} (die) -- radical equation\dotfill\hyperref[e:83]{\pageref*{e:83}}\par

\medskip{\large\textbf{\color{cornflower}X}}\par\nopagebreak
\textbf{$x$-Achse} (die) -- $x$-axis\dotfill\hyperref[e:114]{\pageref*{e:114}}\par

\medskip{\large\textbf{\color{cornflower}Y}}\par\nopagebreak
\textbf{$y$-Achsenabschnitt} (der) -- $y$-intercept\dotfill\hyperref[e:121]{\pageref*{e:121}}\par

\medskip{\large\textbf{\color{cornflower}Z}}\par\nopagebreak
\textbf{Zahl} (die) -- number\dotfill\hyperref[e:1]{\pageref*{e:1}}\par
\textbf{Zähler} (der) -- numerator\dotfill\hyperref[e:29]{\pageref*{e:29}}\par
\textbf{Zeile} (die) -- row\dotfill\hyperref[e:322]{\pageref*{e:322}}\par
\textbf{Ziffer} (die) -- digit\dotfill\hyperref[e:2]{\pageref*{e:2}}\par
\textbf{Zufallsexperiment} (das) -- random experiment\dotfill\hyperref[e:326]{\pageref*{e:326}}\par
\textbf{Zufallsgröße} (die) -- random variable\dotfill\hyperref[e:334]{\pageref*{e:334}}\par
\textbf{Zwölfeck} (das) -- dodecagon\dotfill\hyperref[e:208]{\pageref*{e:208}}\par
\textbf{Zylinder} (der) -- cylinder\dotfill\hyperref[e:229]{\pageref*{e:229}}\par
\end{multicols}

\newpage
\section{Index English -- German}
\begin{multicols}{2}\small\raggedright

\medskip{\large\textbf{\color{cornflower}A}}\par\nopagebreak
\textbf{abscissa} -- die Abszisse\dotfill\hyperref[e:117]{\pageref*{e:117}}\par
\textbf{absolute value} -- der Betrag\dotfill\hyperref[e:53]{\pageref*{e:53}}\par
\textbf{acute angle} -- der spitze Winkel\dotfill\hyperref[e:157]{\pageref*{e:157}}\par
\textbf{add} -- addieren\dotfill\hyperref[e:14]{\pageref*{e:14}}\par
\textbf{adjacent side} -- die Ankathete\dotfill\hyperref[e:188]{\pageref*{e:188}}\par
\textbf{alternate angle} -- der Wechselwinkel\dotfill\hyperref[e:168]{\pageref*{e:168}}\par
\textbf{angle} -- der Winkel\dotfill\hyperref[e:153]{\pageref*{e:153}}\par
\textbf{angle bisector} -- die Winkelhalbierende\dotfill\hyperref[e:156]{\pageref*{e:156}}\par
\textbf{angle of inclination} -- der Steigungswinkel\dotfill\hyperref[e:310]{\pageref*{e:310}}\par
\textbf{antiderivative} -- die Stammfunktion\dotfill\hyperref[e:285]{\pageref*{e:285}}\par
\textbf{approximately} -- ungefähr\dotfill\hyperref[e:69]{\pageref*{e:69}}\par
\textbf{arc} -- der Kreisbogen\dotfill\hyperref[e:221]{\pageref*{e:221}}\par
\textbf{area} -- der Flächeninhalt\dotfill\hyperref[e:210]{\pageref*{e:210}}\par
\textbf{area of a circle} -- die Kreisfläche\dotfill\hyperref[e:223]{\pageref*{e:223}}\par
\textbf{arithmetic sequence} -- die arithmetische Folge\dotfill\hyperref[e:259]{\pageref*{e:259}}\par
\textbf{arm, leg} -- der Schenkel\dotfill\hyperref[e:155]{\pageref*{e:155}}\par
\textbf{assumption, hypothesis} -- die Voraussetzung\dotfill\hyperref[e:348]{\pageref*{e:348}}\par
\textbf{asymptote} -- die Asymptote\dotfill\hyperref[e:271]{\pageref*{e:271}}\par
\textbf{axis of symmetry} -- die Symmetrieachse\dotfill\hyperref[e:171]{\pageref*{e:171}}\par

\medskip{\large\textbf{\color{cornflower}B}}\par\nopagebreak
\textbf{base} -- die Basis\dotfill\hyperref[e:43]{\pageref*{e:43}}\par
\textbf{base} -- die Basis (im Dreieck)\dotfill\hyperref[e:183]{\pageref*{e:183}}\par
\textbf{base} -- die Grundfläche\dotfill\hyperref[e:248]{\pageref*{e:248}}\par
\textbf{base angle} -- der Basiswinkel\dotfill\hyperref[e:184]{\pageref*{e:184}}\par
\textbf{binomial coefficient} -- der Binomialkoeffizient\dotfill\hyperref[e:344]{\pageref*{e:344}}\par
\textbf{binomial distribution} -- die Binomialverteilung\dotfill\hyperref[e:345]{\pageref*{e:345}}\par
\textbf{binomial formula} -- die binomische Formel\dotfill\hyperref[e:93]{\pageref*{e:93}}\par
\textbf{bracket} -- die Klammer\dotfill\hyperref[e:87]{\pageref*{e:87}}\par

\medskip{\large\textbf{\color{cornflower}C}}\par\nopagebreak
\textbf{cancel, reduce} -- kürzen\dotfill\hyperref[e:37]{\pageref*{e:37}}\par
\textbf{Cartesian product} -- das kartesische Produkt\dotfill\hyperref[e:103]{\pageref*{e:103}}\par
\textbf{centre} -- der Mittelpunkt\dotfill\hyperref[e:212]{\pageref*{e:212}}\par
\textbf{chord} -- die Sehne\dotfill\hyperref[e:218]{\pageref*{e:218}}\par
\textbf{circle} -- der Kreis\dotfill\hyperref[e:211]{\pageref*{e:211}}\par
\textbf{circumcircle} -- der Umkreis\dotfill\hyperref[e:194]{\pageref*{e:194}}\par
\textbf{circumference} -- der Kreisumfang\dotfill\hyperref[e:215]{\pageref*{e:215}}\par
\textbf{co-interior angles, allied angles} -- entgegengesetzt liegende Winkel\dotfill\hyperref[e:169]{\pageref*{e:169}}\par
\textbf{coefficient} -- der Koeffizient\dotfill\hyperref[e:74]{\pageref*{e:74}}\par
\textbf{column} -- die Spalte\dotfill\hyperref[e:323]{\pageref*{e:323}}\par
\textbf{combination} -- die Kombination\dotfill\hyperref[e:343]{\pageref*{e:343}}\par
\textbf{complementary angle} -- der Komplementwinkel\dotfill\hyperref[e:164]{\pageref*{e:164}}\par
\textbf{complementary event} -- das Gegenereignis\dotfill\hyperref[e:329]{\pageref*{e:329}}\par
\textbf{complex number} -- die komplexe Zahl\dotfill\hyperref[e:10]{\pageref*{e:10}}\par
\textbf{cone} -- der Kegel\dotfill\hyperref[e:231]{\pageref*{e:231}}\par
\textbf{congruent} -- kongruent\dotfill\hyperref[e:175]{\pageref*{e:175}}\par
\textbf{conic section} -- der Kegelschnitt\dotfill\hyperref[e:312]{\pageref*{e:312}}\par
\textbf{constant of integration} -- die Integrationskonstante\dotfill\hyperref[e:292]{\pageref*{e:292}}\par
\textbf{continuous} -- stetig\dotfill\hyperref[e:268]{\pageref*{e:268}}\par
\textbf{convergent} -- konvergent\dotfill\hyperref[e:262]{\pageref*{e:262}}\par
\textbf{coordinate system} -- das Koordinatensystem\dotfill\hyperref[e:113]{\pageref*{e:113}}\par
\textbf{coordinates} -- die Koordinaten\dotfill\hyperref[e:116]{\pageref*{e:116}}\par
\textbf{corresponding angle} -- der Stufenwinkel\dotfill\hyperref[e:167]{\pageref*{e:167}}\par
\textbf{cosine} -- der Kosinus\dotfill\hyperref[e:133]{\pageref*{e:133}}\par
\textbf{cosine rule, law of cosines} -- der Kosinussatz\dotfill\hyperref[e:141]{\pageref*{e:141}}\par
\textbf{cotangent} -- der Kotangens\dotfill\hyperref[e:135]{\pageref*{e:135}}\par
\textbf{cross product} -- das Kreuzprodukt\dotfill\hyperref[e:307]{\pageref*{e:307}}\par
\textbf{cube} -- der Würfel\dotfill\hyperref[e:226]{\pageref*{e:226}}\par
\textbf{cube root} -- die Kubikwurzel\dotfill\hyperref[e:46]{\pageref*{e:46}}\par
\textbf{cuboid} -- der Quader\dotfill\hyperref[e:227]{\pageref*{e:227}}\par
\textbf{curly brackets, braces} -- die geschweiften Klammern\dotfill\hyperref[e:90]{\pageref*{e:90}}\par
\textbf{curve sketching} -- die Funktionsuntersuchung\dotfill\hyperref[e:280]{\pageref*{e:280}}\par
\textbf{cylinder} -- der Zylinder\dotfill\hyperref[e:229]{\pageref*{e:229}}\par

\medskip{\large\textbf{\color{cornflower}D}}\par\nopagebreak
\textbf{decimal number} -- die Dezimalzahl\dotfill\hyperref[e:13]{\pageref*{e:13}}\par
\textbf{decreasing} -- monoton fallend\dotfill\hyperref[e:143]{\pageref*{e:143}}\par
\textbf{definite integral} -- das bestimmte Integral\dotfill\hyperref[e:288]{\pageref*{e:288}}\par
\textbf{definition} -- die Definition\dotfill\hyperref[e:346]{\pageref*{e:346}}\par
\textbf{degree} -- der Grad\dotfill\hyperref[e:75]{\pageref*{e:75}}\par
\textbf{degree measure} -- das Gradmaß\dotfill\hyperref[e:138]{\pageref*{e:138}}\par
\textbf{denominator} -- der Nenner\dotfill\hyperref[e:30]{\pageref*{e:30}}\par
\textbf{dependent variable} -- die abhängige Variable\dotfill\hyperref[e:107]{\pageref*{e:107}}\par
\textbf{derivative} -- die Ableitung\dotfill\hyperref[e:273]{\pageref*{e:273}}\par
\textbf{determinant} -- die Determinante\dotfill\hyperref[e:325]{\pageref*{e:325}}\par
\textbf{diagonal} -- die Diagonale\dotfill\hyperref[e:197]{\pageref*{e:197}}\par
\textbf{diameter} -- der Durchmesser\dotfill\hyperref[e:214]{\pageref*{e:214}}\par
\textbf{difference} -- die Differenz\dotfill\hyperref[e:18]{\pageref*{e:18}}\par
\textbf{difference quotient} -- der Differenzenquotient\dotfill\hyperref[e:272]{\pageref*{e:272}}\par
\textbf{differential} -- das Differenzial\dotfill\hyperref[e:276]{\pageref*{e:276}}\par
\textbf{differential calculus} -- die Differenzialrechnung\dotfill\hyperref[e:274]{\pageref*{e:274}}\par
\textbf{differential quotient} -- der Differenzialquotient\dotfill\hyperref[e:275]{\pageref*{e:275}}\par
\textbf{differentiate} -- ableiten\dotfill\hyperref[e:278]{\pageref*{e:278}}\par
\textbf{digit} -- die Ziffer\dotfill\hyperref[e:2]{\pageref*{e:2}}\par
\textbf{directrix} -- die Leitlinie\dotfill\hyperref[e:318]{\pageref*{e:318}}\par
\textbf{discontinuous} -- unstetig\dotfill\hyperref[e:269]{\pageref*{e:269}}\par
\textbf{distance} -- der Abstand\dotfill\hyperref[e:152]{\pageref*{e:152}}\par
\textbf{divergent} -- divergent\dotfill\hyperref[e:263]{\pageref*{e:263}}\par
\textbf{divide} -- dividieren\dotfill\hyperref[e:24]{\pageref*{e:24}}\par
\textbf{dividend} -- der Dividend\dotfill\hyperref[e:26]{\pageref*{e:26}}\par
\textbf{divisor} -- der Divisor\dotfill\hyperref[e:27]{\pageref*{e:27}}\par
\textbf{divisor} -- der Teiler\dotfill\hyperref[e:54]{\pageref*{e:54}}\par
\textbf{dodecagon} -- das Zwölfeck\dotfill\hyperref[e:208]{\pageref*{e:208}}\par
\textbf{domain} -- der Definitionsbereich\dotfill\hyperref[e:109]{\pageref*{e:109}}\par
\textbf{dot product, scalar product} -- das Skalarprodukt\dotfill\hyperref[e:306]{\pageref*{e:306}}\par

\medskip{\large\textbf{\color{cornflower}E}}\par\nopagebreak
\textbf{edge} -- die Kante\dotfill\hyperref[e:246]{\pageref*{e:246}}\par
\textbf{element} -- das Element\dotfill\hyperref[e:97]{\pageref*{e:97}}\par
\textbf{ellipse} -- die Ellipse\dotfill\hyperref[e:314]{\pageref*{e:314}}\par
\textbf{empty set} -- die leere Menge\dotfill\hyperref[e:104]{\pageref*{e:104}}\par
\textbf{equal} -- gleich\dotfill\hyperref[e:60]{\pageref*{e:60}}\par
\textbf{equal in area} -- flächengleich\dotfill\hyperref[e:195]{\pageref*{e:195}}\par
\textbf{equation} -- die Gleichung\dotfill\hyperref[e:76]{\pageref*{e:76}}\par
\textbf{equation of a circle} -- die Kreisgleichung\dotfill\hyperref[e:313]{\pageref*{e:313}}\par
\textbf{equation of a straight line} -- die Geradengleichung\dotfill\hyperref[e:309]{\pageref*{e:309}}\par
\textbf{equilateral triangle} -- das gleichseitige Dreieck\dotfill\hyperref[e:185]{\pageref*{e:185}}\par
\textbf{Euler's number} -- die eulersche Zahl\dotfill\hyperref[e:298]{\pageref*{e:298}}\par
\textbf{event} -- das Ereignis\dotfill\hyperref[e:328]{\pageref*{e:328}}\par
\textbf{expand} -- ausmultiplizieren\dotfill\hyperref[e:91]{\pageref*{e:91}}\par
\textbf{expand (a fraction)} -- erweitern\dotfill\hyperref[e:38]{\pageref*{e:38}}\par
\textbf{expected value} -- der Erwartungswert\dotfill\hyperref[e:335]{\pageref*{e:335}}\par
\textbf{exponent} -- der Exponent\dotfill\hyperref[e:44]{\pageref*{e:44}}\par
\textbf{exponential equation} -- die Exponentialgleichung\dotfill\hyperref[e:84]{\pageref*{e:84}}\par
\textbf{exponential function} -- die Exponentialfunktion\dotfill\hyperref[e:130]{\pageref*{e:130}}\par
\textbf{expression} -- der Term\dotfill\hyperref[e:72]{\pageref*{e:72}}\par
\textbf{exterior angle} -- der Außenwinkel\dotfill\hyperref[e:179]{\pageref*{e:179}}\par
\textbf{extract a root} -- radizieren, Wurzel ziehen\dotfill\hyperref[e:52]{\pageref*{e:52}}\par
\textbf{extreme point} -- der Extrempunkt\dotfill\hyperref[e:281]{\pageref*{e:281}}\par

\medskip{\large\textbf{\color{cornflower}F}}\par\nopagebreak
\textbf{face} -- die Seitenfläche\dotfill\hyperref[e:247]{\pageref*{e:247}}\par
\textbf{face diagonal} -- die Flächendiagonale\dotfill\hyperref[e:251]{\pageref*{e:251}}\par
\textbf{factor} -- der Faktor\dotfill\hyperref[e:23]{\pageref*{e:23}}\par
\textbf{factorial} -- die Fakultät\dotfill\hyperref[e:340]{\pageref*{e:340}}\par
\textbf{factorise} -- ausklammern\dotfill\hyperref[e:92]{\pageref*{e:92}}\par
\textbf{first term} -- das Anfangsglied\dotfill\hyperref[e:261]{\pageref*{e:261}}\par
\textbf{focus} -- der Brennpunkt\dotfill\hyperref[e:316]{\pageref*{e:316}}\par
\textbf{for all} -- für alle\dotfill\hyperref[e:352]{\pageref*{e:352}}\par
\textbf{fraction} -- der Bruch\dotfill\hyperref[e:28]{\pageref*{e:28}}\par
\textbf{frustum of a cone} -- der Kegelstumpf\dotfill\hyperref[e:236]{\pageref*{e:236}}\par
\textbf{frustum of a pyramid} -- der Pyramidenstumpf\dotfill\hyperref[e:235]{\pageref*{e:235}}\par
\textbf{full angle} -- der Vollwinkel\dotfill\hyperref[e:162]{\pageref*{e:162}}\par
\textbf{function} -- die Funktion\dotfill\hyperref[e:105]{\pageref*{e:105}}\par

\medskip{\large\textbf{\color{cornflower}G}}\par\nopagebreak
\textbf{geometric sequence} -- die geometrische Folge\dotfill\hyperref[e:260]{\pageref*{e:260}}\par
\textbf{gradient} -- die Steigung\dotfill\hyperref[e:120]{\pageref*{e:120}}\par
\textbf{graph} -- der Funktionsgraph\dotfill\hyperref[e:111]{\pageref*{e:111}}\par
\textbf{greater than} -- größer als\dotfill\hyperref[e:64]{\pageref*{e:64}}\par
\textbf{greater than or equal to} -- größer oder gleich\dotfill\hyperref[e:66]{\pageref*{e:66}}\par
\textbf{greatest common divisor} -- der größte gemeinsame Teiler\dotfill\hyperref[e:56]{\pageref*{e:56}}\par

\medskip{\large\textbf{\color{cornflower}H}}\par\nopagebreak
\textbf{height, altitude} -- die Höhe\dotfill\hyperref[e:191]{\pageref*{e:191}}\par
\textbf{hollow cylinder} -- der Hohlzylinder\dotfill\hyperref[e:237]{\pageref*{e:237}}\par
\textbf{hyperbola} -- die Hyperbel\dotfill\hyperref[e:315]{\pageref*{e:315}}\par
\textbf{hyperboloid of revolution} -- das Rotationshyperboloid\dotfill\hyperref[e:244]{\pageref*{e:244}}\par
\textbf{hypotenuse} -- die Hypotenuse\dotfill\hyperref[e:186]{\pageref*{e:186}}\par

\medskip{\large\textbf{\color{cornflower}I}}\par\nopagebreak
\textbf{identically equal} -- identisch gleich\dotfill\hyperref[e:62]{\pageref*{e:62}}\par
\textbf{if and only if} -- genau dann, wenn\dotfill\hyperref[e:351]{\pageref*{e:351}}\par
\textbf{imaginary unit} -- die imaginäre Einheit\dotfill\hyperref[e:11]{\pageref*{e:11}}\par
\textbf{implies} -- daraus folgt\dotfill\hyperref[e:350]{\pageref*{e:350}}\par
\textbf{improper fraction} -- der unechte Bruch\dotfill\hyperref[e:32]{\pageref*{e:32}}\par
\textbf{incircle, inscribed circle} -- der Inkreis\dotfill\hyperref[e:193]{\pageref*{e:193}}\par
\textbf{increasing} -- monoton steigend\dotfill\hyperref[e:142]{\pageref*{e:142}}\par
\textbf{indefinite integral} -- das unbestimmte Integral\dotfill\hyperref[e:287]{\pageref*{e:287}}\par
\textbf{independent variable} -- die unabhängige Variable\dotfill\hyperref[e:106]{\pageref*{e:106}}\par
\textbf{index (of a root)} -- der Wurzelexponent\dotfill\hyperref[e:49]{\pageref*{e:49}}\par
\textbf{inequality} -- die Ungleichung\dotfill\hyperref[e:77]{\pageref*{e:77}}\par
\textbf{infinite} -- unendlich\dotfill\hyperref[e:70]{\pageref*{e:70}}\par
\textbf{infinite series} -- die unendliche Reihe\dotfill\hyperref[e:266]{\pageref*{e:266}}\par
\textbf{integer} -- die ganze Zahl\dotfill\hyperref[e:6]{\pageref*{e:6}}\par
\textbf{integral} -- das Integral\dotfill\hyperref[e:295]{\pageref*{e:295}}\par
\textbf{integral calculus} -- die Integralrechnung\dotfill\hyperref[e:286]{\pageref*{e:286}}\par
\textbf{integral sign} -- das Integralzeichen\dotfill\hyperref[e:289]{\pageref*{e:289}}\par
\textbf{integrand} -- der Integrand\dotfill\hyperref[e:290]{\pageref*{e:290}}\par
\textbf{integrate} -- integrieren\dotfill\hyperref[e:296]{\pageref*{e:296}}\par
\textbf{interior angle} -- der Innenwinkel\dotfill\hyperref[e:178]{\pageref*{e:178}}\par
\textbf{intersection} -- die Schnittmenge\dotfill\hyperref[e:100]{\pageref*{e:100}}\par
\textbf{interval} -- das Intervall\dotfill\hyperref[e:270]{\pageref*{e:270}}\par
\textbf{inverse trigonometric function} -- die Arkusfunktion\dotfill\hyperref[e:136]{\pageref*{e:136}}\par
\textbf{irrational number} -- die irrationale Zahl\dotfill\hyperref[e:9]{\pageref*{e:9}}\par
\textbf{isosceles triangle} -- das gleichschenklige Dreieck\dotfill\hyperref[e:182]{\pageref*{e:182}}\par

\medskip{\large\textbf{\color{cornflower}K}}\par\nopagebreak
\textbf{kite} -- das Drachenviereck\dotfill\hyperref[e:203]{\pageref*{e:203}}\par

\medskip{\large\textbf{\color{cornflower}L}}\par\nopagebreak
\textbf{lateral surface, curved surface} -- die Mantelfläche\dotfill\hyperref[e:250]{\pageref*{e:250}}\par
\textbf{least common multiple} -- das kleinste gemeinsame Vielfache\dotfill\hyperref[e:57]{\pageref*{e:57}}\par
\textbf{leg} -- die Kathete\dotfill\hyperref[e:187]{\pageref*{e:187}}\par
\textbf{less than} -- kleiner als\dotfill\hyperref[e:63]{\pageref*{e:63}}\par
\textbf{less than or equal to} -- kleiner oder gleich\dotfill\hyperref[e:65]{\pageref*{e:65}}\par
\textbf{like fractions} -- gleichnamige Brüche\dotfill\hyperref[e:34]{\pageref*{e:34}}\par
\textbf{limit} -- der Grenzwert\dotfill\hyperref[e:264]{\pageref*{e:264}}\par
\textbf{limit of integration} -- die Integrationsgrenze\dotfill\hyperref[e:297]{\pageref*{e:297}}\par
\textbf{line segment} -- die Strecke\dotfill\hyperref[e:148]{\pageref*{e:148}}\par
\textbf{linear equation} -- die lineare Gleichung\dotfill\hyperref[e:81]{\pageref*{e:81}}\par
\textbf{linear function} -- die lineare Funktion\dotfill\hyperref[e:119]{\pageref*{e:119}}\par
\textbf{linearly dependent} -- linear abhängig\dotfill\hyperref[e:319]{\pageref*{e:319}}\par
\textbf{logarithm} -- der Logarithmus\dotfill\hyperref[e:95]{\pageref*{e:95}}\par
\textbf{logarithmic equation} -- die logarithmische Gleichung\dotfill\hyperref[e:85]{\pageref*{e:85}}\par
\textbf{logarithmic function} -- die Logarithmusfunktion\dotfill\hyperref[e:131]{\pageref*{e:131}}\par
\textbf{lower sum} -- die Untersumme\dotfill\hyperref[e:294]{\pageref*{e:294}}\par
\textbf{lowest common denominator} -- der Hauptnenner\dotfill\hyperref[e:36]{\pageref*{e:36}}\par

\medskip{\large\textbf{\color{cornflower}M}}\par\nopagebreak
\textbf{magnitude} -- der Betrag eines Vektors\dotfill\hyperref[e:305]{\pageref*{e:305}}\par
\textbf{main diagonal} -- die Hauptdiagonale\dotfill\hyperref[e:324]{\pageref*{e:324}}\par
\textbf{matrix} -- die Matrix\dotfill\hyperref[e:320]{\pageref*{e:320}}\par
\textbf{maximum} -- der Hochpunkt\dotfill\hyperref[e:282]{\pageref*{e:282}}\par
\textbf{mean} -- der Mittelwert\dotfill\hyperref[e:338]{\pageref*{e:338}}\par
\textbf{median} -- der Median\dotfill\hyperref[e:339]{\pageref*{e:339}}\par
\textbf{median} -- die Seitenhalbierende\dotfill\hyperref[e:192]{\pageref*{e:192}}\par
\textbf{minimum} -- der Tiefpunkt\dotfill\hyperref[e:283]{\pageref*{e:283}}\par
\textbf{minuend} -- der Minuend\dotfill\hyperref[e:19]{\pageref*{e:19}}\par
\textbf{mixed number} -- die gemischte Zahl\dotfill\hyperref[e:33]{\pageref*{e:33}}\par
\textbf{much less than} -- viel kleiner als\dotfill\hyperref[e:67]{\pageref*{e:67}}\par
\textbf{multiple} -- das Vielfache\dotfill\hyperref[e:55]{\pageref*{e:55}}\par
\textbf{multiply} -- multiplizieren\dotfill\hyperref[e:21]{\pageref*{e:21}}\par

\medskip{\large\textbf{\color{cornflower}N}}\par\nopagebreak
\textbf{natural exponential function} -- die natürliche Exponentialfunktion\dotfill\hyperref[e:299]{\pageref*{e:299}}\par
\textbf{natural logarithm} -- der natürliche Logarithmus\dotfill\hyperref[e:300]{\pageref*{e:300}}\par
\textbf{natural number} -- die natürliche Zahl\dotfill\hyperref[e:5]{\pageref*{e:5}}\par
\textbf{net} -- das Netz\dotfill\hyperref[e:253]{\pageref*{e:253}}\par
\textbf{not equal} -- ungleich\dotfill\hyperref[e:61]{\pageref*{e:61}}\par
\textbf{$n$th root} -- die $n$-te Wurzel\dotfill\hyperref[e:47]{\pageref*{e:47}}\par
\textbf{null sequence} -- die Nullfolge\dotfill\hyperref[e:258]{\pageref*{e:258}}\par
\textbf{number} -- die Zahl\dotfill\hyperref[e:1]{\pageref*{e:1}}\par
\textbf{numerator} -- der Zähler\dotfill\hyperref[e:29]{\pageref*{e:29}}\par

\medskip{\large\textbf{\color{cornflower}O}}\par\nopagebreak
\textbf{oblique cylinder} -- der schiefe Zylinder\dotfill\hyperref[e:238]{\pageref*{e:238}}\par
\textbf{oblique drawing} -- das Schrägbild\dotfill\hyperref[e:254]{\pageref*{e:254}}\par
\textbf{obtuse angle} -- der stumpfe Winkel\dotfill\hyperref[e:159]{\pageref*{e:159}}\par
\textbf{octagon} -- das Achteck\dotfill\hyperref[e:207]{\pageref*{e:207}}\par
\textbf{opposite side} -- die Gegenkathete\dotfill\hyperref[e:189]{\pageref*{e:189}}\par
\textbf{ordinal number} -- die Ordnungszahl\dotfill\hyperref[e:3]{\pageref*{e:3}}\par
\textbf{ordinate} -- die Ordinate\dotfill\hyperref[e:118]{\pageref*{e:118}}\par
\textbf{origin} -- der Ursprung\dotfill\hyperref[e:115]{\pageref*{e:115}}\par
\textbf{outcome} -- das Ergebnis\dotfill\hyperref[e:327]{\pageref*{e:327}}\par

\medskip{\large\textbf{\color{cornflower}P}}\par\nopagebreak
\textbf{parabola} -- die Parabel\dotfill\hyperref[e:125]{\pageref*{e:125}}\par
\textbf{paraboloid of revolution} -- das Rotationsparaboloid\dotfill\hyperref[e:243]{\pageref*{e:243}}\par
\textbf{parallel} -- parallel\dotfill\hyperref[e:150]{\pageref*{e:150}}\par
\textbf{parallelogram} -- das Parallelogramm\dotfill\hyperref[e:200]{\pageref*{e:200}}\par
\textbf{partial derivative} -- die partielle Ableitung\dotfill\hyperref[e:277]{\pageref*{e:277}}\par
\textbf{pentagon} -- das Fünfeck\dotfill\hyperref[e:206]{\pageref*{e:206}}\par
\textbf{percentage} -- der Prozentsatz\dotfill\hyperref[e:58]{\pageref*{e:58}}\par
\textbf{perimeter} -- der Umfang\dotfill\hyperref[e:209]{\pageref*{e:209}}\par
\textbf{permutation} -- die Permutation\dotfill\hyperref[e:341]{\pageref*{e:341}}\par
\textbf{perpendicular} -- senkrecht\dotfill\hyperref[e:151]{\pageref*{e:151}}\par
\textbf{perpendicular bisector} -- die Mittelsenkrechte\dotfill\hyperref[e:170]{\pageref*{e:170}}\par
\textbf{pi} -- die Kreiszahl\dotfill\hyperref[e:216]{\pageref*{e:216}}\par
\textbf{plane} -- die Ebene\dotfill\hyperref[e:308]{\pageref*{e:308}}\par
\textbf{point} -- der Punkt\dotfill\hyperref[e:146]{\pageref*{e:146}}\par
\textbf{point of inflection} -- der Wendepunkt\dotfill\hyperref[e:284]{\pageref*{e:284}}\par
\textbf{point of intersection} -- der Schnittpunkt\dotfill\hyperref[e:123]{\pageref*{e:123}}\par
\textbf{point of tangency} -- der Berührungspunkt\dotfill\hyperref[e:220]{\pageref*{e:220}}\par
\textbf{polar coordinates} -- die Polarkoordinaten\dotfill\hyperref[e:311]{\pageref*{e:311}}\par
\textbf{polygon} -- das Vieleck\dotfill\hyperref[e:204]{\pageref*{e:204}}\par
\textbf{polyhedron} -- das Polyeder\dotfill\hyperref[e:233]{\pageref*{e:233}}\par
\textbf{polynomial function} -- die ganzrationale Funktion\dotfill\hyperref[e:126]{\pageref*{e:126}}\par
\textbf{position vector} -- der Ortsvektor\dotfill\hyperref[e:302]{\pageref*{e:302}}\par
\textbf{power} -- die Potenz\dotfill\hyperref[e:42]{\pageref*{e:42}}\par
\textbf{power series} -- die Potenzreihe\dotfill\hyperref[e:267]{\pageref*{e:267}}\par
\textbf{prime number} -- die Primzahl\dotfill\hyperref[e:12]{\pageref*{e:12}}\par
\textbf{prism} -- das Prisma\dotfill\hyperref[e:228]{\pageref*{e:228}}\par
\textbf{probability} -- die Wahrscheinlichkeit\dotfill\hyperref[e:330]{\pageref*{e:330}}\par
\textbf{product} -- das Produkt\dotfill\hyperref[e:22]{\pageref*{e:22}}\par
\textbf{proof} -- der Beweis\dotfill\hyperref[e:349]{\pageref*{e:349}}\par
\textbf{proper fraction} -- der echte Bruch\dotfill\hyperref[e:31]{\pageref*{e:31}}\par
\textbf{proper subset} -- die echte Teilmenge\dotfill\hyperref[e:99]{\pageref*{e:99}}\par
\textbf{proportion} -- die Verhältnisgleichung\dotfill\hyperref[e:41]{\pageref*{e:41}}\par
\textbf{proportional to} -- proportional zu\dotfill\hyperref[e:68]{\pageref*{e:68}}\par
\textbf{proportionality} -- die Proportionalität\dotfill\hyperref[e:94]{\pageref*{e:94}}\par
\textbf{pyramid} -- die Pyramide\dotfill\hyperref[e:230]{\pageref*{e:230}}\par
\textbf{Pythagorean theorem} -- der Satz des Pythagoras\dotfill\hyperref[e:190]{\pageref*{e:190}}\par

\medskip{\large\textbf{\color{cornflower}Q}}\par\nopagebreak
\textbf{quadratic equation} -- die quadratische Gleichung\dotfill\hyperref[e:82]{\pageref*{e:82}}\par
\textbf{quadratic function} -- die quadratische Funktion\dotfill\hyperref[e:124]{\pageref*{e:124}}\par
\textbf{quadrilateral} -- das Viereck\dotfill\hyperref[e:196]{\pageref*{e:196}}\par
\textbf{quotient} -- der Quotient\dotfill\hyperref[e:25]{\pageref*{e:25}}\par

\medskip{\large\textbf{\color{cornflower}R}}\par\nopagebreak
\textbf{radian measure} -- das Bogenmaß\dotfill\hyperref[e:139]{\pageref*{e:139}}\par
\textbf{radical equation} -- die Wurzelgleichung\dotfill\hyperref[e:83]{\pageref*{e:83}}\par
\textbf{radicand} -- der Radikand\dotfill\hyperref[e:48]{\pageref*{e:48}}\par
\textbf{radius} -- der Radius\dotfill\hyperref[e:213]{\pageref*{e:213}}\par
\textbf{raise to a power} -- potenzieren\dotfill\hyperref[e:50]{\pageref*{e:50}}\par
\textbf{random experiment} -- das Zufallsexperiment\dotfill\hyperref[e:326]{\pageref*{e:326}}\par
\textbf{random variable} -- die Zufallsgröße\dotfill\hyperref[e:334]{\pageref*{e:334}}\par
\textbf{range} -- der Wertebereich\dotfill\hyperref[e:110]{\pageref*{e:110}}\par
\textbf{ratio} -- das Verhältnis\dotfill\hyperref[e:40]{\pageref*{e:40}}\par
\textbf{rational function} -- die gebrochenrationale Funktion\dotfill\hyperref[e:127]{\pageref*{e:127}}\par
\textbf{rational number} -- die rationale Zahl\dotfill\hyperref[e:7]{\pageref*{e:7}}\par
\textbf{ray} -- der Strahl\dotfill\hyperref[e:149]{\pageref*{e:149}}\par
\textbf{real number} -- die reelle Zahl\dotfill\hyperref[e:8]{\pageref*{e:8}}\par
\textbf{reciprocal} -- der Kehrwert\dotfill\hyperref[e:39]{\pageref*{e:39}}\par
\textbf{rectangle} -- das Rechteck\dotfill\hyperref[e:199]{\pageref*{e:199}}\par
\textbf{reflection} -- die Spiegelung\dotfill\hyperref[e:172]{\pageref*{e:172}}\par
\textbf{reflex angle} -- der überstumpfe Winkel\dotfill\hyperref[e:161]{\pageref*{e:161}}\par
\textbf{regular hexagon} -- das regelmäßige Sechseck\dotfill\hyperref[e:205]{\pageref*{e:205}}\par
\textbf{relative frequency} -- die relative Häufigkeit\dotfill\hyperref[e:331]{\pageref*{e:331}}\par
\textbf{rhombus} -- die Raute\dotfill\hyperref[e:201]{\pageref*{e:201}}\par
\textbf{right angle} -- der rechte Winkel\dotfill\hyperref[e:158]{\pageref*{e:158}}\par
\textbf{right-angled triangle} -- das rechtwinklige Dreieck\dotfill\hyperref[e:181]{\pageref*{e:181}}\par
\textbf{rotation} -- die Drehung\dotfill\hyperref[e:174]{\pageref*{e:174}}\par
\textbf{round} -- runden\dotfill\hyperref[e:59]{\pageref*{e:59}}\par
\textbf{round brackets, parentheses} -- die runden Klammern\dotfill\hyperref[e:88]{\pageref*{e:88}}\par
\textbf{row} -- die Zeile\dotfill\hyperref[e:322]{\pageref*{e:322}}\par

\medskip{\large\textbf{\color{cornflower}S}}\par\nopagebreak
\textbf{scalene triangle} -- das unregelmäßige Dreieck\dotfill\hyperref[e:180]{\pageref*{e:180}}\par
\textbf{secant} -- die Sekante\dotfill\hyperref[e:219]{\pageref*{e:219}}\par
\textbf{sector} -- der Kreissektor\dotfill\hyperref[e:217]{\pageref*{e:217}}\par
\textbf{segment (of a circle)} -- der Kreisabschnitt\dotfill\hyperref[e:222]{\pageref*{e:222}}\par
\textbf{semi-axis} -- die Halbachse\dotfill\hyperref[e:317]{\pageref*{e:317}}\par
\textbf{sequence} -- die Folge\dotfill\hyperref[e:257]{\pageref*{e:257}}\par
\textbf{series} -- die Reihe\dotfill\hyperref[e:265]{\pageref*{e:265}}\par
\textbf{set} -- die Menge\dotfill\hyperref[e:96]{\pageref*{e:96}}\par
\textbf{set difference} -- die Differenzmenge\dotfill\hyperref[e:102]{\pageref*{e:102}}\par
\textbf{similar} -- ähnlich\dotfill\hyperref[e:176]{\pageref*{e:176}}\par
\textbf{sine} -- der Sinus\dotfill\hyperref[e:132]{\pageref*{e:132}}\par
\textbf{sine rule, law of sines} -- der Sinussatz\dotfill\hyperref[e:140]{\pageref*{e:140}}\par
\textbf{solid} -- der Körper\dotfill\hyperref[e:225]{\pageref*{e:225}}\par
\textbf{solid of revolution} -- der Rotationskörper\dotfill\hyperref[e:239]{\pageref*{e:239}}\par
\textbf{solution} -- die Lösung\dotfill\hyperref[e:78]{\pageref*{e:78}}\par
\textbf{solution set} -- die Lösungsmenge\dotfill\hyperref[e:79]{\pageref*{e:79}}\par
\textbf{space diagonal} -- die Raumdiagonale\dotfill\hyperref[e:252]{\pageref*{e:252}}\par
\textbf{sphere} -- die Kugel\dotfill\hyperref[e:232]{\pageref*{e:232}}\par
\textbf{spherical cap, spherical segment} -- das Kugelsegment\dotfill\hyperref[e:240]{\pageref*{e:240}}\par
\textbf{spherical sector} -- der Kugelsektor\dotfill\hyperref[e:241]{\pageref*{e:241}}\par
\textbf{spheroid} -- das Rotationsellipsoid\dotfill\hyperref[e:242]{\pageref*{e:242}}\par
\textbf{square} -- das Quadrat\dotfill\hyperref[e:198]{\pageref*{e:198}}\par
\textbf{square} -- quadrieren\dotfill\hyperref[e:51]{\pageref*{e:51}}\par
\textbf{square brackets} -- die eckigen Klammern\dotfill\hyperref[e:89]{\pageref*{e:89}}\par
\textbf{square matrix} -- die quadratische Matrix\dotfill\hyperref[e:321]{\pageref*{e:321}}\par
\textbf{square root} -- die Quadratwurzel\dotfill\hyperref[e:45]{\pageref*{e:45}}\par
\textbf{square root function} -- die Wurzelfunktion\dotfill\hyperref[e:128]{\pageref*{e:128}}\par
\textbf{standard deviation} -- die Standardabweichung\dotfill\hyperref[e:337]{\pageref*{e:337}}\par
\textbf{straight angle} -- der gestreckte Winkel\dotfill\hyperref[e:160]{\pageref*{e:160}}\par
\textbf{straight line} -- die Gerade\dotfill\hyperref[e:147]{\pageref*{e:147}}\par
\textbf{subscript, index} -- der Index\dotfill\hyperref[e:4]{\pageref*{e:4}}\par
\textbf{subset} -- die Teilmenge\dotfill\hyperref[e:98]{\pageref*{e:98}}\par
\textbf{subtract} -- subtrahieren\dotfill\hyperref[e:17]{\pageref*{e:17}}\par
\textbf{subtrahend} -- der Subtrahend\dotfill\hyperref[e:20]{\pageref*{e:20}}\par
\textbf{sum} -- die Summe\dotfill\hyperref[e:15]{\pageref*{e:15}}\par
\textbf{summand, addend} -- der Summand\dotfill\hyperref[e:16]{\pageref*{e:16}}\par
\textbf{supplementary angle} -- der Nebenwinkel\dotfill\hyperref[e:163]{\pageref*{e:163}}\par
\textbf{supplementary angle} -- der Supplementwinkel\dotfill\hyperref[e:165]{\pageref*{e:165}}\par
\textbf{surface area} -- der Oberflächeninhalt\dotfill\hyperref[e:256]{\pageref*{e:256}}\par
\textbf{symmetric about the origin} -- punktsymmetrisch zum Ursprung\dotfill\hyperref[e:145]{\pageref*{e:145}}\par
\textbf{symmetric about the $y$-axis} -- achsensymmetrisch zur $y$-Achse\dotfill\hyperref[e:144]{\pageref*{e:144}}\par
\textbf{system of equations} -- das Gleichungssystem\dotfill\hyperref[e:80]{\pageref*{e:80}}\par

\medskip{\large\textbf{\color{cornflower}T}}\par\nopagebreak
\textbf{table of values} -- die Wertetabelle\dotfill\hyperref[e:112]{\pageref*{e:112}}\par
\textbf{tangent} -- der Tangens\dotfill\hyperref[e:134]{\pageref*{e:134}}\par
\textbf{tangent} -- die Tangente\dotfill\hyperref[e:224]{\pageref*{e:224}}\par
\textbf{tangent} -- die Tangente\dotfill\hyperref[e:279]{\pageref*{e:279}}\par
\textbf{tetrahedron} -- das Tetraeder\dotfill\hyperref[e:234]{\pageref*{e:234}}\par
\textbf{theorem} -- der Satz\dotfill\hyperref[e:347]{\pageref*{e:347}}\par
\textbf{there exists} -- es gibt\dotfill\hyperref[e:353]{\pageref*{e:353}}\par
\textbf{top face} -- die Deckfläche\dotfill\hyperref[e:249]{\pageref*{e:249}}\par
\textbf{translation} -- die Verschiebung\dotfill\hyperref[e:173]{\pageref*{e:173}}\par
\textbf{trapezium} -- das Trapez\dotfill\hyperref[e:202]{\pageref*{e:202}}\par
\textbf{tree diagram} -- das Baumdiagramm\dotfill\hyperref[e:332]{\pageref*{e:332}}\par
\textbf{triangle} -- das Dreieck\dotfill\hyperref[e:177]{\pageref*{e:177}}\par
\textbf{trigonometric equation} -- die trigonometrische Gleichung\dotfill\hyperref[e:86]{\pageref*{e:86}}\par
\textbf{two-way table} -- die Vierfeldertafel\dotfill\hyperref[e:333]{\pageref*{e:333}}\par

\medskip{\large\textbf{\color{cornflower}U}}\par\nopagebreak
\textbf{union} -- die Vereinigungsmenge\dotfill\hyperref[e:101]{\pageref*{e:101}}\par
\textbf{unit circle} -- der Einheitskreis\dotfill\hyperref[e:137]{\pageref*{e:137}}\par
\textbf{unit vector} -- der Einheitsvektor\dotfill\hyperref[e:303]{\pageref*{e:303}}\par
\textbf{unknown} -- die Unbekannte\dotfill\hyperref[e:73]{\pageref*{e:73}}\par
\textbf{unlike fractions} -- ungleichnamige Brüche\dotfill\hyperref[e:35]{\pageref*{e:35}}\par
\textbf{upper sum} -- die Obersumme\dotfill\hyperref[e:293]{\pageref*{e:293}}\par

\medskip{\large\textbf{\color{cornflower}V}}\par\nopagebreak
\textbf{value of a function} -- der Funktionswert\dotfill\hyperref[e:108]{\pageref*{e:108}}\par
\textbf{variable} -- die Variable\dotfill\hyperref[e:71]{\pageref*{e:71}}\par
\textbf{variable of integration} -- die Integrationsvariable\dotfill\hyperref[e:291]{\pageref*{e:291}}\par
\textbf{variance} -- die Varianz\dotfill\hyperref[e:336]{\pageref*{e:336}}\par
\textbf{variation, arrangement} -- die Variation\dotfill\hyperref[e:342]{\pageref*{e:342}}\par
\textbf{vector} -- der Vektor\dotfill\hyperref[e:301]{\pageref*{e:301}}\par
\textbf{vertex} -- die Ecke\dotfill\hyperref[e:245]{\pageref*{e:245}}\par
\textbf{vertex} -- der Scheitel\dotfill\hyperref[e:154]{\pageref*{e:154}}\par
\textbf{vertex} -- der Scheitelpunkt\dotfill\hyperref[e:129]{\pageref*{e:129}}\par
\textbf{vertically opposite angle} -- der Scheitelwinkel\dotfill\hyperref[e:166]{\pageref*{e:166}}\par
\textbf{volume} -- das Volumen\dotfill\hyperref[e:255]{\pageref*{e:255}}\par

\medskip{\large\textbf{\color{cornflower}X}}\par\nopagebreak
\textbf{$x$-axis} -- die $x$-Achse\dotfill\hyperref[e:114]{\pageref*{e:114}}\par

\medskip{\large\textbf{\color{cornflower}Y}}\par\nopagebreak
\textbf{$y$-intercept} -- der $y$-Achsenabschnitt\dotfill\hyperref[e:121]{\pageref*{e:121}}\par

\medskip{\large\textbf{\color{cornflower}Z}}\par\nopagebreak
\textbf{zero vector} -- der Nullvektor\dotfill\hyperref[e:304]{\pageref*{e:304}}\par
\textbf{zero, root} -- die Nullstelle\dotfill\hyperref[e:122]{\pageref*{e:122}}\par
\end{multicols}

\end{document}
